%
% If you are looking for the PDF file, you can compile the .tex file or donwload the PDF from here:
% http://etd.adm.unipi.it/theses/available/etd-09232013-170716/
% 
% Istruzioni di compilazione (per produrre il pdf): avere un texlive funzionante e i pacchetti che vengono richiesti qui sotto; lanciare i comandi
% pdflatex gelmetti-master_thesis.tex
% makeglossaries gelmetti-master_thesis
% biber gelmetti-master_thesis
% pdflatex gelmetti-master_thesis.tex
% pdflatex gelmetti-master_thesis.tex
% non so se sia necessario lanciare altre volte biber, makeglossaries o pdflatex, di certo male non fa.
%

\documentclass[a4paper,12pt,openright]{book}
\usepackage[utf8]{inputenc}
\usepackage[italian,english]{babel}
\usepackage[binary-units]{siunitx} % Typesetting quantities in a consistent way
\usepackage{chemmacros} 		% \ch{}
\usepackage[backend=biber,labelnumber=true,url=false,isbn=false,maxcitenames=10,firstinits=true,sorting=none]{biblatex}
\usepackage{chemnum}			% Package used to manage index of compounds and schemes
\usepackage{graphicx}
\usepackage{subcaption} %the subfigure and subfig packages are deprecated
\usepackage{csquotes} 
\usepackage[wide,ragged]{sidecap}%[ragged, margincaption, centerbody]
\usepackage{ragged2e} % needed by ragged option of sidecap?
\usepackage{chngcntr} % needed for setting option not resetting footnote counter at chapter end
\usepackage{pgfplots} % graphs
\usepackage{booktabs}
\usepackage{lscape} %for big landscape figures, rotate the page
\usepackage{afterpage} %for making landscape float
\usepackage[section]{placeins} %keeps floats `in their place', preventing them from floating past a "\FloatBarrier" command into another section.
\usepackage{float}
\usepackage{eso-pic} %background image on titlepage
\usepackage{transparent} %transparency for background image on titlepage
\usepackage[format=hang,labelfont=bf,textfont=small]{caption}%personalize figure and table captions
\usepackage[xindy,nonumberlist,acronym,nomain]{glossaries} %table of acronyms
\usepackage{relsize}%needed by smaller in glossaries
\usepackage{multirow}% for tables
\usepackage[bottom,multiple]{footmisc} %bottom insert additional whitespace between the main text and the footnote rule rather than over-stretching the inter-paragraph glue AND to prevent the floats going below footnotes. This has to be loaded before fancyhdr
\usepackage{fancyhdr}%page headers
\usepackage[hmarginratio=1:1,vmarginratio=1:1,a4paper,marginparwidth=1.5cm,textheight=630pt]{geometry}% marginparwidth inserted only for wide sidecaps
\usepackage{amssymb}% for the \square math symbol

\renewcommand{\sfdefault}{cmbr}%for lighter sans serif
\captionsetup[SCfigure]{format=plain}
\pgfplotsset{compat=newest}
\counterwithout{footnote}{chapter} % not resetting footnote counter at chapter end
\newcommand{\HNMR}{$^1$\ch{H}-\textsmaller{NMR}}
\newcommand{\CNMR}{$^{13}$\ch{C}-\textsmaller{NMR}}
\pgfplotsset{tick label style={font=\small},label style={font=\small}}
\let\oldch\ch %move the ch command, needed for redefinition
\renewcommand{\ch}[1]{\textsf{\oldch{#1}}} %add sans serif to \ch

\DeclareSIUnit\Molar{M}
\DeclareSIUnit\rpm{rpm}
\DeclareSIUnit\eq{eq}
\DeclareSIUnit\eV{eV}
\DeclareSIUnit\ppm{ppm}
\DeclareSIUnit\unit{unit}

\cmpd*[cmpd-label = 1S]{ig2-1} \cmpd*[cmpd-label = 2S]{ig2-2} \cmpd*[cmpd-label = P1S]{ig2-4} \cmpd*[cmpd-label = P2S]{ig2-8} \cmpd*[cmpd-label = 1]{ig2-9} \cmpd*[cmpd-label = 2]{ig2-10} \cmpd*[cmpd-label = P2]{ig2-15} \cmpd*[cmpd-label = 3]{ig2-16} \cmpd*[cmpd-label = 4]{ig2-17} \cmpd*[cmpd-label = P2SA]{ig2-19} \cmpd*[cmpd-label = P2ST]{ig2-22} \cmpd*[cmpd-label = 5]{ig2-21} \cmpd*[cmpd-label = P2S-PVP]{sz17}
\newcommand{\tallcell}[2][c]{\begin{tabular}[#1]{@{}c@{}}#2\end{tabular}}
\pgfplotsset{minor grid style={densely dotted,lightgray}}

\bibliography{biblio}

\newcounter{countA}
\newcounter{countB}
\newcounter{countC}
\makeatletter%
\DeclareCiteCommand{\superfootcite}[\cbx@wrap]%
{\gdef\cbx@keys{}}%
{\xappto\cbx@keys{\thefield{entrykey},}}%
{}%
{%
\ifcsundef{cbx@lastin@\cbx@keys @\strfield{postnote}}%
{\csnumgdef{cbx@lastin@\cbx@keys @\strfield{postnote}}{-1}}{}%
\ifnumcomp{\value{countA}}{<}{\csuse{cbx@lastin@\cbx@keys @\strfield{postnote}}}%
{%
\footnotemark[\csuse{cbx@lastfn@\cbx@keys @\strfield{postnote}}]%
}%
{\addtocounter{footnote}{1}\addtocounter{countB}{1}%
\xappto\cbx@cite{\noexpand\footnotemark[\value{countB}]%
\noexpand\let\noexpand\thefootnote\noexpand\relax\noexpand\footnotetext%
{%
\noexpand\textsuperscript{\the\value{countB}}\fullcite{\cbx@keys}}%
\csnumgdef{cbx@lastfn@\cbx@keys @\strfield{postnote}}{\value{\@mpfn}}%
\csnumgdef{cbx@lastin@\cbx@keys @\strfield{postnote}}{\value{instcount}}}%
\setcounter{countA}{\csuse{cbx@lastin@\cbx@keys @\strfield{postnote}}}}%
}%
\newrobustcmd{\cbx@wrap}[1]{#1\cbx@cite\gdef\cbx@cite{}}%
\def\cbx@cite{}%
\makeatother%

\newcommand{\nota}[1]{\addtocounter{countC}{1}\textsuperscript{\roman{countC}}\addtocounter{footnote}{-1}{\let\thefootnote\relax\footnote{\textsuperscript{\roman{countC}} #1}}}




\renewbibmacro*{name:andothers}{% Based on name:andothers from biblatex.def for "et al." in italic
\ifboolexpr{
test {\ifnumequal{\value{listcount}}{\value{liststop}}}
and
test \ifmorenames
}
{\ifnumgreater{\value{liststop}}{1}
{\finalandcomma}
{}%
\andothersdelim\bibstring[\emph]{andothers}}
{}}

\hyphenation{tetra-hydro-furan di-chloro-methane enantio-purity oxy-methyl-ene di-chroic aceto-nitrile iso-propyl alk-oxy-amine poly-thio-phene di-chloride inter-chain photo-luminescence supra-molecular photo-gene-rated hetero-junction poly-thio-phenic di-block thio-phenic poly-thio-phene bi-layer macro-initiator poly-vinyl-pyridine homo-polymer path-length intra-molecular inter-molecular hydro-chloric poly-dispersity poly-disperse}
\makeglossaries

\title{Synthesis and Characterization of Chiral Polythiophenes}
\author{Ilario Gelmetti}
\date{15.10.2013}
\begin{document}
\pagestyle{plain}

\selectlanguage{english}
\frontmatter

{\let\cleardoublepage\clearpage % i don't want openright in frontmatter

\begin{titlepage}\begin{center}


\AddToShipoutPicture*{
\put(0,0){
\parbox[b][\paperheight]{\paperwidth}{%
\vfill
\centering
\transparent{0.02}%stampato risulta randomicamente più scuro o più chiato, on paper this image results much more visible (or less, randomly) than on screen
	{\includegraphics[width=0.9\paperwidth]{img/unipi-logo.pdf}}%
\vfill
}
}
}

	\makebox[\textwidth]{\centering{\includegraphics[width=.4\textwidth]{img/unipi-logo_con_testo.png}}}\\%
	\medskip
	\large{\textsc{Dipartimento di Chimica e Chimica Industriale}}\\
		\rule{5cm}{1pt}\\
	{
\smallskip
	{Master thesis in \textsc{Organic Chemistry}}}\\
\vfill

	\huge{\textsc{Synthesis and Characterization of Chiral Polythiophenes.}}\\
\vfill
	\footnotesize{Author:}\\
	\large{Ilario Gelmetti}\\
		\makebox[.2\textwidth]{\rule{0pt}{.02\textheight}}\\
\end{center}
\begin{small}
\makebox{\parbox[b]{.8\textwidth}{



	\flushleft{\textbf{Supervisors:}}
		\flushright{\begin{tabular}{r c}
		{\large Dr. Silvia Destri}&\makebox[.4\textwidth]{\dotfill}\\
		&\\
		&\\
		{\large Prof. Lorenzo Di Bari}&\makebox[.4\textwidth]{\dotfill}\\
		\end{tabular}}
	\bigskip

	\flushleft{\textbf{Author:}}
		\flushright{\begin{tabular}{r c}
		{\large Ilario Gelmetti}&\makebox[.4\textwidth]{\dotfill}\\
		\end{tabular}}

	\flushleft{\textbf{Examiner:}}
		\flushright{\begin{tabular}{l c}
		{\large Prof. Roberto Solaro}&\makebox[.4\textwidth]{\dotfill}\\
		\end{tabular}}

}}
\bigskip
	\begin{center}
	\rule{3cm}{1pt}\\
	Session: 15 October 2013\\
	Academic year: 2012-2013\\
	\end{center}
\end{small}
\end{titlepage}

\addtocounter{page}{-1}

\tableofcontents
\listoffigures
\listoftables

\newacronym{NMRP}{\textsmaller{NMRP}}{nitroxide-mediated radical polymerization}
\newacronym{GRIM}{\textsmaller{GRIM}}{Grignard metathesis}
\newacronym{KCTP}{\textsmaller{KCTP}}{Kumada catalyst-transfer polycondensation}
\newacronym{HOMO}{\textsmaller{HOMO}}{highest-energy occupied molecular orbital}
\newacronym{LUMO}{\textsmaller{LUMO}}{lowest-energy unoccupied molecular orbital}
\newacronym{PL}{\textsmaller{PL}}{photoluminescence}
\newacronym{UVvis}{\textsmaller{UV}--vis}{ultraviolet--visible}
\newacronym{P3HT}{\textsmaller{P3HT}}{poly(3-hexyl\-thiophene)}
\newacronym{SEC}{\textsmaller{SEC}}{size exclusion chromatography}
\newacronym{NMR}{\textsmaller{NMR}}{nuclear magnetic resonance}
\newacronym{CD}{\textsmaller{CD}}{circular dichroism}
\newacronym{DSC}{\textsmaller{DSC}}{differential scanning calorimetry}
\newacronym{XRD}{\textsmaller{XRD}}{X--ray diffraction}
\newacronym{PCBM}{\textsmaller{PC}$_{60}$\textsmaller{BM}}{phenyl-\ch{C}61-butyric acid methyl ester}
\newacronym{BHJ}{\textsmaller{BHJ}}{bulk heterojunction}
\newacronym{GCMS}{\textsmaller{GC--MS}}{gas chromatography--mass spectrometry}
\newacronym{EI}{\textsmaller{EI}}{electron ionization}
\newacronym{P4VP}{\textsmaller{P4VP}}{poly\-(4-vinyl\-pyridine)}
\newacronym{nidppp}{\ch{Ni(dppp)Cl2}}{di\-chloro\-(1,3-bis(di\-phenyl\-phosph\-ino)\-propane)\-nickel} 
\newacronym{nidppe}{\ch{Ni(dppe)Cl2}}{di\-chloro\-(1,3-bis(di\-phenyl\-phosph\-ino)\-ethane)\-nickel} 
\newacronym{MALDI}{\textsmaller{MALDI--TOF MS}}{matrix-assisted laser desorption-ionization--time-of-flight mass spectrometry}
\newacronym{NMR2D}{\textsmaller{NMR-2D}}{two-dimensional nuclear magnetic resonance}
\newacronym{FTIR}{\textsmaller{FTIR}}{Fourier transform infrared spectroscopy}
\newacronym{XRD2D}{\textsmaller{XRD-2D}}{two-dimensional X-ray diffraction}
\newacronym{BrPhEtTIPNO}{\ch{BrPhEt}-\textsmaller{TIPNO}}{\textit{N-tert}-butyl-\textit{N}-(2-methyl-1-phenyl\-propyl)-\textit{O}-(1-(4-bromo\-phenyl)\-ethyl)\-hydroxyl\-amine}
\newacronym{THF}{\textsmaller{THF}}{tetrahydrofuran}
\newacronym{TCE}{\textsmaller{TCE}}{1,1,2,2-tetrachloroethane}
\newacronym{WAXD}{\textsmaller{WAXD}}{wide-angle X--ray diffraction}
\newacronym{SAXS}{\textsmaller{SAXS}}{small-angle X-ray scattering}
\newacronym{TIPNO}{\textsmaller{TIPNO}}{2,2,5-tri\-methyl-4-phenyl-3-aza\-hexane-3-nitr\-oxide}
\newacronym{PTFE}{\textsmaller{PTFE}}{poly\-tetra\-fluoro\-ethylene}
\newacronym{FTO}{\textsmaller{FTO}}{fluorine tin oxide}
\newacronym{PEDOT:PSS}{\textsmaller{PEDOT:PSS}}{poly\-(3,4-ethylene\-di\-oxy\-thiophene) poly\-(styrene\-sulfonate)}
\newacronym{FF}{\textsmaller{FF}}{fill factor}
\newacronym{PDI}{\textsmaller{PDI}}{polydispersity index}
\newacronym{dppp}{dppp}{1,3-bis\-(di\-phenyl\-phosphino)\-propane}
\newacronym{TLC}{\textsmaller{TLC}}{thin layer chromatography}


\printglossary[type=\acronymtype
]

}%here ends openany and starts openright

\cleardoublepage \vspace*{12cm}
\begin{flushright}
\large {\it Dedicata a mio padre. E a mia madre.}
\end{flushright}

\newgeometry{top=1in, bottom=1in, outer=1in, inner=1in}

\chapter[Abstract]{\centering Abstract}
We synthesized a novel semiconducting, regioregular chiral polythiophene with low polydispersity. 
Two synthetic procedures and two enantiomeric excesses were employed and the corresponding materials (\cmpd+{ig2-4}, enantiopure, about 20 repeating units; \cmpd+{ig2-15}, enantiomerically enriched, about 60 units; \cmpd+{ig2-8}, enantiopure, about 50 units) were deeply characterized from molecular (\acrshort{NMR}, \acrshort{NMR2D}, \acrshort{MALDI} and \acrshort{SEC}), thermal (\acrshort{DSC}), structural (\acrshort{XRD}, \acrshort{XRD2D} and \acrshort{SAXS}) and optical (\acrshort{FTIR}, \acrshort{UVvis}, \acrshort{CD} and \acrshort{PL}) points of view. 
Solutions in good and poor solvents, thin films and powders were examined. 
The process of aggregate formation was followed by means of \acrshort{UVvis}, \acrfull{CD} and \acrfull{PL} spectroscopies in order to study the conformational geometry of the forming structures and investigate their relationship with the origin of the optical activity. The longest enantiopure poly\-thio\-phene was thereafter used as donor component of solar cell active layer and, after proper functionalization, as a macro\-initiator for the polymerization of a diblock copolymer with poly\-vinyl\-pyridine. This copolymer was characterized with \acrshort{NMR} and its aggregation studied with \acrshort{UVvis} and \acrshort{CD} spectroscopies.

{
\vfill
\centering \rule{5cm}{1pt}\\
\vfill
}

\selectlanguage{italian}
\noindent
In questo lavoro di tesi ho sintetizzato un poli\-tiofene regio\-regolare chirale avente bassa poli\-dispersità. Lo specifico semiconduttore organico che ho preparato non è mai stato riportato in letteratura. Sono state utilizzate due differenti procedure di polimerizzazione e due diversi eccessi enantiomerici ed i polimeri ottenuti (\cmpd+{ig2-4}, enantiopuro, circa 20 unità ripetitive; \cmpd+{ig2-15}, enantiomericamente arricchito, circa 60 unità; \cmpd+{ig2-8}, enantiopuro, circa 50 unità) sono stati ampiamente caratterizzati dal punto di vista molecolare (\acrshort{NMR}, \acrshort{NMR2D}, \acrshort{MALDI} e \acrshort{SEC}), termico (\acrshort{DSC}), strutturale (\acrshort{XRD}, \acrshort{XRD2D} e \acrshort{SAXS}) e ottico (\acrshort{FTIR}, \acrshort{UVvis}, \acrshort{CD} e \acrshort{PL}). I polimeri sono stati analizzati in soluzione di buon solvente, in miscela con cattivo solvente, in stato solido come film sottile e come polvere. 
Il processo di aggregazione è stato seguito con spettroscopie \acrshort{UVvis}, dicroismo circolare (\acrshort{CD}) e fotoluminescenza (\acrshort{PL}) per studiare la conformazione delle strutture in crescita e l'origine supramolecolare dell'attività ottica. 
Il politiofene enantiopuro %ottenuto 
a più alto peso molecolare è stato testato come materiale donore in celle solari e, una volta adeguatamente funzionalizzato, è stato impiegato come macroiniziatore per la polimerizzazione di un secondo blocco polimerico di poli\-(vinil\-piridina). Il polimero a blocchi ottenuto è stato caratterizzato all'\acrshort{NMR} e la sua aggregazione è stata seguita con le spettroscopie \acrshort{UVvis} e \acrshort{CD}.

\selectlanguage{english}
\restoregeometry

\fancypagestyle{plain}{% chapter pages use style plain
\renewcommand{\headrulewidth}{0pt} % no top rule
\fancyhf{}% clear all fields
\fancyfoot[L]{\footnotesize{Ilario Gelmetti}}% left footer
\fancyfoot[C]{\footnotesize{\textit{Synthesis and Characterization of Chiral Polythiophenes}}}
\fancyfoot[R]{\thepage}% right footer
}

{
\newpage
\null\vfill
% \vspace*{\fill}
\begin{center}
\includegraphics[width=0.3\textwidth]{img/cc-by-sa.pdf}

\smallskip

\textit{Creative Commons Attribution--ShareAlike license}
\end{center}

\bigskip

\noindent
You are free:


to Share -- to copy, distribute and transmit the work
\indent

to Remix -- to adapt the work
\indent

to make commercial use of the work

\medskip

\noindent
Under the following conditions:
% \smallskip

Attribution -- You must attribute the work in the manner specified by the author or licensor (but not in any way that suggests that they endorse you or your use of the work).
\indent

Share Alike -- If you alter, transform, or build upon this work, you may distribute the resulting work only under the same or similar license to this one. 
}

\mainmatter
\pagestyle{headings}
\chapter{Introduction}
\graphicspath{ {./img/intro/} }
\begin{section}{$\pi$-Conjugated Polymers and Polythiophene}
\paragraph{Conjugated polymers} Conjugated polymers are attracting great interest and research effort due to their semiconducting properties, suitable to substitute classical inorganic semiconductors (for example silicon) in lightweight and flexible devices. Inconveniently conjugated polymers are generally insoluble and infusible materials due to their rigidity and to strong intermolecular $\pi-\pi$ stacking interactions. In order to increase the solubility and to modulate the material properties, conjugated polymers have been substituted with appropriate side groups.

\begin{SCfigure}[][tbp]%conjugated_polymers
\centering
\includegraphics[scale=0.6]{conjugated_polymers.pdf}
\caption{Some conjugated polymers.}
\label{fig:conjugated_polymers}
\end{SCfigure}

\paragraph{Polythiophene} Polythiophene was the first material to give good results in a broad spectrum of applications such as light emitting diodes,\superfootcite{Ohmori1991} field effect transistors\superfootcite{Garnier1994} and organic solar cells\superfootcite{ADMA:ADMA19950070709} due to its convenient optical and electrical properties as well as excellent thermal and chemical stability. Polythiophene properties can be tuned via an easy substitution at one or two of the $\beta$ positions of the thiophene ring and these side chains can add solubility and processability properties to poly\-thio\-phenes. 
The polymer electronic and electrochemical properties can be tuned by varying the electronic properties of the side groups. Also the bulkiness of these groups is important, it can lead, via reciprocal steric repulsion, to a twist in the poly\-thiophenic backbone, reducing the conjugation length and ruining the semiconducting properties. In order to minimize steric interactions, usually only one side group per thiophene unit is bonded in $\beta$ position, doing so the monomeric unit results to be asymmetric. 

\begin{figure}[tbp]%regioregolare-regioirregolare
\begin{subfigure}[b]{0.47\textwidth}
\includegraphics[width=1\textwidth]{regioregolare.pdf}
\subcaption{regioregular}\label{fig:regioregolare}
\end{subfigure}
\qquad
\begin{subfigure}[b]{0.47\textwidth}
\includegraphics[width=1\textwidth]{regioirregolare.pdf}
\subcaption{regioirregular}\label{fig:regioirregolare}
\end{subfigure}
\caption{Regioregularity in monosubstituted polythiophenes.}\label{fig:regioregolare-regioirregolare}
\end{figure}

\paragraph{Polymerization} As shown in Figure~\ref{fig:regioregolare-regioirregolare}, the polymerization of mono\-substituted thiophene can produce polymers in which monomers are randomly connected giving rise to a irregular polymer or polymers in which the repeating units are regularly oriented. All these polymers show different properties and behave as different materials.\superfootcite{Kim2006} 
The polymerization method that allows us to obtain a better regioregularity is a transition metal catalyzed coupling between thiophenic units regioselectively substituted in the $\alpha$ positions with a halogen and a metal. \label{intro-polimerizzazione}
In McCullough method\superfootcite{McCullough1992}\superfootcite{McCullough1993c}\superfootcite{McCullough1993a}\superfootcite{Loewe1999}\superfootcite{Loewe2001} this metal is a magnesium atom and the catalyst is a nickel complex with a wide bite angle\nota{\label{fn:biteangle}\textit{bite angle}, the chelation angle between the ligands of a metal atom, $\widehat{\mathrm{LML}}$.} 
ligand: \acrfull{nidppp}. 
Also in Rieke method\superfootcite{Chen1992}\superfootcite{Chen1995} a regioregular poly\-(3-alkyl\-thio\-phene) is obtained, but employing zinc as the activating metal and \acrfull{nidppe} as catalyst. 
Over the last years the McCullough method was optimized resulting in a quasi-living polymerization\superfootcite{Sheina2004}\superfootcite{Iovu2005}\superfootcite{Yokoyama2004}\superfootcite{Miyakoshi2005} allowing one to gain control on the degree of polymerization, the polydispersity\nota{\label{fn:polydispersity}\textit{polydispersity} or \acrshort{PDI}, a polymer comprising molecules non-uniform with respect to relative molecular mass.} 
and on polymeric terminations. Chain ending control allows us to obtain an asymmetric polythiophene with a functional group on a termination and a different one on the other one. This is useful for joining the polythiophene to other reacting sites. 
Polydispersity and molecular mass control helps one to control aggregation or crystallization in the solid state. 

\end{section}
\begin{section}{Polythiophene and Chirality}
\paragraph{Aggregated state} The solid state structure is a crucial aspect for electronic devices, for example the interchain charge mobility in the material strongly depends on the extent of the $\pi-\pi$ stacking in solid state. \Gls{UVvis} spectroscopy can be used to study the variations of the \acrshort{HOMO}--\acrshort{LUMO} transition which for polythiophenes is in the visible region. The wavelength of this transition depends on the effective conjugation length and on the exciton coupling between stacked chains. 
When aggregation occurs, the chains adopt a more planar conformation, which extends the conjugation length in order to pack. Hence the aggregation process can be spotted in \gls{UVvis} spectra as a decrease in energy of the \acrshort{HOMO}--\acrshort{LUMO} transition, a red shift of the wavelength. 
The aggregation phenomenon can be observed not only passing from solution to solid state, but also, for example, by adding a poor solvent to a solution of polymer. When the solid is regularly ordered also \acrfull{XRD} can be useful for the characterization. But when the solid sample 
is a blend of materials, usually needed for devices preparation, the chaotic soft matter is difficult to characterize on a molecular level. 

\paragraph{Chirality} \label{intro-cd}\Gls{CD} spectroscopy can help us in this operation, but it is sensitive only to chromophores in chiral environments, thus a source of chirality has to be introduced in the material. This approach has proven to give useful information on the aggregate structure of various conjugated polymers\superfootcite{Lakhwani2011}\superfootcite{Babudri2006}\superfootcite{Kane-Maguire2010} and various polythiophenes.\superfootcite{Bidan1996}\superfootcite{Goto2002}\superfootcite{Langeveld-Voss2000}\superfootcite{Goto2002a}\superfootcite{Bouman1995}\superfootcite{Langeveld-Voss1999} 

\paragraph{Chiral group} A convenient way to have chirality in polythiophene chains is to polymerize a thiophene bearing a chiral side group. The introduction of a chiral side group, that has to be branched, will have some influence on aggregation and, by consequence, on the material properties.

\end{section}
\begin{section}{Application of Polythiophene in Photovoltaics}\label{intro-pv}

\paragraph{Organic photovoltaics} Organic photovoltaic cells are promising for lightweight, flexible and cheap (compared to silicon) devices. These organic devices can be rather complex and have a plethora of variables which need to be optimized in order to reach the maximum theorized 10~\% efficiency.\superfootcite{Scharber2006}

\begin{SCfigure}[][tbp]%bilayer
\centering
\includegraphics[width=0.5\textwidth]{bilayer-gray.pdf}
\caption[Structure of a bilayer solar cell.]{Structure of a bilayer solar cell. Only excitons generated near an interface can break in free charges.}
\label{fig:bilayer}
\end{SCfigure}

\paragraph{Active layer morphologies} Early organic solar cells were made only of a single organic material sandwiched between two electrodes,\superfootcite{Chamberlain1983} here the charge generation is dependent on the electric potential coming from electrodes and the efficiencies are poor. 
Then bilayer cells (schematized in Figure~\ref{fig:bilayer}) were developed having a donor (\textit{p}-type semiconductor) and an acceptor (\textit{n}-type material) organic material stratified between electrodes,\superfootcite{Tang1986}\superfootcite{Peumans2000}\superfootcite{Halls1996} in these kind of cells the polythiophene made its entrance in organic photovoltaics.\superfootcite{ADMA:ADMA19950070709} 
The problem with these cells arises from the mechanism of charge photo\-generation in organic materials, which occurs as follows: a photon is absorbed, if the energy is large enough a free electron and a hole are formed, otherwise an exciton\nota{\label{fn:exciton}\textit{exciton}, a bound state of an electron and a hole which are attracted to each other by the electrostatic Coulomb force. It is an electrically neutral quasi-particle that exists in insulators, semiconductors and in some liquids. The exciton is regarded as an elementary excitation of condensed matter that can transport energy without transporting net electric charge. (Wikipedia)} 
is formed. 
A sunlight photon has typically enough energy to create a pair of free charges in inorganic semiconductors. Indeed in organic semiconductors the band gap\nota{\label{fn:bandgap}\textit{band gap energy}, the energy difference between the bottom of the conduction band and the top of the valence band or, from a molecular point of view, the difference between \gls{LUMO} and \gls{HOMO}. } 
is higher and depends also on the molecular properties of the material. The so generated excitons have to reach a donor-acceptor materials interface where the high local electric field breaks them in free charges. However, the diffusion length of an exciton in a semiconducting material (5-\SI{10}{\nm}) prior to its dissipative recombination is typically less than the optical absorption length (\SI{\approx50}{\nm}), then a bilayer approach can't have both high absorption and high exciton dissociation efficiencies. 
A possible solution is to greatly increase the interface inter\-dispersing the acceptor and donor materials and obtaining a bulk heterojunction,\superfootcite{Yu1995}\superfootcite{Saunders2008} schematized in Figure~\ref{fig:bhj}. An excessive dispersion could avoid the semi-crystallization in nano\-domains needed for a good charge mobility. Moreover when two opposite charges are formed in the middle of the active layer, they need two continuous pathways to reach the appropriate electrode, thus the presence of segregated domains negatively affects the efficiency, trapping the charges.

\begin{SCfigure}[][tbp]%bhj
\centering
\includegraphics[width=0.5\textwidth]{bhj.png}
\caption{Structure of a bulk heterojunction solar cell.}
\label{fig:bhj}
\end{SCfigure}

\paragraph{Polythiophene in photovoltaics} \Gls{P3HT} in blend with modified fullerene \gls{PCBM} is the reference standard in organic photovoltaics. The role of polythiophene is to absorb the sunlight, transport the generated exciton to the materials interface, donate an electron to \gls{PCBM} and transport the resulting hole towards the anode, while electrons move through \gls{PCBM}. 
A reduced band gap would be highly favorable for absorbing also the low energy portion of the sunlight spectrum. The conjugation length (and then the regioregularity which depends upon the polymerization method) influences the band gap. In recent years many alternating polymers based on thiophene were synthesized with a low band gap. For these polymers a controlled living polymerization method isn't reported yet.

\paragraph{Block copolymers} Unfortunately the materials dispersed in a bulk heterojunction solar cell are usually immiscible, this means that they tend to segregate in two phases over time, reducing the inter\-facial area. In order to compatibilize the acceptor and donor materials, an amphiphilic surfactant can be added.\superfootcite{Liu2012} 
This would introduce in the system a new molecule which could trap charges or act as an insulating layer. An innovative and promising approach is to use an amphiphilic diblock copolymer as donor material, acting as a surfactant itself. The polar block could be designed to complex the \gls{PCBM} (or other inorganic semiconducting nano\-particles) via electron rich sites\superfootcite{Sary2010a}\superfootcite{Topham2011} or can be covalently bonded to the particles. Tuning the blocks volume (changing the degree of polymerization) and the quantity of \gls{PCBM} is a way to optimize the morphology of the blend. 

\end{section}

\begin{section}{Aim of the Thesis}
Many studies on polythiophenes bearing chiral side groups are already present in the literature. 
There are also many articles on diblock copolymers with an achiral thiophenic block and a poly\-(vinyl\-pyridine) block.\superfootcite{Lohwasser2012}\superfootcite{Sary2010a}\superfootcite{Palaniappan2011}\superfootcite{Mougnier2012} In the literature we found also block copolymers made of two polythiophene blocks, one of these being chiral.\superfootcite{VandenBergh2010}\superfootcite{VandenBergh2008}\superfootcite{Verswyvel2011}\superfootcite{Verswyvel2013} A few articles on alternating block copolymers with chirality on the non thiophenic block can be found.\superfootcite{Henze2003} The aim of this thesis is to synthesize and study a polythiophene with a novel chiral side group and its use with a non chiral poly(4-vinyl\-pyridine) block in a diblock copolymer unreported in the literature. 
A modification of the optical activity of the di\-block copolymer with respect to that of polythiophene homo\-polymer could envisage a possible use of this material as a chiral probe in the aggregation process of bulk hetero\-junction solar cells fabricated using block copolymers as active layer.

Thus we started synthesizing a semiconductor, light harvesting chiral polymer with low polydispersity. Then, using an original convenient strategy, we functionalized a termination of the chiral polythiophene with an initiator for the living \gls{NMRP}, finally we polymerized a second block of poly\-(4-vinyl\-pyridine). The obtained diblock copolymer can be used in various devices like solar photovoltaic cells and in studies on supramolecular chirality.

\end{section}
\begin{section}{Introduction to Synthesis and Characterization}

\begin{subsection}{Monomer and Lateral Chains}

\paragraph{Oxygenated side chain} In order to synthesize a chiral polythiophene a monomer bearing a chiral side chain was used. Many different side chains and monomers could fit our requirements. The source of chirality was chosen as an inexpensive and non toxic\superfootcite{McGinty2010} chiral alcohol: (S)-2-methyl-1-butanol also known as active amyl alcohol, present in nature as a component of fusel alcohol which is a by-product of alcoholic fermentation. Our syntheses produced polythiophenes with a (2-methyl\-butyl\-oxy)\-methyl side chain which is, to the best of our knowledge, unreported in the literature. 
Due to the presence of an oxygen atom in the side chain, electronic levels (\acrshort{HOMO} and \acrshort{LUMO}) of the resulting polythiophenes will be different from electronic levels of the well known \gls{P3HT}.\superfootcite{VanBeek2005}\superfootcite{McCullough1993}\superfootcite{Daoust1991} 
The length of the side chain was chosen to be short enough to allow a good charge mobility in a future material for organic electronics.

\paragraph{Enantiomeric excess} For studying the importance of chirality we synthesized two monomers with different enantiomeric excesses: an enantiopure monomer and a scalemic\nota{\textit{scalemic}, non racemic.}
one having 59~\% optical purity.\nota{\label{fn:opticalpurity}\textit{optical purity}, the ratio of the observed optical rotation of a sample consisting of a mixture of enantiomers to the optical rotation of one pure enantiomer. (IUPAC Gold Book)} We used a scalemic alcohol due to the unavailability of the racemate. 

\end{subsection}
\begin{subsection}{Regioselectivity of Monomer Activation}

\paragraph{Grignard activation} The activation process, which is the first step of the polymerization method, consists in a metalation in position 5 of thiophene ring (see Figure~\ref{fig:regioregolarita} for position numbering). As stated in page~\pageref{intro-polimerizzazione}, both zinc
(Rieke method) and magnesium (McCullough method) can be used for this aim. The latter is the most commonly used because allows a more controlled polymerization,\superfootcite{Sheina2004}\superfootcite{Iovu2005}\superfootcite{Yokoyama2004}\superfootcite{Miyakoshi2005} thus we activated our monomer as a Grignard reagent. 
The metalation on iodine was performed with a Grignard metathesis\nota{\textit{metathesis}, a bimolecular process formally involving the exchange of a bond (or bonds) between similar interacting chemical species so that the bonding affiliations in the products are identical (or closely similar) to those in the reactants. (IUPAC Gold Book)} 
(\acrshort{GRIM}), a reaction of replacement of a halogen atom by a metal from an organometallic compound.\superfootcite{Knochel2003}\superfootcite{Tamborski1971}\superfootcite{Zakharkin1964} 
Both Grignard formation\superfootcite{Krasovskiy2004}\superfootcite{Hauk2006} and subsequent polymerization\superfootcite{Takahashi2010}\superfootcite{Lohwasser2011}\superfootcite{Miyakoshi2006}\superfootcite{Ono2012}\superfootcite{Lanni2009}\superfootcite{Lanni2010} are facilitated by the addition of lithium chloride, but the resulting polymers could have lower regioregularity.\superfootcite{Lohwasser2011} 
Tert-butyl magnesium chloride and iso-propyl magnesium chloride are the most common activating Grignard reagents, the former could give a slightly better selectivity\superfootcite{Loewe2001} towards the less sterically demanding 5 position, but iso-propyl magnesium chloride is reported to be sufficiently regioselective when a iodine-bromine monomer is used.\superfootcite{Takahashi2010}\superfootcite{Lohwasser2011}\superfootcite{Ono2012} 
Iso-propyl magnesium chloride could lead to iso-propyl terminated polymers,\superfootcite{Kiriy2011}\superfootcite{Jeffries-EL2004}\superfootcite{Lohwasser2011}\superfootcite{Miyakoshi2005} this eventuality will be verified by Matrix-assisted laser desorption ionization\nota{\textit{matrix-assisted laser desorption ionization}, 
a soft ionization technique used in mass spectrometry, allowing the analysis of large organic molecules. (Wikipedia)}--time-of-flight mass spectrometry.\nota{\textit{time-of-flight mass spectrometry} or TOF MS, is a method of mass spectrometry in which an ion's mass-to-charge ratio is determined via a time measurement. (Wikipedia)}

\begin{SCfigure}[][tbp]%gaussian-grignard
\centering
\begin{subfigure}[b]{0.35\textwidth}
\centering
\includegraphics[width=\textwidth]{gaussian-grignard-2.png}
\subcaption{}
\label{fig:gaussian-grignard-2}
\end{subfigure}%
~%add desired spacing between images, e. g. ~, \quad, \qquad etc.
%(or a blank line to force the subfigure onto a new line)
\begin{subfigure}[b]{0.28\textwidth}
\centering
\includegraphics[width=\textwidth]{gaussian-grignard-5.png}
\caption{}
\label{fig:gaussian-grignard-5}
\end{subfigure}
\caption[Activated monomers.]{Activated monomers. (\subref{fig:gaussian-grignard-2}) Grignard formation in position 2, \ch{MgCl} substitutes \ch{Br}.\\ (\subref{fig:gaussian-grignard-5}) Grignard formation in position 5, \ch{MgCl} substitutes \ch{I} (optimization Gaussian\-09c01\-/DFT\-/B3LYP\-/LanL2DZ).}\label{fig:gaussian-grignard}
\end{SCfigure}

\paragraph{Influence of side chain on activation} In the case of a thiophene monomer with an alkyl side chain the steric effect by the side group is sufficient to obtain 85 \% regioselectivity even using a dibromide monomer.\superfootcite{Loewe2001}\superfootcite{Iovu2005}\superfootcite{Chen1993} 
But when an oxy\-methylene side group is present, using a dibromide monomer leads to regiorandom polymers\superfootcite{Ohshimizu2008}\superfootcite{Zoombelt2008} while using a iodide-bromide monomer leads to regioregular polymers.\superfootcite{Ohshimizu2008}\superfootcite{Zoombelt2008}\superfootcite{Adachi2006}\superfootcite{Locke2010} 
We suppose that this lack of regioselectivity is related to the formation a five-member ring with the oxygen atom coordinating the magnesium atom in position 2, as shown in Figure~\ref{fig:gaussian-grignard-2}, therefore favoring the formation of the undesired 2-magnesio\-chloro-5-iodo-3-[((S)-2-methyl\-butyl\-oxy)\-methyl]\-thio\-phene. This ring formation can't take place with the product of \gls{GRIM} in position 5, as shown in Figure~\ref{fig:gaussian-grignard-5}, with 2-bromo-5-magnesio\-chloro-3-[((S)-2-methyl\-butyl\-oxy)\-methyl]\-thio\-phene.

\end{subsection}
\begin{subsection}{Polymerization Regioregularity}

\begin{SCfigure}[][tbp]%regioregolarita
\centering
\includegraphics[scale=0.7]
{regioregolarita.pdf}
\caption{Regioregular poly-3-alkylthiophene, \textit{head-to-tail} and numbered thiophene positions.}
\label{fig:regioregolarita}
\end{SCfigure}

\paragraph{Importance of regioregularity} In order to obtain a solar light harvesting and semiconducting polymer we chose to synthesize a regioregularly substituted conjugated polythiophene. The substitution with flexible side chains is necessary to process the otherwise insoluble and infusible unsubstituted polythiophene. Undoped polythiophenes usually absorbs the high energy portion of the solar radiation spectrum, thus a red shift in absorption is welcomed and can be obtained by a reduction of the band gap. 
As stated earlier the band gap width is related both to the electronic molecular properties of the side group and to the extension of conjugation along the polymer. This in turn depends on the coplanarity of the aromatic rings that is avoided by repulsion between side groups. 
In conclusion we want to polymerize a thiophene bearing a single side chain in a regioregular\nota{\label{fn:regioregular}\textit{regioregular}, a polymer in which each repeat unit is derived from the same isomer of the monomer. (Wiktionary)}\superfootcite{Kim2006} 
repetition so that the distance between side chains is maximized and no backbone twist is induced by steric repulsions. As indicated in Figure~\ref{fig:triadi}, a way of referring to the relative orientation of two monomeric units is \textit{head-to-head}, \textit{tail-to-tail} or \textit{head-to-tail}, meaning that the \textit{head} of a monomer, the position 2, nearest to the side chain, or the \textit{tail}, the position 5, farthest from the side group, is directly bonded to the \textit{head} or the \textit{tail} of the successive monomer. 
An ordered planar conformation and tight interdigitation in aggregates can be obtained only with a \textit{head-to-tail} polythiophene, thus a good regioregularity is of crucial importance for aggregation and morphology,\superfootcite{Samitsu2008}\superfootcite{Samitsu2008a}\superfootcite{Goto2000}\superfootcite{Goto2002}\superfootcite{McCullough1993c}\superfootcite{McCullogh1995}\superfootcite{Adachi2011}\superfootcite{Xu1993} but it is also important for intra-chain conduction\superfootcite{Basescu1987} and is necessary for supramolecular expression of side chains chirality.\superfootcite{Bouman1994a} 

\begin{SCfigure}[][tbp]%triadi
\centering
\includegraphics[scale=0.6]
{triadi.pdf}
\caption{Possible triads of monosubstituted thiophene rings.}
\label{fig:triadi}
\end{SCfigure}

\begin{figure}[tbp]%iniziazione-meccanismo
\centering
\includegraphics[width=1\textwidth]
{iniziazione-meccanismo.pdf}
\caption{Mechanism of initiation in nickel catalyzed dehalogenative polycondensation.}
\label{fig:iniziazione-meccanismo}
\end{figure} 

\paragraph{Obtaining regioregularity} Regiocontrolled polymerization of thiophene can be achieved with nickel catalyzed dehalogenative polycondensation\nota{\label{fn:polycondensation}\textit{condensation polymer}, the molecular formula of the monomer differs from that of the structural unit. %(\citeauthor*{Carothers1929}\superfootcite{Carothers1929})
} 
starting from a regioselectively activated monomer.\superfootcite{Chen1992}\superfootcite{Chen1995}\superfootcite{McCullough1993a}\superfootcite{Loewe2001}\superfootcite{Loewe1999}\superfootcite{Mao2004}\superfootcite{McCullogh1995}\superfootcite{McCullough1993c} 
This method produces a polymer with high regioregularity even in presence of some monomer activated in the \textit{wrong} position.\superfootcite{Chen1995}\superfootcite{Loewe2001}\superfootcite{Iovu2005}\superfootcite{Lanni2010}\superfootcite{Lohwasser2011} 
The initiation mechanism is illustrated for \acrfull{P3HT} in Figure~\ref{fig:iniziazione-meccanismo}. Initially a transmetalation occurs, two magnesium atoms on two activated thiophenic units are exchanged with the nickel(II) atom of the catalyst. Then a reductive elimination produces a nickel(0) complex and a symmetric dimer that will be present in a chain extremity as a \textit{tail-to-tail} dyad, a defect in regioregularity.\superfootcite{Verswyvel2011}\superfootcite{Tkachov2010}\superfootcite{Miyakoshi2005} 
At this point the nickel(0) complex can perform an oxidative addition in a thiophenic \ch{C-Br} bond starting the propagation stage.
Then, as shown in Figure~\ref{fig:regioregolarita-meccanismo}, on the \ch{CNi-Br} bond a transmetalation occurs (or, from the nickel point of view, an exchange in anions) involving an activated monomer. A subsequent reductive elimination and oxidative addition complete the propagation mechanism. 
The use of a nickel catalyst, rather than a palladium one, and the use of a chelating ligand with wide bite angle are essential conditions for obtaining a good regioselectivity.\superfootcite{Chen1995}\superfootcite{Choi2010}\superfootcite{Mao2004} Indeed we used \acrlong{nidppp} (\acrshort{nidppp}, \acrshort{dppp} = \ch{Ph2P(CH2)3PPh2}) which is the most widely utilized catalyst for regioregular nickel catalyzed dehalogenative polycondensation.

\begin{figure}[tbp]%regioregolarita-meccanismo
\centering
\includegraphics[width=1\textwidth]
{regioregolarita-meccanismo.pdf}
\caption{Mechanism of propagation in nickel catalyzed dehalogenative polycondensation.}
\label{fig:regioregolarita-meccanismo}
\end{figure} 

\paragraph{Evaluating regioregularity} The regioregularity can be qualitatively evaluated by absorption in solution\superfootcite{Chen1995}\superfootcite{McCullough1993c} and quantitatively by \gls{NMR}.\superfootcite{Jeffries-El2005}\superfootcite{Chen1995}\superfootcite{Chen1992}\superfootcite{Zoombelt2008}\superfootcite{Maior1990}\superfootcite{McCullough1993c}
With our uncommon polythiophenes, evaluating the regioregularity from absorption maxima in chloroform solution is not feasible because of the variance of the data reported in the literature.\superfootcite{Higashihara2011}\superfootcite{VanBeek2005}\superfootcite{Ohshimizu2008}\superfootcite{McCullough1993}\superfootcite{Zoombelt2008}\superfootcite{Saito2005}\superfootcite{Locke2010}\superfootcite{Bouman1994} 
At {\HNMR} the aromatic hydrogen peak and the $\alpha$-methylenic peak of a thiophenic unit are influenced by the orientations of nearby repeating units, giving a different peak for each situation showed in Figure~\ref{fig:triadi}. In our case the signals arising from each situation are singlets from both aromatic and aliphatic protons, thus it's easy to distinguish the presence of nearby signals. This is an advantage over \gls{P3HT} in which the peak from the first methylene in the aliphatic side chain is a triplet. 

\end{subsection}
\begin{subsection}{Polymer Weight Control}

\paragraph{Block copolymers} 
As seen in Section~\ref{intro-pv}, 
to stabilize a \acrfull{BHJ} structure (see Figure~\ref{fig:bhj}) against segregation a possible approach is to use block copolymers as compatibilizer or as main component.\superfootcite{Topham2011} In this thesis we'll synthesize an amphiphilic diblock copolymer poly\-thio\-phene-poly\-(4-vinyl\-pyridine). 
The polythiophene block acts as \textit{p}-type material, electron donor. The poly\-(4-vinyl\-pyridine) doesn't act itself as electron acceptor, indeed it complexes modified fullerenes (usually \acrlong{PCBM}, \acrshort{PCBM}, but also non functionalized fullerene\superfootcite{Laiho2006}) which act as \textit{n}-type semiconductor. The optimization of relative volume fraction of the polymer blocks can produce a bicontinuous interpenetrating network of donor and acceptor materials that would give a better charge mobility.

\paragraph{Importance of weight} A low polydispersity helps the crystallization process, thus helping the polymer to give a high interchain charge mobility. Polydispersity is important also for the aggregation morphology of block polymers.\superfootcite{Maria2008}\superfootcite{Koningsveld1972}\superfootcite{Hong1981}\superfootcite{Whitmore1985} Controlling the molecular weight of the polymer is important in this system because changing block dimension can lead to a different blend morphology.\superfootcite{Topham2011} 
Moreover it is important in rod (rigid, conjugated along the backbone) polymers due to their incapacity of chain stretching to accommodate packing of a statistical distribution of chain sizes within an uniformly sized nano\-domain. In brief, high polydispersity could disturb the crystallization needed for a good charge transportation and introduce disorder in the formation of complex morphologies in block copolymers.\superfootcite{Schlaad2004}

\paragraph{Polymerization mechanism} \label{kctp} Most of the conducting conjugated polymers are synthesized by step-growth\nota{\textit{step-growth polymerization}, a polymerization mechanism where any termination on growing oligomers or monomers can react with any other terminations.} poly\-condensation\superfootcite{conj-sint} resulting in high polydispersity polymers. 
In 2004 a quasi-living\nota{\textit{living polymerization}, a chain polymerization from which chain transfer and chain termination are absent. In many cases, the rate of chain initiation is fast compared with the rate of chain propagation, thus that the number of kinetic-chain carriers is essentially constant throughout the polymerization. (IUPAC Gold Book)} 
chain-growth\nota{\label{fn:chain-growth}\textit{chain-growth polymerization}, a polymerization technique where unsaturated monomer molecules add onto the active site on a growing polymer chain one at a time. Growth of the polymer occurs only at one (or possibly more) ends. Addition of each monomer unit regenerates the active site. (Wikipedia)} 
mechanism was reported\superfootcite{Yokoyama2004}\superfootcite{Sheina2004} for nickel catalyzed de\-halogenative polycondensation of \acrfull{P3HT}. 
The chain-growth character is possible thanks to the loose coordination of \ch{Ni}(0) on the growing polythiophene chain forming a non-diffusive associated pair (as illustrated in Figures~\ref{fig:iniziazione-meccanismo} and \ref{fig:regioregolarita-meccanismo}), thus the oxidative insertion step can occur only in a \ch{C-Br} bond on the growing chain.\superfootcite{Tkachov2010}\superfootcite{Zenkina2008}\superfootcite{Miyakoshi2005}\superfootcite{Jeffries-El2005}\superfootcite{Tkachov2010}\superfootcite{Sheina2004} This recent polymerization method is named \gls{KCTP}. 
Up to now it was extended to a few other polymers different from polythiophene while maintaining a living character.\superfootcite{Ono2012}\superfootcite{Wu2010}\superfootcite{Yokoyama2008} Obviously the catalyst is important and again \gls{nidppp} resulted as the best performing for living poly\-merizations of poly(3-alkyl\-thio\-phenes).\superfootcite{Choi2010}\superfootcite{Senkovskyy2009}\superfootcite{Bronstein2009}

\paragraph{Second block polymerization} \Acrfull{P4VP} was chosen as the polar component in our amphiphilic block copolymer. Various block copolymer with a complexing \acrshort{P4VP} block are already present in the literature.\superfootcite{Sary2010a}\superfootcite{Lohwasser2012}\superfootcite{Palaniappan2011}\superfootcite{Mougnier2012}\superfootcite{Laiho2006}\superfootcite{Maria2008}\superfootcite{Sun2011} The molecular weight of the coil (not rigid) block is important too, thus a living polymerization was used: the chain-growth living \acrfull{NMRP}. 
In the present case we chose the grafting-from\nota{\textit{grafting-from polymerization}, the initiator for a polymerization is chemically bonded on a pre-existing polymer.} approach starting the polymerization of poly(4-vinyl\-pyridine) from a macroinitiator bonded on the end of the polythiophene chains. 
A grafting-onto\nota{\textit{grafting-onto polymerization}, two or more formed polymers are chemically joined to form a multi\-block polymer.} 
approach would have permitted a better characterization of the second block, but obtaining a diblock copolymer (rather than a triblock or more) would have required the asymmetric termination of both polymeric blocks. 
A further significant aspect to be taken into account using this approach is the chelating ability of poly(4-vinyl\-pyridine) that would disturb the joining process chelating any metal or catalyst needed for the reaction. 

\end{subsection}
\begin{subsection}{Polymer Terminations Control}
A chain-growth polymerization is the way to synthesize an asymmetrically terminated macroinitiator. Thankfully quasi-living Kumada catalyst-transfer polycondensation can give asymmetric termination of the polythiophene chains\superfootcite{Jeffries-El2005}\superfootcite{Lohwasser2011} and in the literature we can find many diblock copolymers synthesized with a combination of \gls{KCTP}--\gls{NMRP}.\superfootcite{Zhang2009}\superfootcite{Britze2011}\superfootcite{Kaul2010}\superfootcite{Mougnier2012}\superfootcite{Richard2008}

\paragraph{Polymerization quenching} The quenching process is reported to be crucial in order to obtain a quasi-mono\-dispersed polymer and defined terminations. 
Quenching simply with water\superfootcite{Miyakoshi2004} or methanol\superfootcite{Lohwasser2011} can cause dimerization via disproportionation of the poly\-thio\-phene--\ch{Ni}(0)--\ch{Br} complex, followed by reductive elimination of poly\-thio\-phene--poly\-thio\-phene with doubled weight. Quenching with aqueous hydrochloric acid is reported to prevent dimerization.\superfootcite{Miyakoshi2004}\superfootcite{Lohwasser2011} 
Quenching with certain Grignard reagents can even react \textit{in situ} with a single chain end to give asymmetrically functionalized poly\-thio\-phenes.\superfootcite{Jeffries-El2005}\superfootcite{Jeffries-EL2004} We preferred to quench with aqueous hydrochloric acid because it promised for a reliable \ch{H/Br} termination.\superfootcite{Lohwasser2011}

\end{subsection}
\begin{subsection}{Optical Properties}
\paragraph{Band gap} We expect a polythiophene with a rather usual band gap energy (\SI{\approx1.9}{\eV}) but a slightly different \gls{LUMO} 
energy and \gls{HOMO} energy. 
These energy levels are an essential parameter of thin film photovoltaic devices,\superfootcite{Scharber2006}\superfootcite{Cremer2006}\superfootcite{Koster2006} in fact it is well known that open circuit voltage ($\mathrm{V_{oc}}$, which is a generally accepted estimate for the built-in potential) has an upper limit derived from the difference between \gls{HOMO} of the electron donor material and \gls{LUMO} of the electron acceptor material.\superfootcite{Brabec2001}\superfootcite{Gadisa2004} An electron withdrawing substituent in the donor material could move the electronic levels increasing the oxidation potential, this would increase the $\mathrm{V_{oc}}$. 
In the literature on polythiophenes the presence of a \ch{-CH2-O-R} side group is reported as beneficial for various photovoltaic devices\superfootcite{VanBeek2005}\superfootcite{Zoombelt2008} and for the conducting properties of materials.\superfootcite{Bryce1987} 

\paragraph{UV-vis} 
The study of solvato\-chromism reveals the formation of aggregates and the aggregation quality. Both an interaction of chains giving an exciton coupling and an extension of the conjugation length, result in a red-shift of the \gls{UVvis} absorption peak. When a certain amount of non-solvent is added to a solution of polymer, aggregation occurs giving a colloidal dispersion. The conditions for obtaining an aggregate can deeply influence the resulting morphology.\superfootcite{Kiriy2003}\superfootcite{Kim2010}\superfootcite{Goto2002}\superfootcite{Salatelli2010} 
In order to verify the solid state aggregation and morphology, we employed various sample preparation methods: spin coating, drop casting and dispersion in \ch{KCl} disk. 

These three methods should allow us to obtain different semi\-crystallinity extents: spin coating should avoid an organization of polymer chains, drop casting produces both an amorphous mixed with a crystalline solid, \ch{KCl} disks preparation employs dry polymer, thus should result in the most well packed structure. The presence of a vibronic fine structure suggests, but don't assure, the presence of some crystalline domains in the material.

\end{subsection}
\begin{subsection}{Photoluminescence and Chirality}
\paragraph{Photoluminescence} One of the aims of introducing chirality in our polymer was to influence the aggregation morphology, avoiding a parallel chain stacking, in order to reduce exciton quenching related to inter chain coupling. We suppose that studying photoluminescence quenching can help us to understand also exciton quenching.\superfootcite{Theander1999}\superfootcite{Kaneto1989} It was already reported that an aggregation morphology in which chains adopt a non parallel relative orientation results in highly fluorescent materials.\superfootcite{Zahn2002}\superfootcite{Deans2000}%
\superfootcite{Bredas1999}

\paragraph{Chirality as a probe} Talking of bulk heterojunction solar cells we are referring to soft-matter, quite disordered, difficult to characterize in the solid state. As stated in page~\pageref{intro-cd}, introducing some extent of chirality in our material allows us to use circular dichroism spectroscopy to characterize also aggregated states. 
In an assembly with chirally oriented identical transition dipole moments in close spatial proximity, their interaction brings about a splitting, called Davydov splitting,\nota{\textit{Davydov splitting}, the splitting of bands in the electronic or vibrational spectra of crystals due to the presence of more than one (interacting) equivalent molecular entity in the unit cell. (IUPAC Gold Book)} %
which is reflected in two distinct \gls{CD} bands 
with opposite rotational strength.\superfootcite{Berova2007} 
The sum of these two Cotton signals\nota{\label{fn:cottoneffect}\textit{Cotton effect}, the characteristic change in optical rotatory dispersion and/or circular dichroism in the vicinity of an absorption band of a substance. In a wavelength region where the light is absorbed, the absolute magnitude of the optical rotation at first varies rapidly with wavelength, crosses zero at absorption maxima and then again varies rapidly with wavelength but in opposite direction. (Wikipedia)} 
is a bisignate couplet. 
In polythiophenes the chromophore, for $\pi-\pi^*$ transitions, is located on the conjugated backbone. We can imagine the chromophore as a segment of the chain, as long as the conjugation length. Usually bisignate couplet signals arise from interchain interactions\superfootcite{Langeveld-Voss1999} but seldom also from intra-chain interactions.\superfootcite{Kiriy2003}\superfootcite{Nilsson2004}\superfootcite{Langeveld-Voss1996} The sign of the couplet can indicate the relative orientation of the chromophores.\superfootcite{Langeveld-Voss1998} 

\paragraph{Chirality in good solvent} Chiral polythiophenes in a good solvent usually exhibit only very weak electronic circular dichroism\superfootcite{Salatelli2010} because when the chain is well solvated chirality in side groups doesn't influence the $\pi-\pi^*$ transitions of the backbone. 

\paragraph{Chirality in poor solvent} When the conjugated chains interact with each other, the side groups rotational freedom is reduced and their chirality influences the organization of the material via steric repulsion. Adding a poor solvent or a non solvent to a solution of polymer, we can observe the rising of an induced circular dichroism during the formation of aggregated states; this can be caused by main chain twists or by a supramolecular chiral organization. 
The percent of non solvent can have a non trivial effect on the obtained aggregate\superfootcite{Bidan1996}\superfootcite{Goto2000} which can be studied via circular dichroism spectroscopy. Moreover changing solvent and/or non solvent could change dramatically the kind of aggregation and the \gls{CD} spectra of the supramolecular assembly.\superfootcite{Salatelli2010}\superfootcite{Goto2002}\superfootcite{Langeveld-Voss1998a} We employed mixtures of good and poor solvents like: chloroform-methanol, chloroform-acetonitrile, tetrahydrofuran-methanol and tetrahydrofuran-acetonitrile.

\end{subsection}
\end{section}

\chapter{Synthesis and Characterization}
\graphicspath{ {./img/results/} }
\begin{section}{Monomer Synthesis and Characterization}
Two monomers were chosen for their simple and inexpensive synthesis: 
2-bromo-5-iodo-3-[2-((S)-2-methyl\-butyl\-oxy)\-ethyl]\-thio\-phene and 2-bromo-5-iodo-3-[((S)-2-methyl\-butyl\-oxy)\-methyl]\-thio\-phene \cmpd+{ig2-2}. During the synthesis of the first monomer we encountered some problems thus we abandoned the former in favor of the latter monomer, which promised an effortless synthesis.
\begin{subsection}{Side Chain}

\begin{SCfigure}[][tbp]%syn1-eterificazione
\centering
\includegraphics[scale=0.7]
{syn1-eterificazione.pdf}
\caption{Etherification reaction.}
\label{fig:syn1-eterificazione}
\end{SCfigure}

Enantiopure (S)-2-methyl-1-butanol and scalemic 2-methyl-1-butanol alcohols were deprotonated in tetrahydrofuran using one equivalent of sodium hydride, then their alkoxides were reacted in a nucleophilic substitution with commercial 
2-bromo-3-(bromo\-methyl)\-thio\-phene giving 2-bromo-3-[(2-methyl\-butyl\-oxy)\-methyl]\-thio\-phene \cmpd+{ig2-9} and 2-bromo-3-[((S)-2-methyl\-butyl\-oxy)\-methyl]\-thio\-phene \cmpd+{ig2-1}. Refer to Section~\ref{sec:ig2-9} for complete experimental procedure.\superfootcite{Lee2011} Purity was confirmed by \gls{TLC} and \gls{GCMS}. 
The electronic effect of the oxy\-genated side group on the thio\-phene ring was estimated via computational calculations. Comparing molecular orbital energies of 3-propyl\-thio\-phene with 3-(meth\-oxy\-methyl)\-thio\-phene, the oxy\-gen atom resulted to have an electron withdrawing effect on the thio\-phenic ring (optimization and energy calculation Gaussian09c01\-/DFT\-/B3LYP\-/cc-pVDZ). This is confirmed by cyclic voltammetry on similar polymers reported in the literature.\superfootcite{Zoombelt2008}

\end{subsection}
\begin{subsection}{Iodination}

\begin{SCfigure}[][tbp]%syn2-iodurazione
\centering
\includegraphics[width=0.95\textwidth]
{syn2-iodurazione.pdf}
\caption{Regioselective iodination.}
\label{fig:syn2-iodurazione}
\end{SCfigure}
The chosen polymerization strategy requires a iodinated-brominated thio\-phene monomer, thus the previously synthesized ether was iodinated selectively in position 5. 

In a dichloromethane solution of 2-bromo-3-[(2-methyl\-butyl\-oxy)\-methyl]\-thio\-phene \cmpd+{ig2-9} or 2-bromo-3-[((S)-2-methyl\-butyl\-oxy)\-methyl]\-thio\-phene \cmpd+{ig2-1}, iodine and iodo\-benzene di\-acetate were added to obtain the desired monomer 2-bromo-5-iodo-3-[((S)-2-methyl\-butyl\-oxy)\-methyl]\-thio\-phene \cmpd+{ig2-10} or 2-bromo-5-iodo-3-[((S)-2-methyl\-butyl\-oxy)\-methyl]\-thio\-phene \cmpd+{ig2-2}, as indicated in Figure~\ref{fig:syn2-iodurazione}. 
Refer to Section~\ref{sec:ig2-10} for complete experimental procedure.\superfootcite{Locke2010} The structure was verified by {\HNMR}, {\CNMR} and \gls{FTIR} characterization. Partial confirmation of {\HNMR} and {\CNMR} data can be found in \citeauthor*{Locke2010}\superfootcite{Locke2010} looking for molecule S6. Purity was confirmed by \gls{TLC} and \gls{GCMS}.

The quantities of iodine and iodo\-benzene di\-acetate are crucial to obtain only the mono\-iodinated product. 
The presence of a iodine atom allows us to activate selectively the desired position without competition from the bromine atom, as we'll see in the next section. 
\end{subsection}
\begin{subsection}{Metalation Reaction}
\label{sec:monomero-regioselettivita}

\begin{SCfigure}[][tbp]%syn3-attivazione
\centering
\includegraphics[scale=0.7]
{syn3-attivazione.pdf}
\caption{Metalation activation reaction.}
\label{fig:syn3-attivazione}
\end{SCfigure}
In order to evaluate the regioselectivity of Grignard metathesis we reacted the monomer \cmpd+{ig2-10} with one equivalent of isopropyl magnesium bromide, both in presence and in absence of lithium chloride, and quenched with a solution of hydrochloric acid. Gas chromatography--mass spectrometry (GCMS) analysis showed a complete and regioselective hydrolysis when no lithium chloride is present.
If \ch{LiCl} is added, the product, after washing and extraction, appears to contain a small quantity of immiscible byproduct. \Gls{GCMS} analysis showed a decrease in regio and chemioselectivity. Chromatograms are reported in appendix.

\end{subsection}
\end{section}
\clearpage
\begin{section}{Polymer Synthesis}
For the polymerization of poly\{3-[((S)-2-methyl\-butyl\-oxy)\-methyl]\-thio\-phene\} \cmpd+{ig2-4} we followed the procedure by \citeauthor*{Miyakoshi2005}\superfootcite{Miyakoshi2005} designed for the polymerization of \gls{P3HT}. 

\begin{figure}[tbp]%syn4-polimerizzazione
\centering
\includegraphics[width=1\textwidth]
{syn4-polimerizzazione.pdf}
\caption{Polymerization method.}
\label{fig:syn4-polimerizzazione}
\end{figure}

Polythio\-phene \cmpd+{ig2-4} was synthesized without addition of lithium chloride. To 2-bromo-5-iodo-3-[((S)-2-methyl\-butyl\-oxy)\-methyl]\-thio\-phene \cmpd+{ig2-2} in dry tetrahydrofuran at \SI{0}{\celsius} under nitrogen atmosphere one equivalent of isopropyl-magnesium chloride was added dropwise. After one hour a suspension of the catalyst \acrfull{nidppp} was poured in reaction mixture and the polymerization remained at room temperature for a day. The reaction was quenched adding \SI{5}{\Molar} hydrochloric acid, resulting in a low weight polymer (see Figure~\ref{fig:ig2-4-ig2-8-sec-chcl3} and Table~\ref{tab:sec}, 
$M_n\approx$~\SI{3.4}{\kg\per\mole} thus $\approx 20$ repeating units).

Afterward we changed the polymerization conditions by following the procedure of 
\citeauthor*{Higashihara2011}\superfootcite{Higashihara2011} designed for the synthesis of poly\-thio\-phene with triethylene glycol side chains to prepare poly\{3-[(2-methyl\-butyl\-oxy)\-methyl]\-thio\-phene\} \cmpd+{ig2-15} and poly\{3-[((S)-2-methyl\-butyl\-oxy)\-methyl]\-thio\-phene\} (\cmpd+{ig2-8}, same repeating unit of \cmpd+{ig2-4}). This procedure is the most common one for monomers bearing hydrophilic oxyethylene side chains,\superfootcite{VanBeek2005}\superfootcite{Hammer2011}\superfootcite{Ohshimizu2008}\superfootcite{Zoombelt2008}\superfootcite{McCullough1993} but procedures with different phosphine ligand are also reported.\superfootcite{Adachi2006}\superfootcite{Yokozawa2008} 
Polymers \cmpd+{ig2-15} and \cmpd+{ig2-8} were synthesized starting respectively from 2-bromo-5-iodo-3-[(2-methyl\-butyl\-oxy)\-methyl]\-thio\-phene \cmpd+{ig2-10} and 2-bromo-5-iodo-3-[((S)-2-methyl\-butyl\-oxy)\-methyl]\-thio\-phene \cmpd+{ig2-2} in presence of lithium chloride. These syntheses resulted in medium weight polymers (see Figure~\ref{fig:ig2-4-ig2-8-sec-chcl3} and Table~\ref{tab:sec}, $M_n\approx$~\SI{10}{\kg\per\mole}, $\approx 50$ repeating units) thus \gls{NMR} signals from initial \textit{tail-to-tail} dyad and from terminations should be negligible if compared with main peaks. 
\cmpd+{ig2-15} and \cmpd+{ig2-8} were carefully purified via various Soxhlet extractions in order to remove catalyst and possible oligomer impurities (see experimental procedure at page~\pageref{sec:ig2-15}).

\begin{subsection}{Regioregularity}
\label{sec:regioregolarita}

As shown in Figure~\ref{fig:triadi} a thio\-phenic unit can be in four different environments, each giving different \gls{NMR} signals. In a regioregular \textit{head-to-tail} poly\-thio\-phene, every backbone ring is in the same environment. Instead a regiorandom polymer would give rise to four different singlets for aromatic protons.\superfootcite{Jeffries-El2005}\superfootcite{Chen1995}\superfootcite{Chen1992}\superfootcite{Zoombelt2008}\superfootcite{Maior1990}\superfootcite{McCullough1993c}

\begin{SCfigure}[][tbp]%ig2-15-nmr-h
\includegraphics[width=0.9\textwidth]{ig2-15-nmr-h.pdf}
\caption[{\HNMR} of polymer \cmpd+{ig2-15}.]{{\HNMR} (\SI{600}{\MHz}) of polymer \cmpd+{ig2-15} in \gls{TCE}-d2.}
\label{fig:ig2-15-nmr-h}
\end{SCfigure}

\begin{figure}[tbp]%regioregolarita-nmr
\parbox[b]{0.47\linewidth}{
\begin{subfigure}[b]{0.47\textwidth}
\parbox[b]{0\linewidth}{\subcaption{}\label{fig:regioregolarita-nmr-ig2-4}}
\includegraphics[width=1\textwidth]{regioregolarita-nmr-ig2-4.pdf}
\end{subfigure}}
\qquad
\parbox[b]{0.47\linewidth}{
\begin{subfigure}[b]{0.47\textwidth}
\parbox[b]{0\linewidth}{\subcaption{}\label{fig:regioregolarita-nmr-ig2-4-c}}
\includegraphics[width=1\textwidth]{regioregolarita-nmr-ig2-4-c.pdf}
\end{subfigure}
}

\parbox[b]{0.47\linewidth}{
\begin{subfigure}[b]{0.47\textwidth}
\parbox[b]{0\linewidth}{\subcaption{}\label{fig:regioregolarita-nmr-ig2-15}}
\includegraphics[width=1\textwidth]{regioregolarita-nmr-ig2-15.pdf}
\end{subfigure}}
\qquad
\parbox[b]{0.47\linewidth}{
\begin{subfigure}[b]{0.47\textwidth}
\parbox[b]{0\linewidth}{\subcaption{}\label{fig:regioregolarita-nmr-ig2-15-c}}
\includegraphics[width=1\textwidth]{regioregolarita-nmr-ig2-15-c.pdf}
\end{subfigure}
}

\parbox[b]{0.47\linewidth}{
\begin{subfigure}[b]{0.47\textwidth}
\parbox[b]{0\linewidth}{\subcaption{\vspace{15pt}}\label{fig:regioregolarita-nmr-ig2-8}}
\includegraphics[width=1\textwidth]{regioregolarita-nmr-ig2-8.pdf}
\end{subfigure}}
\qquad
\parbox[b]{0.47\linewidth}{
\begin{subfigure}[b]{0.47\textwidth}
\parbox[b]{0\linewidth}{\subcaption{\vspace{15pt}}\label{fig:regioregolarita-nmr-ig2-8-c}}
\includegraphics[width=1\textwidth]{regioregolarita-nmr-ig2-8-c.pdf}
\end{subfigure}
}

\caption[Part of {\HNMR} and {\CNMR} spectra in d2-TCE.]{Part of {\HNMR} and {\CNMR} spectra in d2-\gls{TCE}. 
(\subref{fig:regioregolarita-nmr-ig2-4}) and (\subref{fig:regioregolarita-nmr-ig2-4-c}) spectra of low weight polymer \cmpd+{ig2-4}; (\subref{fig:regioregolarita-nmr-ig2-15}) and (\subref{fig:regioregolarita-nmr-ig2-15-c}) spectra of polymer \cmpd+{ig2-15}; (\subref{fig:regioregolarita-nmr-ig2-8}) and (\subref{fig:regioregolarita-nmr-ig2-8-c}) spectra of polymer \cmpd+{ig2-8}.}\label{fig:regioregolarita-nmr}
\end{figure}

{\HNMR} and {\CNMR} spectra of \cmpd+{ig2-4} show many secondary peaks (see Figures~\ref{fig:regioregolarita-nmr-ig2-4} and \ref{fig:regioregolarita-nmr-ig2-4-c}), some of these peaks will be assigned to the end groups of the polymer, see Section~\ref{sec:sintesi-terminazione}. 

{\HNMR} spectra from polymers \cmpd+{ig2-15} and \cmpd+{ig2-8} (see Figures~\ref{fig:regioregolarita-nmr-ig2-15}, \ref{fig:regioregolarita-nmr-ig2-8}, \ref{fig:ig2-15-nmr-h} and appendix) show only a singlet for aromatic hydrogen at \SI{7.21}{\ppm}.

In the region of aliphatic protons in {\HNMR} spectrum from \cmpd+{ig2-15} (see Figures~\ref{fig:regioregolarita-nmr-ig2-15} and \ref{fig:ig2-15-nmr-h}) we can see only main peaks while in {\HNMR} spectrum from \cmpd+{ig2-8} (see Figures~\ref{fig:regioregolarita-nmr-ig2-8} and appendix%
) there are small secondary peaks indicating a lower molecular weight or a higher percent of hydrogen substituted chain ends. For a full assignment of {\HNMR} and {\CNMR} signals refer to Table~\ref{tab:ig2-4-nmr2d}.
{\CNMR} spectra of these polymers (see Figures~\ref{fig:regioregolarita-nmr-ig2-15-c} and \ref{fig:regioregolarita-nmr-ig2-8-c}, complete spectra in appendix) show only four resonances in the aromatic region, corresponding to the four carbons of the thio\-phene ring repeating unit with \textit{head-to-tail} orientation,\superfootcite{Maior1990}\superfootcite{McCullough1993} thus the regioregularity seems complete.\superfootcite{McCullough1993c} 
Some very small peaks from impurities, noticeable in Figure~\ref{fig:regioregolarita-nmr-ig2-8-c}, are present also in Figure~\ref{fig:regioregolarita-nmr-ig2-4-c} (corresponding to peaks reported in Table~\ref{tab:ig2-4-nmr2d-10}).
\end{subsection}
\begin{subsection}{Polydispersity and Molecular Weight}
\label{sec:weight}

\afterpage{\clearpage
\begin{figure}[tbhp]%ig2-4-ig2-8-sec-chcl3
\centering
\begin{tikzpicture}
\begin{axis}[axis x line=bottom,axis y line=left,xlabel=Slice Log MW,ylabel=dwt/d(logM),width=1\textwidth,height=6cm,cycle list name=linestyles*]
\addplot file {img/spectra/ig2-4-sec-chcl3.txt};
\addlegendentry{\cmpd+{ig2-4}};
\addplot file {img/spectra/ig2-8-sec-chcl3.txt};
\addlegendentry{\cmpd+{ig2-8}};
\end{axis}
\end{tikzpicture}
\caption[SEC chromatograms at \SI{35}{\celsius} in \ch{CHCl3} of polymer \cmpd+{ig2-4} and \cmpd+{ig2-8}.]{SEC chromatograms at \SI{35}{\celsius} in \ch{CHCl3} of polymer \cmpd+{ig2-4}: $M_{peak}=$~\SI{3.9}{\kg\per\mole}, $M_n=$~\SI{3.4}{\kg\per\mole}, $M_w=$~\SI{4.4}{\kg\per\mole}, $M_z=$~\SI{5.7}{\kg\per\mole}, $M_{z+1}=$~\SI{7.1}{\kg\per\mole}, $M_w/M_n=$~1.32, $M_z/M_w=$~1.29 and of polymer \cmpd+{ig2-8}: $M_{peak}=$~\SI{12.9}{\kg\per\mole}, $M_n=$~\SI{10.6}{\kg\per\mole}, $M_w=$~\SI{14.8}{\kg\per\mole}, $M_z=$~\SI{20.2}{\kg\per\mole}, $M_{z+1}=$~\SI{26.7}{\kg\per\mole}, $M_w/M_n=$~1.39, $M_z/M_w=$~1.37. Conditions for recording both chromatograms were the same.}
\label{fig:ig2-4-ig2-8-sec-chcl3}
\end{figure}
\begin{table}[tbhp]%sec
\centering
\caption[Data from size exclusion chromatograms of polymers \cmpd+{ig2-4}, \cmpd+{ig2-15} and \cmpd+{ig2-8}.]{Data from size exclusion chromatograms of polymers \cmpd+{ig2-4}, \cmpd+{ig2-15} and \cmpd+{ig2-8} in chloroform and ortho-di\-chloro\-benzene (o-\textsmaller{DCB}).}\label{tab:sec}
\begin{tabular}{c|c|c|c|c}
\toprule
Polymer & Method & $M_n$ (\SI{}{\g\per\mol}) & $M_w$ (\SI{}{\g\per\mol}) & $M_w/M_n$\\\cmidrule{1-5}
\cmpd+{ig2-4}\textsuperscript{a} & \ch{CHCl3}, \SI{35}{\celsius} & 3359 & 4420 & 1.32 \\
\cmpd+{ig2-8}\textsuperscript{a} &\ch{CHCl3}, \SI{35}{\celsius} & 10550 & 14873 & 1.41 \\
\cmpd+{ig2-8}\textsuperscript{a} & o-\textsmaller{DCB}, \SI{145}{\celsius} & 7556 & 12511 & 1.66 \\
\cmpd+{ig2-15}\textsuperscript{a} & o-\textsmaller{DCB}, \SI{145}{\celsius} & 7700 & 13351 & 1.73 \\
\cmpd+{ig2-4}\textsuperscript{b} & \acrshort{THF}, \SI{35}{\celsius} & 3590 & 4620 & 1.29 \\
\cmpd+{ig2-8}\textsuperscript{b} & \acrshort{THF}, \SI{35}{\celsius} & 10825 & 14335 & 1.32 \\
\cmpd+{ig2-15}\textsuperscript{b} & \acrshort{THF}, \SI{35}{\celsius} & 14800 & 20130 & 1.36 \\
\bottomrule
\end{tabular}
\smallskip 

{\textsuperscript{a} measured in Milan; \textsuperscript{b} measured in Catania.}
\end{table}
}

It is known that the use of a folded random coil polystyrene as a standard for size exclusion chromatography leads to an overestimation of the polymeric weight of our rigid rod poly\-thio\-phenes.\superfootcite{Wong2012} 

After purification of \cmpd+{ig2-4} the yield was low (maybe due to losses during purification) and the resulting polymer was rather an oligomer ($\mathrm{M_n}\approx$~\SI{3.3}{\kg\per\mole}, corresponding to $\approx18$ repeating units, polydispersity index $M_w/M_n=$~1.32), as seen by \gls{SEC} in Figure~\ref{fig:ig2-4-ig2-8-sec-chcl3}. This shortness could be caused by the presence of an oxy\-gen atom in our side chain. This oligomer was used for \gls{NMR2D} characterization reported in Section~\ref{sec:terminazioni-nmr2d}.

According to its bi\-modal peak in size-exclusion chromatogram reported in Figure~\ref{fig:ig2-4-ig2-8-sec-chcl3}, polymer \cmpd+{ig2-8} contains two families of polymers: the most abundant one with molar mass~\SI{\approx 13}{\kg\per\mole} and the minor one with molar mass~\SI{\approx 27}{\kg\per\mole}. This indicates the presence of a dimerization reaction during the overnight polymerization. 
We can exclude a dimerization during the quenching performed with aqueous hydrochloric acid.\superfootcite{Miyakoshi2004}\superfootcite{Lohwasser2011} Size-exclusion chromatography at \SI{145}{\celsius} in ortho-di\-chloro\-benzene was performed as a check (see Table~\ref{tab:sec}). The resulting polydispersity index (\acrshort{PDI}) measured in chloroform $M_w/M_n = 1.39$ is better than the reported \gls{PDI} $M_w/M_n=1.68$ of a poly\-thio\-phene polymerized with \gls{GRIM}\-/\gls{KCTP} and a very similar side chain.\superfootcite{Zoombelt2008}
Number average molecular weight obtained in our synthesis ($M_n =$~\SI{10.6}{\kg\per\mole}) corresponds to around 58 repeating units and is compatible with the monomer to catalyst feed ratio (45:1) but lower than the molar mass from the reference article ($M_n =$~\SI{16.5}{\kg\per\mole}).\superfootcite{Higashihara2011} 
By the way, the resulting molar mass is higher than $M_n$ of \cmpd+{ig2-4}, we attribute this enhancement to the addition of lithium chloride. The yield in \cmpd+{ig2-8} (estimated from recovered weight after purification) is $\approx 78~\%$. 

Results from \gls{SEC} at room temperature in tetrahydrofuran of polymer \cmpd+{ig2-15} indicate a higher molecular weight than for polymer \cmpd+{ig2-8} and a broader polydispersity, this is in accordance with data obtained from \gls{SEC} characterizations at \SI{145}{\celsius} in ortho-di\-chloro\-benzene (see Table~\ref{tab:sec}). In the chromatograms recorded in \gls{THF} we can see the presence of peaks at higher molecular weight related to aggregated chains, thus confirming the \gls{THF} as a non perfect solvent. 
The yield in \cmpd+{ig2-15} (estimated from recovered weight after purification) is $\approx 61~\%$.

\end{subsection}
\begin{subsection}{Chain Termini}
\label{sec:sintesi-terminazione}

Quenching a living polymerization with \SI{5}{\Molar} aqueous hydrochloric acid should inactivate the catalyst and insert a hydrogen atom in the growing end,\superfootcite{Lohwasser2011} as shown in Figure~\ref{fig:syn5-quenching}.

\afterpage{\clearpage
\begin{SCfigure}%maybe here I should put a [H]
\centering
\includegraphics[width=0.2\textwidth]{monomero-numerato.pdf}\caption{Our repeating unit with numbering.}\label{fig:monomero-numerato}
\end{SCfigure}
\begin{table}%maybe here I should put a [H]
\centering
\caption[Groups of {\HNMR} and {\CNMR} signals of \cmpd+{ig2-4} from NMR-2D spectra.]{Groups of {\HNMR} and {\CNMR} signals of \cmpd+{ig2-4} from \gls{NMR2D} spectra. 
}\label{tab:ig2-4-nmr2d}
\begin{subtable}[]{0.49\textwidth}
\centering
\caption{Relative intensity: 1}\label{tab:ig2-4-nmr2d-main}
\begin{tabular}{l|c|l|c}
\toprule
\ch{H} & $\delta$ (ppm) & \ch{C} & $\delta$ (ppm)\\\cmidrule{1-4}
& & {{\ch{C}$_{\textrm 5}$}} & 132.91 \\
& & {{\ch{C}$_{\textrm 2}$}} & 133.68 \\
& & {{\ch{C}$_{\textrm 3}$}} & 136.54 \\
{{\ch{H}$_{\textrm 15}$}} & 7.22 & {{\ch{C}$_{\textrm 4}$}} & 129.28 \\
{{\ch{H}$_{\textrm 16, 17}$}} & 4.54 & {{\ch{C}$_{\textrm 6}$}} & 66.71 \\
{{\ch{H}$_{\textrm 18, 19}$}} & 3.39, 3.31 & {{\ch{C}$_{\textrm 8}$}} & 75.88 \\
{{\ch{H}$_{\textrm 20, 21, 22}$}} & 0.92 & {{\ch{C}$_{\textrm 12}$}} & 16.73 \\
{{\ch{H}$_{\textrm 23}$}} & 1.69 & {{\ch{C}$_{\textrm 9}$}} & 34.86 \\
{{\ch{H}$_{\textrm 24, 25}$}} & 1.15, 1.46 & {{\ch{C}$_{\textrm 10}$}} & 26.25 \\
{{\ch{H}$_{\textrm 26, 27, 28}$}} & 0.88 & {{\ch{C}$_{\textrm 11}$}} & 11.38 \\
\bottomrule
\end{tabular}

\end{subtable}
\begin{subtable}[]{0.49\textwidth}
\centering
\caption{Relative intensity: 0.18}\label{tab:ig2-4-nmr2d-18}
\begin{tabular}{l|c|l|c}
\toprule
\ch{H} & $\delta$ (ppm) & \ch{C} & $\delta$ (ppm)\\\cmidrule{1-4}
{{\ch{H}$_{\textrm 13}$}} & 7.36 $J=5.2$ & {{\ch{C}$_{\textrm 5}$}} & 126.33 \\
& & {{\ch{C}$_{\textrm 2}$}} & 129.60? \\
& & {{\ch{C}$_{\textrm 3}$}} & 139.02?\\
{{\ch{H}$_{\textrm 15}$}} & 7.16 $J=5.1$ & {{\ch{C}$_{\textrm 4}$}} & 128.92 \\
{{\ch{H}$_{\textrm 16, 17}$}} & 4.38 & {{\ch{C}$_{\textrm 6}$}} & 66.41 \\
{{\ch{H}$_{\textrm 18, 19}$}} & 3.16, 3.24 & {{\ch{C}$_{\textrm 8}$}} & 75.72 \\
{{\ch{H}$_{\textrm 20, 21, 22}$}} & 0.85 & {{\ch{C}$_{\textrm 12}$}} & 16.58 \\
{{\ch{H}$_{\textrm 23}$}} & 1.60 & {{\ch{C}$_{\textrm 9}$}} & 34.76 \\
{{\ch{H}$_{\textrm 24, 25}$}} & 1.38, 1.06 & {{\ch{C}$_{\textrm 10}$}} & 26.16 \\
{{\ch{H}$_{\textrm 26, 27, 28}$}} & 0.88 & {{\ch{C}$_{\textrm 11}$}} & 11.32 \\
\bottomrule
\end{tabular}
\end{subtable}
\bigskip

\bigskip

\begin{subtable}[]{0.49\textwidth}
\centering
\caption{Relative intensity: 0.16}\label{tab:ig2-4-nmr2d-16}
\begin{tabular}{l|c|l|c}

\toprule
\ch{H} & $\delta$ (ppm) & \ch{C} & $\delta$ (ppm)\\\cmidrule{1-4}
& & {{\ch{C}$_{\textrm 3, 2}$}} & \tallcell{139.50,\\ 130.76} \\
{{\ch{H}$_{\textrm 15}$}} & 7.24 & {{\ch{C}$_{\textrm 4}$}} & 127.71 \\
{{\ch{H}$_{\textrm 16, 17}$}} & 4.33 & {{\ch{C}$_{\textrm 6}$}} & 66.45 \\
{{\ch{H}$_{\textrm 18, 19}$}} & 3.16, 3.24 & {{\ch{C}$_{\textrm 8}$}} & 75.72 \\
{{\ch{H}$_{\textrm 20, 21, 22}$}} & 0.85 & {{\ch{C}$_{\textrm 12}$}} & 16.58 \\
{{\ch{H}$_{\textrm 23}$}} & 1.60 & {{\ch{C}$_{\textrm 9}$}} & 34.76 \\
{{\ch{H}$_{\textrm 24, 25}$}} & 1.38, 1.06 & {{\ch{C}$_{\textrm 10}$}} & 26.16 \\
{{\ch{H}$_{\textrm 26, 27, 28}$}} & 0.88 & {{\ch{C}$_{\textrm 11}$}} & 11.32 \\
\bottomrule
\end{tabular}
\end{subtable}
\begin{subtable}[]{0.49\textwidth}
\centering
\caption{Relative intensity: $\approx0.10$}\label{tab:ig2-4-nmr2d-10}
\begin{tabular}{l|c|l|c}
\toprule
\ch{H} & $\delta$ (ppm)& \ch{C} & $\delta$ (ppm)\\\cmidrule{1-4}
& & {{\ch{C}$_{\textrm 2, 3, 5}$}} & \tallcell{134.59,\\ 132.20,\\ 136.77?} \\
{{\ch{H}$_{\textrm 15}$}} & 7.19 & {{\ch{C}$_{\textrm 3}$}} & 126.81 \\
{{\ch{H}$_{\textrm 16, 17}$}} & 4.52 & {{\ch{C}$_{\textrm 6}$}} & 66.71\\
\bottomrule
\end{tabular}
\end{subtable}
\end{table}
\clearpage
}

In the aromatic region of {\HNMR} spectrum from \cmpd+{ig2-4} we can see (see Figure~\ref{fig:regioregolarita-nmr-ig2-4}%
) two doublets at 7.15 and \SI{7.35}{\ppm} with scalar coupling $J=$~\SI{5.2}{\Hz}. The presence of two hydrogen atoms on a thio\-phene ring bearing also a side group, implies a hydrogen terminated chain end. An hydrogen in the $\alpha$ position on a thio\-phene ring is expected to be generated by the quenching process of the living chain; this reaction should produce a thio\-phene ring with the 2 and 4 positions (for position numbering refer to Figure~\ref{fig:regioregolarita}) substituted with hydrogen atoms. 
Looking at the amplitude of the $J$-coupling we can verify the relative position of these two protons. 
\label{h-terminale} From the literature on thio\-phene we know that the width of the $J$-coupling for two non vicinal hydrogen atoms in position 2 and 4 is $J\approx$~\SI{1.4}{\Hz}.\superfootcite{Miyakoshi2005} The observed $J$-coupling of \SI{5.2}{\Hz} indicates\superfootcite{Zoombelt2009} the position of hydrogen atoms being vicinal, thus positions 4 and 5. This means that a reversed regioregularity in the ending thio\-phene ring was present. 

\begin{SCfigure}[][tbp]%syn5-quenching
\centering
\includegraphics[scale=0.6]
{syn5-quenching.pdf}
\caption{Quenching process.}
\label{fig:syn5-quenching}
\end{SCfigure}

\label{sec:terminazioni-nmr2d}
Analysis of \gls{NMR2D} maps ($^1$\ch{H}-COSY;\nota{\textit{COSY}, two-dimensional correlation spectroscopy.} $^1$\ch{H}-TOCSY;\nota{\textit{TOCSY}, two-dimensional total correlation spectroscopy, also called HOHAHA (homo\-nuclear Hartmann Hahn).} 
$^1$\ch{H},$^{13}$\ch{C}-HSQC;\nota{\textit{HSQC}, hetero\-nuclear single quantum coherence spectroscopy.} $^1$\ch{H},$^{13}$\ch{C}-HMBC\nota{\textit{HMBC}, hetero\-nuclear multiple bond correlation.}) confirms this hypothesis and allows us to distinguish four groups of signals related to four thio\-phenic units present in \cmpd+{ig2-4}. 
The most abundant unit is assigned to \textit{head-to-tail} repetition and is reported in Table~\ref{tab:ig2-4-nmr2d-main}. The previously identified \ch{H} terminated \textit{tail} ending is reported in Table~\ref{tab:ig2-4-nmr2d-18}. Other groups of signals could be due to different terminations or regioregularity defects.

The presence of this \ch{H} terminated \textit{tail} monomer in the chain end could be due to a partial formation of the unintended regioisomer 2-magnesio\-chloro-5-iodo-3-[((S)-2-methyl\-butyl\-oxy)\-methyl]\-thio\-phene (see Section~\ref{sec:monomero-regioselettivita}). This isomer is less reactive\superfootcite{Chen1995}\superfootcite{Loewe2001}\superfootcite{Iovu2005}\superfootcite{Lanni2010}\superfootcite{Lohwasser2011} than 2-bromo-5-magnesio\-chloro-3-[((S)-2-methyl\-butyl\-oxy)\-methyl]\-thio\-phene. 
This \textit{wrong} terminus combined with the low yield make evident that the presence of an oxy\-gen in side chain has a strong negative influence on this living polymerization.

\cmpd+{ig2-15} and \cmpd+{ig2-8} were polymerized with shorter reaction times (\SI{12}{\hour}) and addition of lithium chloride. From {\HNMR} and {\CNMR} characterization these polymers seems very long and regioregular, terminations signals are so weak to make their identification difficult.
In {\CNMR} spectra of \cmpd+{ig2-15} no aromatic carbon signals can be distinguished from the noise except the ones from regioregular repetition (see Figure~\ref{fig:regioregolarita-nmr-ig2-15-c}). In {\CNMR} spectra of \cmpd+{ig2-8} we can recognize also the presence of signals reported for \cmpd+{ig2-4} in Table~\ref{tab:ig2-4-nmr2d-10} (see Figure~\ref{fig:regioregolarita-nmr-ig2-8-c}). 

A reliable quantitative analysis of terminations using \gls{NMR} spectroscopy, both $^1$\ch{H} and $^{13}$\ch{C}, is prevented by the rigid nature of the conjugated polymer backbone that causes an heterogeneous correlation time $\tau_c$. This fact implies a huge experimental work to find the best relaxation time.

Further termination characterization was performed on sample \cmpd+{ig2-8} with \acrfull{MALDI} and \gls{SEC}\-/\gls{MALDI} by the Institute of Chemistry and Technology of Polymers (\textsmaller{ICTP-CNR}, Catania). 
Linear mode mass spectrometry (spectra in appendix) from two different matrices (di\-thranol and $\alpha$-cyano-4-hydr\-oxy\-cinnamic acid) was performed and results are reported in Table~\ref{tab:maldi}. Surprisingly the terminations were completely different from the \ch{H/Br} termination expected from the literature,\superfootcite{Lohwasser2011} this could be related to a quenching of the polymerization 
due to some oxy\-gen or monomers activated with inverse regioselectivity due to the presence of \ch{LiCl}. 
Another hypothesis is that the quenching wasn't effective because of the not so good solubility of the prepared polymer in the polymerization solvent, tetrahydrofuran at room temperature, that could lead to aggregation prior to quenching agent addition thus reducing the quenching efficiency. 

Preliminary results from \gls{SEC}\-/\gls{MALDI} of compound \cmpd+{ig2-8} are reported in Figure~\ref{fig:ig2-8-sec-maldi}.

\begin{table}[tbp]%maldi
\centering
\caption[MALDI-TOF MS results from \cmpd+{ig2-8} in linear mode.]{\Gls{MALDI} results from polymer \cmpd+{ig2-8} in linear mode from dithranol and $\alpha$-cyano-4-hydr\-oxy\-cinnamic acid matrices.}\label{tab:maldi}
\begin{tabular}{c|c|c|c|c}
\toprule
\multirow{2}{*}{terminations}&\multicolumn{2}{c}{\tallcell{relative intensity in\\ dithranol}}&\multicolumn{2}{c}{\tallcell{relative intensity in\\ $\alpha$-cyano-4-hydr\-oxy-\\cinnamic acid}}\\ 
&low mass\textsuperscript{a}&high mass\textsuperscript{b}&low mass\textsuperscript{a}&high mass\textsuperscript{b}\\\cmidrule{1-5}
\ch{H-/-OH}&30&30&30&35\\
\ch{Br-/-I}&70&25&100&40\\
\ch{Br-/-Ni(dppp)Cl}&25&5&30&5\\
\ch{Br-/-OH}&100&50&85&50\\
\ch{Br-/-H}&20&45&30&40\\
\ch{H-/-I}&20&5&40&10\\
\ch{Br-/-Br}&90&100&75&100\\
\ch{Br-/-NiCl}&30&28&40&20\\
\bottomrule
\end{tabular}

\smallskip

\textsuperscript{a} 2.5 $<$ mass \SI{}{\kg\per\mol}$<$ 3.2; \textsuperscript{b} 3.2 $<$ mass \SI{}{\kg\per\mol}
\end{table}

\begin{SCfigure}[][tbp]%ig2-8-sec-maldi
\centering
\includegraphics[width=0.8\textwidth]{ig2-8-sec-maldi.png}
\caption[Preliminary results from combined SEC\-/MALDI-TOF MS of \cmpd+{ig2-8}.]{Preliminary results from combined \gls{SEC}\-/\gls{MALDI} of \cmpd+{ig2-8}.}
\label{fig:ig2-8-sec-maldi}
\end{SCfigure}

\end{subsection}
\end{section}
\clearpage
\begin{section}[Optical and Morphological Characterization]{Polymer Optical and Morphological Characterization}
\begin{subsection}{Vibrational Spectroscopy Characterization}
\label{sec:ir}

The main absorption bands observed in the infrared spectra of monomer \cmpd+{ig2-10}, polymers samples and their assignments are listed in Table~\ref{tab:ir}. All of the poly\-thio\-phenes have similar infrared spectra but noticeable differences are evident. Spectra are reported in Figure~\ref{fig:ir}.

The infrared band in the region around \SI{3060}{\per\cm} are related to aromatic \ch{C-H}$_\beta$ ($\beta$ 
is referred to proton in 3 or 4 position) stretching.\superfootcite{Hotta1987} \ch{C-H}$_\alpha$ (\ch{H} in position 2 or 5) 
origins a stretching band around \SI{3105}{\per\cm}, this signal is present in \gls{FTIR} spectrum of \cmpd+{ig2-4} as a shoulder thus confirming the presence of \ch{H}$_\alpha$ as spotted by \gls{NMR} spectra discussed in page~\pageref{h-terminale}.

\textsmaller{IR} bands around 1517 and \SI{1461}{\per\cm} were assigned respectively to the \ch{C=C} antisymmetric stretching vibration and \ch{C=C} symmetric stretching.\superfootcite{Scott1969} 

\noindent
The $I_{antisymm}/I_{symm}$ ratio is reported to increase as the polymerization degree increases thus it is a rough indication of the conjugation length.\superfootcite{Furukawa1987} In our polymers the lowest $I_{a}/I_{s}$ ratio is recorded for \cmpd+{ig2-4} in accordance with his shortest length. 

The \textsmaller{IR} band at \SI{1260}{\per\cm} could be related to the conformation of the poly\-thio\-phene chain.\superfootcite{Winokur1989}

\textsmaller{IR} band in the region 830-\SI{840}{\per\cm} is assigned to aromatic \ch{CH} out-of-plane bending of substituted thio\-phene rings.\superfootcite{Hotta1987} The energy of this band is reported to be indicative of the poly\-thio\-phene regioregularity,\superfootcite{Chen1995}\superfootcite{Caronna1997} in our polymers this energy is comparable with similar regioregular poly\-thio\-phenes.\superfootcite{Ohshimizu2008}

The vibration at \SI{471}{\per\cm} is only present in monomer \cmpd+{ig2-10} and seems to involve halogen atoms.\superfootcite{Horak1966}

\begin{table}%ir
\centering
\caption[Characteristic FTIR peaks wavenumbers (\SI{}{\per\cm}) for compounds \cmpd+{ig2-10}, \cmpd+{ig2-4}, \cmpd+{ig2-15} and \cmpd+{ig2-8} films on \ch{KBr} disk.]{Characteristic \gls{FTIR} peaks wavenumbers (\SI{}{\per\cm}) for compounds \cmpd+{ig2-10}, \cmpd+{ig2-4}, \cmpd+{ig2-15} and \cmpd+{ig2-8} films on \ch{KBr} disk.}\label{tab:ir}
\begin{tabular}{c|c|c|c|c}
\toprule
\cmpd+{ig2-10} & \cmpd+{ig2-4} & \cmpd+{ig2-15} & \cmpd+{ig2-8} & Assignment \\ \cmidrule{1-5}
3093	(w)	&	3064	(w)	&	3061	(w)	&	3062	(w)	&	 arom.\ \ch{CH} stretching	\\
2959	(s)	&	2961	(s)	&	2961	(s)	&	2961	(s)	&	 \ch{CH3} degen.\ stretching	\\
2928	(s)	&	2930	(s)	&	2927	(s)	&	2928	(s)	&	 \ch{CH2} antisymm.\ stretching	\\
2873	(s)	&	2874	(s)	&	2874	(s)	&	2874	(s)	&	 \ch{CH3} symm.\ stretching	\\
2858	(s)	&	2858	(s)	&	2858	(s)	&	2858	(s)	&	 \ch{CH2} symm.\ stretching	\\
		&	2805	(w)	&	2805	(w)	&	2805	(w)	&		\\
2733	(w,b)	&	2732	(w,b)	&	2731	(w,b)	&	2731	(w,b)	&		\\
2600	(w,b)	&			&			&			&		\\
		&	2080	(w)	&	2080	(w)	&	2080	(w)	&		\\
1730	(w)	&			&	1736	(m)	&	1730	(w)	&		\\
1675	(w)	&	1670	(m,b)	&	1673	(m)	&	1665	(w)	&		\\
1537	(w)	&	1517	(m)	&	1516	(m)	&	1516	(m)	&	 \tallcell{\ch{C=C} thio\-phene\\ antisymm.\ stretching}	\\
1460	(s)	&	1461	(s)	&	1462	(s)	&	1462	(s)	&	 \tallcell{\ch{C=C} thio\-phene symm.\ \\stretching,\\ \ch{CH3} degen.\ deform.}	\\
1413	(s)	&	1415	(m,b)	&	1414	(m)	&	1414	(m)	&		\\
1377	(m)	&	1378	(s)	&	1378	(s)	&	1378	(s)	&	 \ch{CH3} symm.\ deform.	\\
1260	(w)	&	1260	(w)	&	1261	(m)	&	1260	(w)	&		\\
1178	(s)	&	1180	(m,b)	&	1181	(m)	&	1180	(m)	&		\\
1098	(s)	&	1091	(s)	&	1091	(s)	&	1091	(s)	&	 \ch{CO} antisymm.\ stretching	\\
994	(m)	&			&			&			&	 ring breathing	\\
833	(m)	&	838	(m,b)	&	841	(m)	&	841	(m)	&	 arom.\ \ch{H} out-of-plane bending	\\
		&	801	(w)	&	801	(m)	&	802	(w)	&		\\
471	(m)	&			&			&			&	\tallcell{\ch{C-halogen} sensitive,\\ out-of-plane}	\\

\bottomrule
\end{tabular} 
\end{table}

\newgeometry{left=1in,right=1in,bottom=1in,top=1in}
\begin{landscape}%ir

\thispagestyle{empty}

\centering
\begin{figure}%ir
\centering
\begin{tikzpicture}
\begin{axis}[enlarge x limits=false,enlarge y limits=false,width=1\hsize,height=1\textwidth,x dir=reverse,xminorgrids=true,xmajorgrids=true,ymajorgrids=true,minor x tick num=4,/pgf/number format/1000 sep={},max space between ticks=80pt,xlabel=Wavenumber (cm$^{-1}$),ylabel=Normalized transmittance (\%),legend pos=south west]
\addplot[dashdotted] table {img/spectra/ig2-10-ir.txt};
\addlegendentry{\cmpd+{ig2-10}};
\addplot[densely dotted] table {img/spectra/ig2-4-ir.txt};
\addlegendentry{\cmpd+{ig2-4}};
\addplot[dashed] table[y expr=(\thisrowno{1}*0.94+0.05)] {img/spectra/ig2-15-ir.txt};
\addlegendentry{\cmpd+{ig2-15}};
\addplot[line width=0.2pt] table {img/spectra/ig2-8-ir.txt};
\addlegendentry{\cmpd+{ig2-8}};
\end{axis}
\end{tikzpicture}
\caption[FT-IR of \cmpd+{ig2-10}, \cmpd+{ig2-4}, \cmpd+{ig2-15} and \cmpd+{ig2-8}.]{\gls{FTIR} of pure liquid \cmpd+{ig2-10} and cast films of \cmpd+{ig2-4}, \cmpd+{ig2-15} and \cmpd+{ig2-8} on \ch{KBr} disk. 
}
\label{fig:ir}
\end{figure}
\end{landscape}

\restoregeometry
\end{subsection}
\clearpage

\begin{subsection}{Ultraviolet--Visible Spectroscopy Characterization}
\begin{subsubsection}{In a Good Solvent}

\Acrfull{UVvis} spectra of our polymers in chloroform solution were recorded and reported in Figure~\ref{fig:uvvis}. The main peak belongs to \gls{HOMO}-\gls{LUMO} transition due to $\pi \rightarrow \pi^*$ excitation of conjugated poly\-thio\-phene chromophore. A narrow and weak peak around \SI{265}{\nm} \label{peak-265} could be due to 
non-conjugated rings, common in aromatics.\superfootcite{Lohwasser2012}

\begin{figure}[tbp]%uvvis
\begin{tikzpicture}
\begin{axis}[axis x line=bottom,axis y line=left,enlarge x limits=false,enlarge y limits=false,width=1\textwidth,height=6cm,xlabel=Wavelength (nm),ylabel=Molar absorptivity (\SI{}{\per\Molar\per\cm}),cycle list name=linestyles*,xmin=240,xmax=600]
\addplot table[y expr=\thisrowno{1}/0.0000753] {img/results/ig2-4-uvvis-137-chcl3.txt};
\addlegendentry{\cmpd+{ig2-4} in \ch{CHCl3}};
\addplot table[y expr=\thisrowno{1}/0.0000753] {img/results/ig2-15-uvvis-137-chcl3.txt};
\addlegendentry{\cmpd+{ig2-15} in \ch{CHCl3}};
\addplot table[y expr=\thisrowno{1}/0.0000753] {img/results/ig2-8-uvvis-137-chcl3-for-meoh.txt};
\addlegendentry{\cmpd+{ig2-8} in \ch{CHCl3}};
\end{axis}
\end{tikzpicture}
\caption[UV-vis spectra of \cmpd+{ig2-4}, \cmpd+{ig2-15} and \cmpd+{ig2-8} in good solvents.]{\gls{UVvis} spectra of \cmpd+{ig2-4}, \cmpd+{ig2-15} and \cmpd+{ig2-8} in good solvents. Concentration \SI{137}{\mg\per\liter}. Pathlength \SI{1}{\mm}. Molar absorptivity is referred to the quantity of thio\-phene rings; $\epsilon = \mathrm{abs} / ( l (\mathrm{cm}) \cdot c (\mathrm{mol/L}))$.}
\label{fig:uvvis}
\end{figure}

From the reported \gls{UVvis} spectra in chloroform we can notice a red-shift in the main peaks of \cmpd+{ig2-15} (\SI{447}{\nm}) and \cmpd+{ig2-8} (\SI{445}{\nm}) vs.\ \cmpd+{ig2-4} (\SI{430}{\nm}) reflecting the longer conjugation length of \cmpd+{ig2-15} and \cmpd+{ig2-8}. When the \gls{UVvis} spectrum of \cmpd+{ig2-15} is recorded in tetrahydrofuran only a negligible difference in the maximum absorption wavelength (\SI{449}{\nm}) can be observed. 

Chloroform, ortho-di\-chloro\-benzene and tetra\-chloro\-ethane (\acrshort{TCE}) are good solvents for our polymers, tetrahydrofuran is a borderline solvent as it dissolves our polymers but with a slow dissolution kinetic. Also dichloromethane isn't a perfect solvent for our polymers. Some poor solvents are: \label{poorsolvents} hexane, dimethylformamide, toluene (toluene is a poor solvent at room temperature, but a good solvent at high temperature), methanol, acetonitrile and acetone.

Peak energies are comparable with other regioregular poly\-thio\-phene with side chain starting with methylene-oxy\-gen group.\superfootcite{Higashihara2011}\superfootcite{Zoombelt2008}\superfootcite{VanBeek2005}\superfootcite{McCullough1993}

\label{epsilon}
The molar absorption coefficient referred to the monomeric units for \cmpd+{ig2-4} is \SI{\approx3.8e3}{\per\Molar\per\cm} (\SI{380}{\square\meter\per\mole}), for \cmpd+{ig2-15} is \SI{\approx5.6e3}{\per\Molar\per\cm} (\SI{560}{\square\meter\per\mole}) and for \cmpd+{ig2-8} is \SI{\approx7.0e3}{\per\Molar\per\cm} (\SI{700}{\square\meter\per\mole}).

\end{subsubsection}
\begin{subsubsection}{In Mixtures of Good and Poor Solvents}

\begin{figure}[tbp]%ig2-8-vis-chcl3-meoh
\centering
\includegraphics[width=0.7\textwidth]{ig2-8-vis-chcl3-meoh.jpg}
\caption[Photography of \cmpd+{ig2-8} in chloro\-form-methanol mixtures.]{Photography of \cmpd+{ig2-8} in chloro\-form-methanol mixtures. From left to right (and RGB color descriptions): in \ch{CHCl3} (red 59~\%, green 41~\%), \ch{MeOH} 20~\% (red 61~\%, green 39~\%), \ch{MeOH} 30~\% (red 67~\%, green 33~\%), \ch{MeOH} 40~\% (red 73~\%, green 27~\%), \ch{MeOH} 50~\% (red 78~\%, green 22~\%), \ch{MeOH} 60~\% (red 80~\%, green 20~\%) and \ch{MeOH} 70~\% (red 82~\%, green 18~\%).}
\label{fig:ig2-8-vis-chcl3-meoh}
\end{figure}

Adding some poor solvent to a solution of polymer leads to an aggregation, this can be simply identified by looking at the color of the solutions as shown in Figure~\ref{fig:ig2-8-vis-chcl3-meoh}. 
The obtained aggregates spontaneously precipitate within some hours or days forming a flocculate, this can affect the recording of the spectra.

We noticed that \cmpd+{ig2-4} (fixed concentration \SI{137}{\mg\per\liter}, \SI{\approx0.75}{\milli\Molar} in monomeric units) in chloroform shows a small absorbance shift when a poor solvent like methanol was added (\ch{MeOH} 0~\% \SI{430}{\nm}, \ch{MeOH} 50~\% \SI{433}{\nm}, \ch{MeOH} 60~\% \SI{435}{\nm} and \ch{MeOH} 70~\% \SI{434}{\nm}) while an extensive aggregation and precipitation occurred as revealed by the scattering baseline of sample in 70~\% of methanol (spectra are reported in appendix). 
The absence of a significant peak shift indicates the lack of exciton coupling or conjugation extension in aggregated form compared to solubilized form (maybe the conjugation length is limited by the very small molecular weight, see \gls{SEC} chromatogram in Figure~\ref{fig:ig2-4-ig2-8-sec-chcl3}). At high methanol concentrations a new peak at \SI{330}{\nm} rises, this peak is not assigned. 

As stated before, acetonitrile is a poor solvent for our polymers. Adding acetonitrile to a chloroform solution of polymer \cmpd+{ig2-15} until 1:1 ratio of good-poor solvent resulted in a red-shift of \SI{17}{\nm} (\SI{447}{\nm} $\rightarrow$ \SI{464}{\nm}) due to chain conjugation extension and exciton coupling in the aggregated state. 
Adding acetonitrile to a tetrahydrofuran solution caused a red-shift of \SI{17}{\nm} (\SI{449}{\nm} $\rightarrow$ \SI{466}{\nm}).
When methanol is used as a poor solvent together with tetrahydrofuran a \SI{17}{\nm} red-shift occur adding just 30~\% of non-solvent (\SI{449}{\nm} $\rightarrow$ \SI{466}{\nm}). 
Adding 10~\% more methanol resulted in further \SI{5}{\nm} shift (\SI{466}{\nm} $\rightarrow$ \SI{471}{\nm}). These data suggests that tetrahydrofuran is a worse solvent than chloroform and methanol is worse than acetonitrile.

We notice that when a non-solvent is added the peak at $\approx$~\SI{265}{\nm} becomes weaker, thus supporting the hypothesis that this peak is related to non-conjugated thio\-phene rings (expressed in page~\pageref{peak-265}).

\Gls{UVvis} measurement for \cmpd+{ig2-8} were performed at constant polymer concentration (\SI{137}{\mg\per\liter}, \SI{0.75}{\milli\Molar} in monomeric units, \SI{\approx1e-5}{\Molar} in polymeric chains) both in chloro\-form-methanol and chloro\-form-aceto\-nitrile using a \SI{1}{\mm} cuvette. 
In chloro\-form-methanol until 30~\% \ch{MeOH} no variation in the spectra was observed (\SI{445}{\nm} $\rightarrow$ \SI{444}{\nm}) then at 40~\% \ch{MeOH} a shift of \SI{14}{\nm} can be noticed (\SI{445}{\nm} $\rightarrow$ \SI{459}{\nm}) and the peak at short wavelengths (\SI{265}{\nm}) starts to weaken, see Figure~\ref{fig:ig2-8-uvvis-meoh}. 
Adding more methanol (while keeping constant polymer concentration) the peak energy varies slightly (\SI{459}{\nm} $\rightarrow$ \SI{462}{\nm}) but the absorptivity increases from \SI{8300}{\per\Molar\per\cm} up to \SI{13000}{\per\Molar\per\cm}. Repeating the measurement of the sample with 70~\% \ch{MeOH} after a day resulted in no variation of the absorption spectra (after re\-dispersion of the sample).

The fact that a clear iso\-sbestic point is missing in Figure~\ref{fig:ig2-8-uvvis-meoh} suggests that the transition isn't a simple one-step process. 

\begin{figure}[tbp]%ig2-8-uvvis-meoh
\begin{tikzpicture}
\begin{axis}[axis x line=bottom,axis y line=left,enlarge x limits=false,enlarge y limits=false,width=1\textwidth,height=8cm,xlabel=Wavelength (nm),ylabel=Molar absorptivity (\SI{}{\per\Molar\per\cm}),xmin=240,cycle list name=linestyles*,legend pos=north west,legend style={font=\footnotesize}]
\addplot table[y expr=\thisrowno{1}/0.0000753] {img/results/ig2-8-uvvis-137-chcl3-for-meoh.txt};
\addlegendentry{\cmpd+{ig2-8} in \ch{CHCl3}};
\addplot table[y expr=\thisrowno{1}/0.0000753] {img/results/ig2-8-uvvis-137-chcl3-70-meoh-30.txt};
\addlegendentry{\ch{CHCl3} 70 \% - \ch{MeOH} 30 \%};
\addplot table[y expr=\thisrowno{1}/0.0000753] {img/results/ig2-8-uvvis-137-chcl3-60-meoh-40.txt};
\addlegendentry{\ch{CHCl3} 60 \% - \ch{MeOH} 40 \%};
\addplot table[y expr=\thisrowno{1}/0.0000753] {img/results/ig2-8-uvvis-137-chcl3-50-meoh-50.txt};
\addlegendentry{\ch{CHCl3} 50 \% - \ch{MeOH} 50 \%};
\addplot table[y expr=\thisrowno{1}/0.0000753] {img/results/ig2-8-uvvis-137-chcl3-40-meoh-60.txt};
\addlegendentry{\ch{CHCl3} 40 \% - \ch{MeOH} 60 \%};
\addplot[densely dotted] table[y expr=\thisrowno{1}/0.0000753] {img/results/ig2-8-uvvis-137-chcl3-30-meoh-70.txt};
\addlegendentry{\ch{CHCl3} 30 \% - \ch{MeOH} 70 \%};
\end{axis}
\end{tikzpicture} \caption[UV-vis spectra of \cmpd+{ig2-8} in chloro\-form-methanol mixtures.]{\gls{UVvis} spectra of \cmpd+{ig2-8} in chloro\-form-methanol mixtures. Concentration \SI{137}{\mg\per\liter}. Pathlength \SI{1}{\mm}.}
\label{fig:ig2-8-uvvis-meoh}
\end{figure}

In chloro\-form-aceto\-nitrile \cmpd+{ig2-8} didn't show the absorption intensity enhancement shown in chloro\-form-methanol. Passing from chloroform solution to 40~\% acetonitrile a \SI{23}{\nm} red-shift was observed (\SI{444}{\nm} $\rightarrow$ \SI{467}{\nm}). 
Adding more acetonitrile the peak wavelength didn't change, no iso\-sbestic point can be seen. \label{uvvis-ch3cn} The longer wavelength obtained from aggregation in acetonitrile compared with aggregation in methanol (see Table~\ref{tab:bandgap} on page~\pageref{tab:bandgap}) could indicate a different aggregation geometry, this will be further investigated with circular dichroism spectroscopy in Section~\ref{sec:cd-poor}. 

The largest red-shift we observed with our polymers (\SI{23}{\nm}) is still smaller than some of the values reported in the literature for aggregated regioregular poly\-thio\-phenes.\superfootcite{McCullough1993}\superfootcite{Goto2002}\superfootcite{Kiriy2003}\superfootcite{Shi2006} This could point to a different organization with a lesser degree of planarization or $\pi-\pi$ stacking.

\end{subsubsection}
\begin{subsubsection}{In Solid State}

\begin{table}%bandgap
\centering
\caption[UV-vis data and estimated band gaps for samples of \cmpd+{ig2-4}, \cmpd+{ig2-15} and \cmpd+{ig2-8}.]{\gls{UVvis} data and estimated band gaps for polymers \cmpd+{ig2-4}, \cmpd+{ig2-15} and \cmpd+{ig2-8}.}\label{tab:bandgap}
\begin{tabular}{c|c|c|c|c}
\toprule
polymer&processing&\tallcell{\gls{UVvis} peak\\(\SI{}{\nm})}&\tallcell{\gls{UVvis} onset\\ (\SI{}{\nm})}&\tallcell{Band gap\\(\SI{}{\eV})} \\ \cmidrule{1-5}
\multirow{3}{*}{\cmpd+{ig2-4}}&\ch{CHCl3} solution	&430	&523	&2.37	 		 	\\ 
&spin \SI{700}{\rpm} from \ch{CHCl3}	&$\approx462$	&574	&2.16	 	\\
&in \ch{KCl}&$\approx464$	&&	\\ \cmidrule{1-1}
\multirow{4}{*}{\cmpd+{ig2-15}}&\ch{CHCl3} solution	&446	&531	&2.33	 		 	\\ 
&\gls{THF} solution		&450&533	&2.33	 \\
&cast from \ch{CH2Cl2}&$\approx486$	&589	&2.10		\\ 
&in \ch{KCl}&$\approx493$	&&	\\\cmidrule{1-1}
\multirow{8}{*}{\cmpd+{ig2-8}}&\ch{CHCl3} solution	&444	&528	&2.35	 		 	\\
&cast from \ch{CHCl3}	&$\approx471$	&584&2.12\\
&cast from \gls{THF}		&$\approx473$	&581&2.13 \\
&cast from \ch{CH2Cl2}	&$\approx469$	&577&2.15	\\
&spin \SI{700}{\rpm} from \ch{CHCl3}	&$\approx476$	&571&2.17	\\
&spin \SI{700}{\rpm} from \gls{THF}	&$\approx476$	&566&2.19	\\
&spin \SI{1000}{\rpm} from \ch{CHCl3}&$\approx492$	&587&2.11	\\
&in \ch{KCl}&$\approx495$	&&	\\

\bottomrule
\end{tabular}
\end{table}

\label{ig2-4-uvvis}
\cmpd+{ig2-4} was analyzed in a thin film spin coated (\SI{700}{\rpm}) from chloroform solution. While only a small red-shift was observed passing from chloroform solution to chloro\-form-methanol mixtures (\SI{430}{\nm} $\rightarrow$ \SI{434}{\nm}), in the \gls{UVvis} spectrum of this thin film (reported in appendix), a \SI{32}{\nm} red-shift was observed (\SI{430}{\nm} $\rightarrow$ \SI{462}{\nm}). This fact suggests that the aggregation obtained in 70~\% \ch{MeOH} doesn't imply a good packing of chains while this takes place in solid state.

\begin{figure}[tbp]%ig2-15-uvvis-film
\begin{tikzpicture}
\begin{axis}[axis x line=bottom,axis y line=left,enlarge x limits=false,enlarge y limits=false,yticklabels={,,},width=1\textwidth,height=6cm,xlabel=Wavelength (nm),ylabel=Absorbance,xmin=240,xmax=650,legend pos=north west,legend style={font=\footnotesize},cycle list name=linestyles*]
\addplot table[y expr=\thisrowno{1}/1.04169] {img/results/ig2-15-uvvis-chcl3.txt};
\addlegendentry{\cmpd+{ig2-15} in \ch{CHCl3}};
\addplot table[y expr=\thisrowno{1}/1.8789] {img/results/ig2-15-uvvis-thf-55-ch3cn-45.txt};
\addlegendentry{\gls{THF} 55 \% - \ch{CH3CN} 45 \%};
\addplot table[y expr=(\thisrowno{1}-0.15)/0.219712] {img/results/ig2-15-uvvis-cast-ch2cl2.txt};
\addlegendentry{cast from \ch{CH2Cl2}};
\end{axis}
\end{tikzpicture} \caption[UV-vis spectra of polymer \cmpd+{ig2-15} in solution and thin film.]{Normalized UV-vis spectra of polymer \cmpd+{ig2-15} in chloroform, in tetra\-hydro\-furan-aceto\-nitrile solution and thin film cast from a dichloromethane solution.}
\label{fig:ig2-15-uvvis-film}
\end{figure}

\label{ig2-15-uvvis-film}
In Figure~\ref{fig:ig2-15-uvvis-film} a strong red-shift is observed for \cmpd+{ig2-15} cast from dichloromethane (chloroform solution peak \SI{447}{\nm}, \gls{THF} 55~\% - \ch{MeOH} 45~\% \SI{466}{\nm}, cast from dichloromethane \SI{\approx486}{\nm}).

\label{ig2-8-film}
For \cmpd+{ig2-8} many films were analyzed, both from spin coating and from casting. 
The obtained films vary in thickness and the wavelength of the maximum absorption peak in their spectra also vary. 
Thin films from casting have a maximum absorption wavelength around \SI{\approx471}{\nm} (cast from \ch{CHCl3} \SI{\approx471}{\nm}, cast from \gls{THF} \SI{\approx473}{\nm}, cast from \ch{CH2Cl2} \SI{\approx469}{\nm}) while spin coated thin films are generally slightly red-shifted to \SI{\approx476}{\nm} (spin \SI{700}{\rpm} from \ch{CHCl3} \SI{\approx476}{\nm}, spin \SI{700}{\rpm} from \gls{THF} \SI{\approx476}{\nm}) with only a case of film resulting in a very strong red-shift to \SI{\approx492}{\nm} (spin \SI{1000}{\rpm} from \ch{CHCl3}). These \gls{UVvis} spectra are reported in appendix. 
These values are somewhat different from peak wavelengths of colloidal dispersions in solvent-non-solvent mixtures (\ch{CHCl3} - \ch{MeOH} \SI{\approx462}{\nm}, \ch{CHCl3} - \ch{CH3CN} \SI{\approx468}{\nm}, \gls{THF} - \ch{CH3CN} \SI{467}{\nm}, see Table~\ref{tab:bandgap}). The out\-lier film (spin \SI{1000}{\rpm} from \ch{CHCl3}) also showed a weaker absorption at short wavelengths (\SI{265}{\nm}). The fact that all spin coated films are red-shifted relatively to all cast films 
means that a fast spin coating forces the chains to stretch and flatten on the quartz substrate while a slow deposition allows the chains to aggregate in a preferred organization with shorter conjugation length or a worse packing. 

Thin films from similar polythio\-phenes bearing unbranched side chains\superfootcite{Zoombelt2008} showed a lower $\pi-\pi^*$ transition energy (\SI{500}{\nm}). This could be related to a worse exciton coupling in our chiral polythio\-phenes or to a worse packing caused by the branching of our side chains.

\label{uvvis-pentane}
Dipping our films in a side groups-selective solvent could give a reorganization of side chains helping the reaching of a more ordered arrangement of the macromolecules.\superfootcite{Caronna1997}\superfootcite{Salatelli2010} Indeed no change in main peak absorption was observed when spin coated films were dipped in pentane; only the \SI{265}{\nm} weak peak, where present, decreased in intensity (that could be related to ordering of non-conjugated rings, see page~\pageref{peak-265}). 

The fact that we could not observe a vibronic fine structure in any of the spectra of our samples, suggests a lack of crystallinity,\superfootcite{Brown2003}\superfootcite{Zoombelt2008} as will be further 
investigated in Section~\ref{xrd}.

The band gap in solution can be estimated from onset of \gls{HOMO}-\gls{LUMO} transition (a $\pi \rightarrow \pi^*$ transition) in \gls{UVvis} spectra. The utilized onset point was the point where line tangent to the peak's right inflection point intersect the baseline. The results are reported in Table~\ref{tab:bandgap}. 

Polymers were analyzed also in \ch{KCl} disks (dried from chloroform with a nitrogen flux, ground with \ch{KCl} and pressed). The obtained disks weren't homogeneous thus a strong scattering was the issue with these measurements. %was noticed in Figure~\ref{fig:uvvis-kcl}.
In disks from \cmpd+{ig2-4} no wavelength change was observed comparing to the spin coated thin film. 
Both \cmpd+{ig2-15} and \cmpd+{ig2-8} have a main peak wavelength around \SI{495}{\nm} and we can notice (see Table~\ref{tab:bandgap}) a red-shift referring to \SI{700}{\rpm} spin coated and cast films. The facts that \cmpd+{ig2-8} peak position in \ch{KCl} isn't significantly different from \SI{1000}{\rpm} spin coated film and that peak of \cmpd+{ig2-4} in \ch{KCl} isn't different from \cmpd+{ig2-4} spin coated prompt us to believe that drying the samples with a nitrogen flux had the same effect of a fast spin coating on polymers. But we can't exclude that an aggregation change could be due to sintering of the polymer due to the huge pressure used to produce \ch{KCl} disks.\superfootcite{Loi2001} 

\end{subsubsection}
\end{subsection}
\clearpage
\begin{subsection}{Photoluminescence Spectroscopy Characterization}
\label{sec:pl}

\begin{subsubsection}{In a Good Solvent}

\begin{figure}[tbp]%pl-vs-uvvis
\begin{tikzpicture}
\begin{axis}[axis x line=bottom,axis y line=left,enlarge x limits=false,enlarge y limits=false,width=1\textwidth,height=8cm,xlabel=Wavelength (nm),ylabel=Absorption (a.\ u.),xmin=300,xmax=800,yticklabels={,,},legend entries={\cmpd{ig2-4} in \ch{CHCl3},\cmpd{ig2-15} in \ch{CHCl3},\cmpd{ig2-8} in \ch{CHCl3}}]
\addplot[solid] table[y expr=\thisrowno{1}*100] {img/results/ig2-4-uvvis-chcl3.txt};
\addplot[dashed] table[y expr=\thisrowno{1}*127] {img/results/ig2-15-uvvis-chcl3.txt};
\addplot[dotted] table[y expr=\thisrowno{1}*215] {img/results/ig2-8-uvvis-137-chcl3-for-meoh.txt};
\node at (axis cs:600,15) {emission};
\node at (axis cs:430,15) {absorption};
\end{axis}
\begin{axis}[axis x line=none,axis y line=right,enlarge x limits=false,enlarge y limits=false,width=1\textwidth,height=8cm,ylabel=Photoluminescence (a.\ u.),xmin=300,xmax=800]
\addplot[solid] table {img/results/ig2-4-pl-5480-chcl3.txt};
\addplot[dashed] table {img/results/ig2-15-pl-5480-chcl3.txt};
\addplot[dotted] table {img/results/ig2-8-pl-5480-chcl3.txt};
\end{axis}
\end{tikzpicture} \caption[Photoluminescence and absorption of \cmpd+{ig2-4}, \cmpd+{ig2-15} and \cmpd+{ig2-8} chloroform solutions.]{Photoluminescence and absorption of \cmpd+{ig2-4}, \cmpd+{ig2-15} and \cmpd+{ig2-8} chloroform solutions. Concentration for photoluminescence samples is \SI{5.48}{\g\per\liter}. Absorption spectra are rescaled to \gls{PL} intensity.}
\label{fig:pl-vs-uvvis}
\end{figure}

In a good solvent like chloroform, \cmpd+{ig2-15} and \cmpd+{ig2-8} have similar photoluminescence spectra (see Figure~\ref{fig:pl-vs-uvvis}) when stimulated at \SI{450}{\nm}. This was expected because in solution the chirality of side chains shouldn't influence to a big extent the conjugated backbone which is responsible for absorption and emission. 
\cmpd+{ig2-4} in chloroform shows a lower fluorescence efficiency (see Figure~\ref{fig:pl-vs-uvvis}).
For each polymer increasing by 40 times the concentration had no effect on fluorescence spectra except for a small decrease of emission maybe due to interaction between chains, not yet an aggregation, or due to auto-absorption (as shown in Figure~\ref{fig:pl-vs-uvvis}, the superposition of absorption and emission spectra is tiny, thus an auto-absorption by F\"{o}rster resonance energy transfer should have a limited importance). 
All the polymers show fluorescence maximum at about the same wavelength (\SI{575}{\nm}) with an intense shoulder at lower energies (\SI{600}{\nm}). Usually in poly\-thio\-phenes this shoulder is assigned to a small extent of aggregation present even in good solvent or to a vibronic band. No similar vibronic shoulder was seen in absorption spectra. Another possibility is that the two peaks could originate from different molecular structures, for example from two different molecular weights (see Section~\ref{sec:weight} for results on weight), with the emission from the low weight fraction pumping energy in the absorbance of the high weight fraction via a F\"{o}rster resonance energy transfer. Moreover we can hypothesize the presence of different chain conformations in solution. 

The Stokes shift\nota{\textit{Stokes shift} is the difference (in wavelength or frequency units) between positions of the band maxima of the absorption and emission spectra (fluorescence and Raman being two examples) of the same electronic transition. (Wikipedia)} 
is about \SI{0.63}{\eV} (absorption \SI{\approx445}{\nm}, \SI{2.79}{\eV}, emission \SI{\approx575}{\nm}, \SI{2.16}{\eV}), comparable with other poly\-thio\-phenes.\superfootcite{Langeveld-Voss1996}\superfootcite{Xu1993}\superfootcite{Theander1999}\superfootcite{Ohshimizu2008}\superfootcite{Higashihara2011}\superfootcite{Bouman1994a}

\end{subsubsection}
\begin{subsubsection}{In Mixtures of Good and Poor Solvents}

\begin{figure}[tbp]%ig2-8-pl365-chcl3-meoh
\centering
\includegraphics[width=0.7\textwidth]{ig2-8-pl365-chcl3-meoh.jpg}
\caption[Photography of photoluminescence of \cmpd+{ig2-8} in chloro\-form-methanol mixtures.]{Photography of photoluminescence of \cmpd+{ig2-8} in chloro\-form-methanol mixtures illuminated at \SI{365}{\nm}. From the left to the right: in \ch{CHCl3}, \ch{MeOH} 20~\%, \ch{MeOH} 30~\%, \ch{MeOH} 40~\%, \ch{MeOH} 50~\%, \ch{MeOH} 60~\% and \ch{MeOH} 70~\%.}
\label{fig:ig2-8-pl365-chcl3-meoh}
\end{figure}

\begin{figure}[tbp]%ig2-15-pl
\begin{tikzpicture}
\begin{axis}[axis x line=bottom,axis y line=left,enlarge x limits=false,enlarge y limits=false,width=1\textwidth,height=6cm,xlabel=Wavelength (nm),ylabel=Photoluminescence (a.\ u.),xmin=480,cycle list name=linestyles*]
\addplot table {img/results/ig2-15-pl-137-chcl3.txt};
\addlegendentry{\cmpd+{ig2-15} in \ch{CHCl3}};
\addplot table {img/results/ig2-15-pl-137-chcl3-75-meoh-25.txt};
\addlegendentry{in \ch{CHCl3} 75 \% - \ch{MeOH} 25 \%};
\addplot table {img/results/ig2-15-pl-137-chcl3-50-meoh-50.txt};
\addlegendentry{in \ch{CHCl3} 50 \% - \ch{MeOH} 50 \%};
\addplot table {img/results/ig2-15-pl-137-chcl3-30-meoh-70.txt};
\addlegendentry{in \ch{CHCl3} 30 \% - \ch{MeOH} 70 \%};
\end{axis}
\end{tikzpicture} \caption[Photoluminescence of \cmpd+{ig2-4} in chloro\-form-methanol mixtures.]{Photoluminescence of \cmpd+{ig2-4} in chloro\-form-methanol mixtures. Concentration \SI{137}{\mg\per\liter}.}
\label{fig:ig2-15-pl}
\end{figure}

\begin{figure}[tbp]%ig2-8-pl
\begin{tikzpicture}
\begin{axis}[axis x line=bottom,axis y line=left,enlarge x limits=false,enlarge y limits=false,width=1\textwidth,height=6cm,xlabel=Wavelength (nm),ylabel=Photoluminescence (a.\ u.),xmin=480,cycle list name=linestyles*]
\addplot table {img/results/ig2-8-pl-137-chcl3.txt};
\addlegendentry{\cmpd+{ig2-8} in \ch{CHCl3}};
\addplot table {img/results/ig2-8-pl-137-chcl3-75-meoh-25.txt};
\addlegendentry{in \ch{CHCl3} 75 \% - \ch{MeOH} 25 \%};
\addplot table {img/results/ig2-8-pl-137-chcl3-50-meoh-50.txt};
\addlegendentry{in \ch{CHCl3} 50 \% - \ch{MeOH} 50 \%};
\addplot table {img/results/ig2-8-pl-137-chcl3-30-meoh-70.txt};
\addlegendentry{in \ch{CHCl3} 30 \% - \ch{MeOH} 70 \%};
\end{axis}
\end{tikzpicture} \caption[Photoluminescence of \cmpd+{ig2-8} in chloro\-form-methanol mixtures]{Photoluminescence of \cmpd+{ig2-8} in chloro\-form-methanol mixtures. Concentration \SI{137}{\mg\per\liter}.}
\label{fig:ig2-8-pl}
\end{figure}

\begin{figure}[tbp]%ig2-8-ig2-15-pl
\begin{tikzpicture}
\begin{axis}[axis x line=bottom,axis y line=left,enlarge x limits=false,enlarge y limits=false,width=1\textwidth,height=6cm,xlabel=Wavelength (nm),ylabel=Photoluminescence (a.\ u.),xmin=480,legend style={font=\footnotesize},cycle list name=linestyles*]
\addplot table {img/results/ig2-15-pl-137-chcl3-30-meoh-70.txt};
\addlegendentry{\cmpd+{ig2-15} in \ch{CHCl3} 30 \% - \ch{MeOH} 70 \%};
\addplot table {img/results/ig2-8-pl-137-chcl3-50-meoh-50.txt};
\addlegendentry{\cmpd+{ig2-8} in \ch{CHCl3} 50 \% - \ch{MeOH} 50 \%};
\addplot table {img/results/ig2-8-pl-137-chcl3-30-meoh-70.txt};
\addlegendentry{\cmpd+{ig2-8} in \ch{CHCl3} 30 \% - \ch{MeOH} 70 \%};
\end{axis}
\end{tikzpicture} \caption[Photoluminescence of \cmpd+{ig2-15} and \cmpd+{ig2-8} in chloro\-form-methanol mixtures.]{Comparison of photoluminescence of \cmpd+{ig2-15} and \cmpd+{ig2-8} in chloro\-form-methanol mixtures. Spectra are not rescaled, concentration is \SI{137}{\mg\per\liter} for both \cmpd+{ig2-15} and \cmpd+{ig2-8}.}
\label{fig:ig2-8-ig2-15-pl}
\end{figure}

As it's clear by looking at Figure~\ref{fig:ig2-8-pl365-chcl3-meoh}, the addition of methanol to polymer solutions in chloroform produces a quenching in the fluorescence due to the formation of aggregated states with very low quantum efficiency.

Adding 50~\% of methanol to a chloroform solution of polymer \cmpd+{ig2-4} induces only a small change in photoluminescence spectra, just a reduction of the shoulder at longer wavelength.

Adding 25~\%, 50~\%, 70~\% of methanol to a chloroform solution of \cmpd+{ig2-15} resulted in a progressive quenching of photoluminescence, both for the main emission peak and of the shoulder (see Figure~\ref{fig:ig2-15-pl}). 
Measuring again these samples after an hour resulted in a diminished fluorescence (approximately of 10 - 25~\%) maybe due the greater degree of ordering reached by the aggregates. The absence of a strong quenching of photoluminescence indicates the presence of non well-packed chains.

In every spectra the position of the peaks is unvaried compared with photoluminescence of polymer in chloroform, but the relative intensity of the main peak and the shoulder varies. Adding methanol both peak show a decreased intensity but the quenching of the higher energy peak is more important than for the shoulder. In 70~\% of methanol the two peaks from \cmpd+{ig2-15} have the same fluorescence intensity. 

\cmpd+{ig2-8} in a mixture of chloroform with 25~\% methanol gives a fluorescence spectra identical to \cmpd+{ig2-15} in 25~\% methanol. The further addition of methanol to solutions of \cmpd+{ig2-8} (see Figure~\ref{fig:ig2-8-pl}) has a rather different effect: a dramatic quenching of fluorescence and a strong red-shift of the emission. 
This behavior can be noted also from Figure~\ref{fig:ig2-8-pl365-chcl3-meoh} where \cmpd+{ig2-8} is illuminated with a common \SI{365}{\nm} lamp. In Figure~\ref{fig:ig2-8-ig2-15-pl} is reported a comparison between \gls{PL} of \cmpd+{ig2-15} and \cmpd+{ig2-8} at high methanol percents. The new peak at long wavelengths could be assigned to aggregated species, it was not distinguishable in the previous spectra because of the more intense emission of non aggregated chains. 
The non enantio\-purity of \cmpd+{ig2-15} can be regarded as a degree of disorder that can interfere with the ordering needed for the formation of aggregates.

\end{subsubsection}
\begin{subsubsection}{In Solid State}

Up to now no measurement of photoluminescence on thin film was performed because of their very weak stimulated emission.

\end{subsubsection}
\end{subsection}
\clearpage
\begin{subsection}{Circular Dichroism Spectroscopy Characterization}
\begin{subsubsection}{In a Good Solvent}

Both \cmpd+{ig2-15} in tetrahydrofuran and \cmpd+{ig2-8} in chloroform showed a non zero circular dichroism signal. The signal is extremely weak but we can recognize a bisignate peak 
with the inflection point located around the \gls{UVvis} main absorption peak. These signals may recall a 
positive bisignate circular dichroism couplet (the sign of a bisignate couplet is defined as the sign of the long wavelength part of this couplet; the bisignate peak inflection point and amplitude were obtained fitting the peak with a sine waveform, see page~\pageref{fitting}). 

\end{subsubsection}
\begin{subsubsection}{In Mixtures of Good and Poor Solvents}
\label{sec:cd-poor}

\begin{table}%soluzioni-cd
\centering
\caption[UV-vis and CD data of \cmpd+{ig2-15} and \cmpd+{ig2-8} solutions.]{\gls{UVvis} and \gls{CD} data of \cmpd+{ig2-15} and \cmpd+{ig2-8} solutions.}\label{tab:soluzioni-cd}
\begin{tabular}{c|c|c|c|c}
\toprule
processing&\tallcell{UV-vis\\ peak\\(\SI{}{nm})}&\tallcell{CD\\ inflection\\(\SI{}{nm})}&\tallcell{CD\\ amplitude\\(\SI{}{mdeg})}&\tallcell{molar\\ circular\\ dichroism\\ ($\Delta\epsilon$, \\ \SI{}{\per\Molar\per\cm})}\\ \cmidrule{1-5}
\cmpd+{ig2-15}					&\multicolumn{4}{c}{}\\ \cmidrule{1-1}
\ch{CHCl3} solution						&446&	&0	&0\\
\gls{THF} solution							&450&	&	 &	\\
\ch{CHCl3} 70\% \ch{MeOH} 30\%	&446&$\approx460$&0.5&0.2	\\
\ch{CHCl3} 70\% \ch{MeOH} 30\% \SI{30}{\minute}&$\approx446$&457&1	 &0.4	\\
\ch{CHCl3} 50\% \ch{MeOH} 50\%	&$\approx466$	&459&7.7	 &3.1	\\
\ch{CHCl3} 40\% \ch{CH3CN} 60\%&$\approx469$&454&2.3	 &0.9	\\ 
\ch{CHCl3} 30\% \ch{CH3CN} 70\%&$\approx469$&460&4.8	 &1.9	\\
\ch{CHCl3} 20\% \ch{CH3CN} 80\%&$\approx469$&460&4.8	 &1.9	\\
\ch{CHCl3} 10\% \ch{CH3CN} 90\%&$\approx470$	&460&4.1	 &1.7	\\
\gls{THF} 60\% \ch{CH3CN} 40\%	&$\approx466$&458&5.2	 &2.1	\\
\gls{THF} 50\% \ch{CH3CN} 50\%	&469&458&4	 &1.6	\\
\gls{THF} 50\% \ch{CH3CN} 50\% \SI{274}{\mg/\liter}&466&458&11.4	 &2.3	\\
\gls{THF} 50\% \ch{CH3CN} 50\% \SI{548}{\mg/\liter}&469&459&26.2	 &2.6	\\
\gls{THF} 40\% \ch{CH3CN} 60\%	&$\approx468$&456&2.8	&1.1	\\
\gls{THF} 20\% \ch{CH3CN} 80\%	&$\approx470$&459&3	 &1.2\\ \cmidrule{1-5}
\cmpd+{ig2-8}					&\multicolumn{4}{c}{}\\ \cmidrule{1-1}
\ch{CHCl3} solution 						&444	&	&0	&0	\\
\ch{CHCl3} 70\% \ch{MeOH} 30\%	&$\approx443$&&0&0\\
\ch{CHCl3} 70\% \ch{MeOH} 30\% \SI{10}{\minute}&$\approx447$&441&6.1&2.5\\
\ch{CHCl3} 70\% \ch{MeOH} 30\% \SI{30}{\minute}&443&445&3.5&1.4\\
\ch{CHCl3} 70\% \ch{MeOH} 30\% \SI{1}{\day}&$\approx443$&443&22.9&9.4\\
\ch{CHCl3} 60\% \ch{MeOH} 40\%	&459&456&13.8&5.5\\
\ch{CHCl3} 50\% \ch{MeOH} 50\%	&461&457&15.6&6.4\\
\ch{CHCl3} 40\% \ch{MeOH} 60\%	&462&456&19.6&7.9\\
\ch{CHCl3} 30\% \ch{MeOH} 70\%	&462&456&24&9.7\\
\ch{CHCl3} 40\% \ch{CH3CN} 60\%	&467&458&7.3&2.9\\
\ch{CHCl3} 30\% \ch{CH3CN} 70\%	&465&458&12.4&5.0\\
\ch{CHCl3} 20\% \ch{CH3CN} 80\%	&467&459&10.1&4.1\\
\ch{CHCl3} 10\% \ch{CH3CN} 90\%	&469&459&6.5&2.6\\
\gls{THF} 60\% \ch{CH3CN} 40\%	&$\approx466$&457&6.4&2.6\\
\gls{THF} 50\% \ch{CH3CN} 50\%	&$\approx468$&457&5.5&2.2\\
\gls{THF} 40\% \ch{CH3CN} 60\%	&$\approx467$&457&7.0&2.8\\
\gls{THF} 20\% \ch{CH3CN} 80 \%	&$\approx467$&458&8.2&3.3\\
\bottomrule
\end{tabular}

\smallskip

{Where not specified a \SI{137}{\mg\per\liter} concentration was used.}
\end{table}

In chloro\-form-methanol mixtures, \cmpd+{ig2-4} showed at first (100~\% chloroform) a dichroic monosignate negative peak in the region of the main \gls{UVvis} absorption, then increasing methanol content, a bisignate peak appeared.

This positive bisignate circular dichroism couplet, that is what we detect for all our polymers in the aggregated form, could rise from intramolecular or intermolecular interactions. Referring to the latter case a positive couplet is reported as being produced by two close identical chromophores with a right handed relative orientation,\superfootcite{Berova2007} this is confirmed by experiments on oligo\-thio\-phene with a known relative orientation.\superfootcite{Langeveld-Voss1998} 

\begin{figure}[tbp]%ig2-15-cd-meoh
\begin{tikzpicture}
\begin{axis}[axis x line=bottom,axis y line=left,enlarge x limits=false,enlarge y limits=true,width=1\textwidth,height=6cm,xlabel=Wavelength (nm),ylabel=Molar circular dichroism $\Delta\epsilon$,xmin=240,legend pos=north west,cycle list name=linestyles*]
\addplot table[y expr=\thisrowno{1}*0.403] {img/results/ig2-15-cd-137-chcl3-70-meoh-30.txt};
\addlegendentry{in \ch{CHCl3} 70 \% - \ch{MeOH} 30 \%};
\addplot table[y expr=\thisrowno{1}*0.403] {img/results/ig2-15-cd-137-chcl3-70-meoh-30-30min.txt};
\addlegendentry{in \ch{CHCl3} 70 \% - \ch{MeOH} 30 \% \SI{30}{\minute}};
\addplot table[y expr=\thisrowno{1}*0.403] {img/results/ig2-15-cd-137-chcl3-50-meoh-50.txt};
\addlegendentry{in \ch{CHCl3} 50 \% - \ch{MeOH} 50 \%};
\end{axis}
\end{tikzpicture} \caption[Circular dichroism of \cmpd+{ig2-15} in chloro\-form-methanol mixtures.]{Circular dichroism of \cmpd+{ig2-15} in chloro\-form-methanol mixtures. Concentration \SI{137}{\mg\per\liter}, pathlength \SI{1}{\mm}. Molar circular dichroism is referred to the quantity of thio\-phene rings; $\Delta\epsilon = \theta (\mathrm{mdeg}) / (32982 \cdot l (\mathrm{cm}) \cdot c (\mathrm{mol/L}))$.}
\label{fig:ig2-15-cd-meoh}
\end{figure}

In contrast with measurement on \cmpd+{ig2-4}, \cmpd+{ig2-15} doesn't show any monosignate circular dichroism (CD) peak. As shown in Figure~\ref{fig:ig2-15-cd-meoh}, adding 30~\% of methanol to a chloroform solution generates only an extremely weak bisignate signal at first (amplitude \SI{0.5}{mdeg}), but it was sufficient to wait half an hour to see a clear dichroic signal rising. Adding more methanol resulted in strong enhancement of circular dichroism. 
Looking at the bisignate couplet inflection points and \gls{UVvis} peak wavelengths 
we notice no variation in inflection points (\SI{457}{\nm} for \ch{MeOH} 30~\% \SI{30}{\minute}, \SI{459}{\nm} for \ch{MeOH} 50~\%) while a \SI{20}{\nm} red shift is observed in \gls{UVvis} spectra. 
All the significant numeric data for \gls{UVvis} and \gls{CD} characterization in solution are reported in Table~\ref{tab:soluzioni-cd}.

\begin{figure}[tbp]%ig2-15-cd-ch3cn
\begin{tikzpicture}
\begin{axis}[axis x line=bottom,axis y line=left,enlarge x limits=false,enlarge y limits=true,width=1\textwidth,height=8cm,xlabel=Wavelength (nm),ylabel=Molar circular dichroism $\Delta\epsilon$,xmin=240,legend pos=north west,cycle list name=linestyles*]
\addplot table[y expr=\thisrowno{1}*0.403] {img/results/ig2-15-cd-137-chcl3-40-ch3cn-60.txt};
\addlegendentry{in \ch{CHCl3} 40 \% - \ch{CH3CN} 60 \%};
\addplot table[y expr=\thisrowno{1}*0.403] {img/results/ig2-15-cd-137-chcl3-30-ch3cn-70.txt};
\addlegendentry{in \ch{CHCl3} 30 \% - \ch{CH3CN} 70 \%};
\addplot table[y expr=\thisrowno{1}*0.403] {img/results/ig2-15-cd-137-chcl3-20-ch3cn-80.txt};
\addlegendentry{in \ch{CHCl3} 20 \% - \ch{CH3CN} 80 \%};
\addplot table[y expr=\thisrowno{1}*0.403] {img/results/ig2-15-cd-137-chcl3-10-ch3cn-90.txt};
\addlegendentry{in \ch{CHCl3} 10 \% - \ch{CH3CN} 90 \%};
\end{axis}
\end{tikzpicture} \caption[Circular dichroism of \cmpd+{ig2-15} in chloro\-form-aceto\-nitrile mixtures.]{Circular dichroism of \cmpd+{ig2-15} in chloro\-form-aceto\-nitrile mixtures. Concentration \SI{137}{\mg\per\liter}, pathlength \SI{1}{\mm}.}
\label{fig:ig2-15-cd-ch3cn}
\end{figure}

In chloro\-form-aceto\-nitrile mixtures passing from 60~\% to 70~\% of non solvent a clear shift of the bisignate couplet from \cmpd+{ig2-15} was noticed (\SI{454}{\nm} $\rightarrow$ \SI{460}{\nm}) joined with an increase of \gls{CD} intensity. Adding more non solvent the signal remains quite the same (see Figure~\ref{fig:ig2-15-cd-ch3cn}). During this addition the \gls{UVvis} spectra remains almost unchanged (see Table~\ref{tab:soluzioni-cd}). Interesting is the presence of a monosignate peak in correspondence with the \gls{UVvis} signal at \SI{270}{\nm}, likely related to non conjugated rings absorption (see page~\pageref{peak-265}).

\begin{figure}[tbp]%ig2-15-cd-thf-ch3cn
\begin{tikzpicture}
\begin{axis}[axis x line=bottom,axis y line=left,enlarge x limits=false,enlarge y limits=true,width=1\textwidth,height=8cm,xlabel=Wavelength (nm),ylabel=Molar circular dichroism $\Delta\epsilon$,xmin=240,legend style={font=\footnotesize},legend pos=north west,cycle list name=linestyles*]
\addplot table[y expr=\thisrowno{1}*0.403] {img/results/ig2-15-cd-137-thf-60-ch3cn-40.txt};
\addlegendentry{in \gls{THF} 60 \% - \ch{CH3CN} 40 \%};
\addplot table[y expr=\thisrowno{1}*0.403] {img/results/ig2-15-cd-137-thf-50-ch3cn-50.txt};
\addlegendentry{in \gls{THF} 50 \% - \ch{CH3CN} 50 \%};
\addplot table[y expr=\thisrowno{1}*0.201] {img/results/ig2-15-cd-274-thf-50-ch3cn-50.txt};
\addlegendentry{in \gls{THF} 50 \% - \ch{CH3CN} 50 \% \SI{274}{\mg\per\liter}};
\addplot table[y expr=\thisrowno{1}*0.101] {img/results/ig2-15-cd-548-thf-50-ch3cn-50.txt};
\addlegendentry{in \gls{THF} 50 \% - \ch{CH3CN} 50 \% \SI{548}{\mg\per\liter}};
\addplot table[y expr=\thisrowno{1}*0.403] {img/results/ig2-15-cd-137-thf-40-ch3cn-60.txt};
\addlegendentry{in \gls{THF} 40 \% - \ch{CH3CN} 60 \%};
\addplot[densely dotted] table[y expr=\thisrowno{1}*0.403] {img/results/ig2-15-cd-137-thf-20-ch3cn-80.txt};
\addlegendentry{in \gls{THF} 20 \% - \ch{CH3CN} 80 \%};
\end{axis}
\end{tikzpicture} \caption[Circular dichroism of \cmpd+{ig2-15} in tetra\-hydro\-furan-aceto\-nitrile mixtures.]{Circular dichroism of \cmpd+{ig2-15} in tetra\-hydro\-furan-aceto\-nitrile mixtures. Where not specified a concentration of \SI{137}{\mg\per\liter} was used. Pathlength \SI{1}{\mm}.}
\label{fig:ig2-15-cd-thf-ch3cn}
\end{figure}

\label{nonlinearita}
In order to discern if the interactions responsible of the \gls{CD} signal are intra-molecular or inter-molecular, the concentration dependence of the \gls{CD} spectrum of \cmpd+{ig2-15} in tetra\-hydro\-furan-aceto\-nitrile 50~\% - 50~\% was investigated (see Figure~\ref{fig:ig2-15-cd-thf-ch3cn} and data in Table~\ref{tab:soluzioni-cd}). 
A non linear dependence on concentration was found both for circular dichroism and for absorption intensities (confirmed by $g$ calculations, graph reported in appendix). The non linearity of \gls{UVvis} absorption is unexplained but the cooperative effect of the concentration on bisignate couplet amplitude suggests that the transition occurs inter-molecularly. 
Increasing the acetonitrile percent in tetra\-hydro\-furan-aceto\-nitrile mixtures results in a strong decrease of the \gls{CD} signal (see Figure~\ref{fig:ig2-15-cd-thf-ch3cn}) accompanied by a very small decrease in the \gls{UVvis} absorption peak, this is not unexpected and could be related to an extensive aggregation or to a change of the aggregate structure with coincident loss of the helical structure.\superfootcite{Zahn2002}

Adding 20~\% of methanol to a chloroform solution of \cmpd+{ig2-8} didn't generate any \gls{CD} signal, even after one day. Adding 30~\% of methanol resulted at first in a flat \gls{CD} spectrum, waiting \SI{10}{\minute} a \SI{6}{mdeg} \gls{CD} signal centered at \SI{441}{\nm} (inflection point of the bisignate couplet) was found (see Figure~\ref{fig:ig2-8-cd-meoh-1} and Table~\ref{tab:soluzioni-cd}) accompanied by a small shift in \gls{UVvis} absorption ($443 \rightarrow$~\SI{447}{\nm}). 
After waiting \SI{30}{\minute} this signal weakened at about half the intensity and clearly shifted to \SI{445}{\nm} while \gls{UVvis} absorption peak returned on the same wavelength of the first measurement. Finally after one day the signal was much more intense and centered at \SI{443}{\nm} (inflection point of the bisignate couplet) with a \gls{UVvis} peak at \SI{443}{\nm}. 
This indicates that some consecutive conformational changes are taking place in the material, each producing similar bisignate peak in \gls{CD} spectra. But the most important concept is that all these changes in \gls{CD} spectra are occurring without an extensive aggregation and planarization that would lead to a strong red shift in \gls{UVvis} absorption peak.
When 40~\% of methanol is added a \SI{16}{\nm} red shift is observed in the \gls{UVvis} main absorption peak ($443 \rightarrow$ \SI{459}{\nm}) and an usual bisignate dichroic peak centered at \SI{456}{\nm} (inflection point of the bisignate couplet) is observed. Both the shift in \gls{UVvis} and the shift in \gls{CD} spectra (check Figure~\ref{fig:ig2-8-cd-meoh-2} and Table~\ref{tab:soluzioni-cd}) are indicators of an extensive aggregation in the sample. 
This is supported by the quenching of \gls{UVvis} peak at \SI{265}{\nm} passing from 30~\% to 40~\% methanol (see Figure~\ref{fig:ig2-8-uvvis-meoh}, the peak at \SI{265}{\nm} probably is related to non conjugated thio\-phene rings, see page~\pageref{peak-265}). 
So a very similar di\-chroic signal can rise from both loosely and closely aggregated chains. Adding more methanol no shifts are observed, only an increase in \gls{CD} signal intensity. 

\begin{figure}[tbp]%ig2-8-cd-meoh-1
\begin{tikzpicture}
\begin{axis}[axis x line=bottom,axis y line=left,enlarge x limits=false,enlarge y limits=true,width=1\textwidth,height=8cm,xlabel=Wavelength (nm),ylabel=Molar circular dichroism $\Delta\epsilon$,xmin=240,legend pos=north west,cycle list name=linestyles*]
\addplot table[y expr=\thisrowno{1}*0.403] {img/results/ig2-8-cd-137-chcl3-70-meoh-30.txt};
\addlegendentry{in \ch{CHCl3} 70 \% - \ch{MeOH} 30 \%};
\addplot table[y expr=\thisrowno{1}*0.403] {img/results/ig2-8-cd-137-chcl3-70-meoh-30-10min.txt};
\addlegendentry{in \ch{CHCl3} 70 \% - \ch{MeOH} 30 \% \SI{10}{\minute}};
\addplot table[y expr=\thisrowno{1}*0.403] {img/results/ig2-8-cd-137-chcl3-70-meoh-30-30min.txt};
\addlegendentry{in \ch{CHCl3} 70 \% - \ch{MeOH} 30 \% \SI{30}{\minute}};
\addplot table[y expr=\thisrowno{1}*0.403] {img/results/ig2-8-cd-137-chcl3-70-meoh-30-1day.txt};
\addlegendentry{in \ch{CHCl3} 70 \% - \ch{MeOH} 30 \% \SI{1}{\day}};
\end{axis}
\end{tikzpicture} \caption[Circular dichroism of \cmpd+{ig2-8} in chloro\-form-methanol mixtures.]{Circular dichroism of \cmpd+{ig2-8} in chloro\-form-methanol mixtures. Concentration \SI{137}{\mg\per\liter}. Pathlength \SI{1}{\mm}.}
\label{fig:ig2-8-cd-meoh-1}
\end{figure}

\begin{figure}[tbp]%ig2-8-cd-meoh-2
\begin{tikzpicture}
\begin{axis}[axis x line=bottom,axis y line=left,enlarge x limits=false,enlarge y limits=true,width=1\textwidth,height=8cm,xlabel=Wavelength (nm),ylabel=Molar circular dichroism $\Delta\epsilon$,xmin=240,legend pos=north west,cycle list name=linestyles*]
\addplot table[y expr=\thisrowno{1}*0.403] {img/results/ig2-8-cd-137-chcl3-70-meoh-30-10min.txt};
\addlegendentry{in \ch{CHCl3} 70 \% - \ch{MeOH} 30 \% \SI{10}{\minute}};
\addplot table[y expr=\thisrowno{1}*0.403] {img/results/ig2-8-cd-137-chcl3-60-meoh-40.txt};
\addlegendentry{in \ch{CHCl3} 60 \% - \ch{MeOH} 40 \%};
\addplot table[y expr=\thisrowno{1}*0.403] {img/results/ig2-8-cd-137-chcl3-50-meoh-50.txt};
\addlegendentry{in \ch{CHCl3} 50 \% - \ch{MeOH} 50 \%};
\addplot table[y expr=\thisrowno{1}*0.403] {img/results/ig2-8-cd-137-chcl3-40-meoh-60.txt};
\addlegendentry{in \ch{CHCl3} 40 \% - \ch{MeOH} 60 \%};
\addplot table[y expr=\thisrowno{1}*0.403] {img/results/ig2-8-cd-137-chcl3-30-meoh-70.txt};
\addlegendentry{in \ch{CHCl3} 30 \% - \ch{MeOH} 70 \%};
\end{axis}
\end{tikzpicture} \caption[Circular dichroism of \cmpd+{ig2-8} in chloro\-form-methanol mixtures.]{Circular dichroism of \cmpd+{ig2-8} in chloro\-form-methanol mixtures. Concentration \SI{137}{\mg\per\liter}, pathlength \SI{1}{\mm}.}
\label{fig:ig2-8-cd-meoh-2}
\end{figure}

\cmpd+{ig2-8} in chloro\-form-aceto\-nitrile mixtures was tested giving the usual bisignate peak. In these \gls{CD} measurements we found a maxima for bisignate couplet intensity at 70~\% of acetonitrile (see Table~\ref{tab:soluzioni-cd}, this phenomenon is the same observed for \cmpd+{ig2-15} in Figure~\ref{fig:ig2-15-cd-thf-ch3cn}); the addition of less or more acetonitrile resulted in a weaker dichroic signal without a change in \gls{UVvis} absorption (for all the samples, the concentration was fixed to \SI{137}{\mg\per\liter}). 
As noticed at page~\pageref{uvvis-ch3cn}, \gls{UVvis} main peak in chloro\-form-aceto\-nitrile mixtures is \SI{\approx5}{\nm} red shifted in comparison with peak from chloro\-form-methanol mixtures while a smaller shift was observed in \gls{CD} signal inflection point.

Peaks positions in \gls{UVvis} and \gls{CD} spectra in tetra\-hydro\-furan-aceto\-nitrile mixtures of \cmpd+{ig2-8} are similar to spectra in chloro\-form-aceto\-nitrile (see Table~\ref{tab:soluzioni-cd}). The intensity of the classical bisignate \gls{CD} peaks vary only to a small extent changing solvent composition.

Looking to every presented \gls{CD} spectrum we can notice that the integral of each bisignate couplet is positive, the inflection point of the peaks was always located at positive values by $\approx10~\%$ of the total signal amplitude. This positive-biased signal can be interpreted as the sum of the conservative bisignate excitonic couplet and an overlying positive monosignate term. 
Another feature recognizable in every spectrum is a shift of a few nanometers (\SI{\approx10}{\nm}) towards high energies of bisignate couplet inflection point compared to the \gls{UVvis} absorption peak. This could suggest for the birth of \gls{CD} signal from non well planarized chains.

\label{majority}
Comparing the intensities of \gls{CD} signals, measured at the same concentration and solvent mixture, from \cmpd+{ig2-15} and \cmpd+{ig2-8} in Table~\ref{tab:soluzioni-cd} we can notice that \cmpd+{ig2-8} always shows a higher molar circular dichroism than \cmpd+{ig2-15}. A mean of the ratio of intensities show that \cmpd+{ig2-15} \gls{CD} signals are around 50~\% less intense than \gls{CD}s from \cmpd+{ig2-8}. 
Knowing that the chiral alcohol used for the synthesis of \cmpd+{ig2-10} and \cmpd+{ig2-15} has an optical purity of 59~\% (optical purity is a good estimation for the enantiomeric excess) we can affirm that no majority rule\nota{\textit{majority rules}, cooperative effects of the monomer units along the polymer backbone that results in a nonlinear relation between the specific optical rotation and the enantiomeric excess of chiral units present in the polymer.} 
effect can be seen\superfootcite{Langeveld-Voss1999}.

Looking at chiral poly\-thio\-phenes bearing a side chain with a different number of carbon atoms in between backbone and (S)-2-methyl\-butyl\-oxy group we can notice the odd-even effect.\superfootcite{Bouman1994}\superfootcite{Bouman1994a} The odd-even effect indicates an inversion of the bisignate couplet when changing the chiral center distance by one atom.\superfootcite{Lermo1999}

The maximum and minimum measured dissymmetry factors (also known as chiral anisotropy factor) $g(\lambda) = \Delta\epsilon(\lambda) / \epsilon(\lambda)$ for \cmpd+{ig2-15} are \SI{\approx4.7e-4}{} and \SI{\approx-3.2e-4}{} respectively at 512 and \SI{412}{\nm} from solution in \ch{CHCl3} 50~\% - \ch{MeOH} 50~\%. For \cmpd+{ig2-8} maximum and minimum $g$ values are \SI{\approx6.5e-4}{} and \SI{\approx-4.7e-4}{} respectively at 515 and \SI{405}{\nm}. 
These values are lower but comparable with data from other similar poly\-thio\-phenes reported in the literature.\superfootcite{Salatelli2010}\superfootcite{Andreani1998}\superfootcite{Langeveld-Voss1998a} This, together with the absence of fine structure in \gls{CD} spectra, suggests that chirality isn't coming from intermolecular interactions of well organized structures, rather these weak \gls{CD} spectra could be related to intramolecular interactions. 
\end{subsubsection}
\begin{subsubsection}{In Solid State}

Samples in mixtures of good and poor solvents are a handy way of studying aggregation but the real interest is on the solid state. 

Unfortunately obtaining good and homogeneous solid state samples is a time consuming task and often linear dichroism artifacts are present and interfere with \gls{CD} measurement. 
Only spectra with low distortion by linear dichroism were presented and in order to further reduce the linear dichroism artifacts, a mean on at least two different orientations (\SI{0}{\degree} and \SI{90}{\degree}) was reported in each case. Moreover in solid state it's more difficult to obtain reproducible results, due to the dependence on sample preparation, doping\superfootcite{Goto2002a} and even thermal history.\superfootcite{Bouman1995}
All the numeric data for \gls{UVvis} and \gls{CD} characterization in solid state is reported in Table~\ref{tab:film-cd}. 

\begin{table}[tbp]%film-cd
\centering
\caption[UV-vis and CD data for \cmpd+{ig2-4}, \cmpd+{ig2-15} and \cmpd+{ig2-8} films on quartz.]{\gls{UVvis} and \gls{CD} data for \cmpd+{ig2-4}, \cmpd+{ig2-15} and \cmpd+{ig2-8} films on quartz.}\label{tab:film-cd}
\begin{tabular}{c|c|c|c|c|c}
\toprule
polymer&processing&\tallcell{UV-vis\\ peak\\(\SI{}{nm})}&\tallcell{UV-vis\\abs.}&\tallcell{CD\\ inflection\\(\SI{}{nm})}& \tallcell{CD\\ amplitude\\(\SI{}{mdeg})} \\ \cmidrule{1-6}
\multirow{4}{*}{\cmpd+{ig2-4}}&\ch{CHCl3} solution				&430&	&	&	\\ 
&\ch{CHCl3} 50 \% - \ch{MeOH} 50 \%	&$\approx433$	& &&\\
&spin \SI{700}{\rpm} from \ch{CHCl3}	&$\approx462$	&0.28	&	449 		&1.3 	\\ 
&in \ch{KCl}			&$\approx464$	&	&			&		\\ \cmidrule{1-1}
\multirow{4}{*}{\cmpd+{ig2-15}}&\ch{CHCl3} solution				&446&	&	&\\
&\ch{CHCl3} 50 \% - \ch{MeOH} 50 \%	&$\approx466$	& &459&\\
&cast from \ch{CH2Cl2}			&$\approx486$	&0.29	&	469 		&6.6		\\
&in \ch{KCl}			&$\approx493$	&	&			&		\\ 
\cmidrule{1-1}
\multirow{8}{*}{\cmpd+{ig2-8}}&\ch{CHCl3} solution 				&444			&	&&	\\
&\ch{CHCl3} 50 \% - \ch{MeOH} 50 \%	&461&0.76&457&15.6\\
&cast from \gls{THF}					&$\approx473$	&0.40&463&1.8 \\
&cast from \ch{CH2Cl2}				&$\approx469$	&0.45&462&4.6	\\
&spin \SI{700}{\rpm} from \ch{CHCl3}	&$\approx476$	&0.17&468&0.6	\\
&spin \SI{700}{\rpm} from \gls{THF}		&$\approx476$	&0.10&453&0.5	\\
&spin \SI{1000}{\rpm} from \ch{CHCl3}&$\approx492$	&0.13&474&0.6	\\
&in \ch{KCl}			&$\approx495$	&	&			&		\\ 
\bottomrule
\end{tabular} 

\smallskip

{\gls{UVvis} peak position is approximated due to the huge full width at half maximum.}

\end{table}

Measuring circular dichroism of \cmpd+{ig2-4} thin film spin coated from a chloroform solution (\SI{700}{\rpm}) we can see a negatively biased bisignate couplet similar to the \gls{CD} spectrum obtained from the \cmpd+{ig2-4} solution in \ch{CHCl3} 40~\% - \ch{MeOH} 60~\%. 
The \gls{CD} signal from thin film is at shorter wavelength compared to its \gls{UVvis} absorption peak (\gls{CD} inflection point \SI{449}{\nm}, \gls{UVvis} peak \SI{\approx462}{\nm}, see Table~\ref{tab:film-cd}).

\cmpd+{ig2-15} thin film cast from dichloromethane showed a \gls{CD} signal with dissymmetry factor $g$ higher than samples in solution. We already noted (at page~\pageref{ig2-15-uvvis-film}) a \SI{20}{\nm} red shift in \gls{UVvis} peak from \ch{CHCl3} 50~\% - \ch{MeOH} 50~\% mixture to cast film (\SI{466}{\nm} $\rightarrow$ \SI{486}{\nm}); in \gls{CD} spectrum a red shift of \SI{10}{\nm} can be observed (\SI{459}{\nm} $\rightarrow$ \SI{469}{\nm}).

\begin{figure}[tbp]%ig2-8-cd-film
\begin{tikzpicture}
\begin{axis}[axis x line=bottom,axis y line=left,enlarge x limits=false,enlarge y limits=true,width=1\textwidth,height=8cm,xlabel=Wavelength (nm),ylabel=Dissymmetry factor $g$,xmin=260,xmax=550,legend style={font=\footnotesize},legend pos=south east,cycle list name=linestyles*]
\addplot table {img/results/ig2-8-g-cast-ch2cl2.txt};
\addlegendentry{cast from \ch{CH2Cl2}};
\addplot table {img/results/ig2-8-g-cast-thf.txt};
\addlegendentry{cast from \gls{THF}};
\addplot table {img/results/ig2-8-g-spin700-thf.txt};
\addlegendentry{spin \SI{700}{\rpm} from \gls{THF}};
\addplot table {img/results/ig2-8-g-spin700-chcl3.txt};
\addlegendentry{spin \SI{700}{\rpm} from \ch{CHCl3}};
\addplot table {img/results/ig2-8-g-spin1000-chcl3.txt};
\addlegendentry{spin \SI{1000}{\rpm} from \ch{CHCl3}};
\end{axis}
\end{tikzpicture} \caption[Dissymmetry factor of \cmpd+{ig2-8} in thin films.]{Dissymmetry factor $g$ (also known as chiral anisotropy factor) of polymer \cmpd+{ig2-8} in thin film from casting or spin coating. $g(\lambda) = \Delta\epsilon(\lambda) / \epsilon(\lambda)$.}
\label{fig:ig2-8-cd-film}
\end{figure}

Many films of \cmpd+{ig2-8} were produced, both by casting and spin coating. Both \gls{UVvis} and \gls{CD} peaks are red shifted comparing each film to each sample in good-poor solvent mixture. Cast films from tetrahydrofuran and dichloromethane showed similar \gls{UVvis} absorption 
and \gls{CD} peak inflection point (see Figure~\ref{fig:ig2-8-cd-film} and Table~\ref{tab:film-cd}) but the film from dichloromethane showed less scattering and a higher bisignate peak intensity.

Films from spin coating at slow rotational speeds (\SI{700}{\rpm}) resulted in similarly shaped \gls{UVvis} spectra while the position of 
\gls{CD} signals changes by varying the deposition solvent (see Table~\ref{tab:film-cd}). Spin coating at higher speed resulted in a strongly red shifted absorption peak (\SI{492}{\nm}) and a consequently red shifted \gls{CD} peak inflection point (\SI{474}{\nm}). This strong red shift was associated with chain stretching on quartz substrate due to high rotational speed.

Dipping the spin coated films in pentane induced no change in \gls{UVvis} main peak wavelength (see Section~\ref{uvvis-pentane}). In the literature this process is reported to change the \gls{CD} solubilizing the lateral chains and allowing a reorganization,\superfootcite{Caronna1997}\superfootcite{Salatelli2010} but our spectra were influenced only slightly (spectra in appendix).

\begin{SCfigure}[][tbp]%cd-kcl
\begin{tikzpicture}
\begin{axis}[axis x line=bottom,axis y line=left,enlarge x limits=false,enlarge y limits=true,width=0.85\textwidth,height=6cm,xlabel=Wavelength (nm),ylabel=Circular Dichroism (mdeg),xmin=240,xmax=580,legend pos=north west,cycle list name=linestyles*]
\addplot table {img/results/ig2-4-cd-kcl.txt};
\addlegendentry{\cmpd+{ig2-4} in \ch{KCl}};
\addplot table {img/results/ig2-15-cd-kcl-3.txt};
\addlegendentry{\cmpd+{ig2-15} in \ch{KCl}};
\addplot table {img/results/ig2-8-cd-kcl-2.txt};
\addlegendentry{\cmpd+{ig2-8} in \ch{KCl}};
\end{axis}
\end{tikzpicture} \caption{Circular dichroism spectra of polymers \cmpd+{ig2-4}, \cmpd+{ig2-15} and \cmpd+{ig2-8} in \ch{KCl} disks.}
\label{fig:cd-kcl}
\end{SCfigure}

Then bulk solid polymer samples were used to make \ch{KCl} disks. These resulted in poorly dispersed solid giving a huge scattering, nevertheless we can see a weak \gls{CD} (see Figure~\ref{fig:cd-kcl}). Only \cmpd+{ig2-8} didn't show the bisignate couplet, but this can be related to poor quality of obtained disks. 
\end{subsubsection}
\end{subsection}
\begin{subsection}{Thermal Characterization}

\Gls{DSC} characterization was performed. Heating both polymers \cmpd+{ig2-15} and \cmpd+{ig2-8} the only observed thermal transition is the fusion at \SI{198}{\celsius}. By cooling the samples at \SI{10}{\celsius\per\minute} only \cmpd+{ig2-15} showed a crystallization peak at \SI{180}{\celsius} while \cmpd+{ig2-8} didn't crystallize (\gls{DSC} curves are reported in appendix).

\end{subsection}
\clearpage
\begin{subsection}{X-ray Crystallography Characterization}
\label{xrd}

In the literature is reported that the replacement of the second methylene group in the side chain of \gls{P3HT} by an oxy\-gen atom prohibits a close packing of the chains.\superfootcite{Zoombelt2008} In the present case the interdigitation should be further hampered by the branching of our side chain (which is needed for the introduction of chirality). Polythio\-phenes with poly\-ether side groups are reported to be amorphous from \gls{XRD} studies.\superfootcite{Lin2012} 

\begin{figure}[tbp]%xrd
\begin{tikzpicture}
\begin{axis}[axis x line=bottom,axis y line=left,enlarge x limits=false,enlarge y limits=false,width=1\textwidth,height=8cm,yticklabels={,,},xlabel=Degrees $2\theta$,ylabel=Intensity (a.\ u.),xmin=2]
\addplot[densely dotted] table {img/results/ig2-4-xrd.txt};
\addlegendentry{\cmpd{ig2-4}};
\addplot[line width=0.5pt] table[y expr=\thisrowno{1}-150] {img/results/ig2-8-xrd.txt};
\addlegendentry{\cmpd{ig2-8}};
\addplot[line width=1.5pt] table[y expr=\thisrowno{1}/35-350] {img/results/ig2-8-xrd-cast-ch2cl2-Perp_32_00001IP.txt};
\addlegendentry{\cmpd{ig2-8} cast from \ch{CH2Cl2}};
\end{axis}
\end{tikzpicture} \caption{X-ray diffraction of \cmpd+{ig2-4} and \cmpd+{ig2-8} as powder, and of \cmpd+{ig2-8} as cast film from synchrotron radiation.}
\label{fig:xrd}
\end{figure}

Figure~\ref{fig:xrd} shows \gls{WAXD} patterns of \cmpd+{ig2-4} and \cmpd+{ig2-8} in bulk solid and \cmpd+{ig2-8} cast from \ch{CH2Cl2}. The last spectrum was obtained from synchrotron radiation grazing incidence X--ray diffraction (at \SI{0.5}{\degree}, the reported spectra is the in-plane scattering). This thin film was chosen for analysis because of the high dissymmetry factor $g$ showed. The patterns exhibited an amorphous halo centered around \SI{21}{\degree}, more intense for the thin film (this is in part caused by the quartz substrate) than for bulk solid. 
Four clear peaks are present in each \gls{XRD} spectra, typical peaks of a regioregularly substituted poly\-thio\-phene well interdigitated and stacked in lamellae.\superfootcite{Chen1995}\superfootcite{Winokur1989} The first peak at a lower diffraction angle of $2\theta$ \SI{\approx5.4}{\degree}, shows the distance of the chains in the interdigitation direction (Miller index (100)), the second and the third ($2\theta =$ \SI{\approx10.5}{\degree} and $2\theta =$ \SI{\approx15.5}{\degree}) represent second and third-order reflections ((200) and (300)). 
From these values using the Bragg's law we can derive a mean distance of \SI{1.68}{\nm} between chains in the interdigitation direction. This value is quite similar to that reported for poly\-thio\-phene with an hexyl side group (\gls{P3HT}).\superfootcite{Chen1995}\superfootcite{Winokur1989} 
The fourth peak at \SI{\approx23.5}{\degree} is related to the inter\-layer $\pi-\pi$ stacking direction ((020), \SI{0.38}{\nm}). 
The thin film was annealed heating up to \SI{100}{\celsius} and cooling slowly at \SI{1}{\celsius\per\minute}. \Gls{WAXD} analysis showed no difference in the (100) peak before and after annealing (spectra reported in appendix). 

\begin{figure}[tbp]%ig2-8-xrd2d-cast-ch2cl2
\centering
\includegraphics[width=0.8\textwidth]{ig2-8-xrd2d-cast-ch2cl2.png}
\caption{Bidimensional X-ray diffraction of \cmpd+{ig2-8} cast from \ch{CH2Cl2}.}
\label{fig:ig2-8-xrd2d-cast-ch2cl2}
\end{figure}

The thin film of \cmpd+{ig2-8} cast from \ch{CH2Cl2} was analyzed also in \gls{XRD2D} and the resulting image is reported in Figure~\ref{fig:ig2-8-xrd2d-cast-ch2cl2}. The presence of arcs indicates the absence of orientation on quartz substrate. 

Solid \cmpd+{ig2-8} was analyzed also using \gls{SAXS} and no long range spacing was observed. This information is contrary to the hypothesis of chirality originating from a large supramolecular assembly made of tilted lamellae in crystalline solid.\superfootcite{Goto2002} This is in accordance with the absence of vibronic bands in \gls{CD} spectra that should be present if the chirality rises from well ordered structures. We don't exclude the presence of a chiral supramolecular organization without ordered repetition in amorphous domains. 

\end{subsection}
\clearpage
\begin{subsection}{Photovoltaic Device Tests}
\label{cella}

\cmpd+{ig2-8} was tested in solar cells in blend with \acrfull{PCBM} nanoparticles by the Centre for Hybrid and Organic Solar
Energy (\textsmaller{CHOSE}) of the University of Rome Tor Vergata. Preliminary results are reported here, tests on \cmpd+{ig2-15} are in course. 

\begin{figure}[tbp]%ig2-11-pcbm-uvvis
\begin{tikzpicture}
\begin{axis}[axis x line=bottom,axis y line=left,enlarge x limits=false,enlarge y limits=false,width=1\textwidth,height=6cm,xlabel=Wavelength (nm),ylabel=Absorbance,legend pos=north east,cycle list name=linestyles*]
\addplot table {img/results/ig2-11-pcbm-uvvis-500rpm.txt};
\addlegendentry{\SI{500}{\rpm} (\SI{65}{\nm})};
\addplot table {img/results/ig2-11-pcbm-uvvis-1000rpm.txt};
\addlegendentry{\SI{1000}{\rpm} (\SI{55}{\nm})};
\addplot table {img/results/ig2-11-pcbm-uvvis-1500rpm.txt};
\addlegendentry{\SI{1500}{\rpm} (\SI{45}{\nm})};
\end{axis}
\end{tikzpicture} \caption[UV-vis spectra of thin films from spin coating of \cmpd+{ig2-8}:PCBM blend.]{\Gls{UVvis} spectra of thin films from spin coating of \cmpd+{ig2-8}:\gls{PCBM} blend at various rotation speeds (and film thickness). Spinning duration \SI{120}{\s}. Support soda-lime glass.}
\label{fig:ig2-11-pcbm-uvvis}
\end{figure}

The polymer was dissolved in ortho-di\-chloro\-benzene (1~\% weight) with \gls{PCBM} in a 1:1 ratio. 
This solution was used for producing thin films via spin coating. The rotation speed was varied (\SI{500}{\rpm}, \SI{1000}{\rpm} and \SI{1500}{\rpm}) in order to obtain different film thicknesses. 
\Gls{UVvis} absorption spectra (and thickness) of the obtained thin films are reported in Figure~\ref{fig:ig2-11-pcbm-uvvis}. The obtained thickness were much smaller than expected, thus we utilized \SI{400}{\rpm} for the cell production, that is the lower limit needed to obtain an uniform deposition of the film via spin coating. 
We can notice an intense absorption peak at $\lambda = $ \SI{335}{\nm} due to the \gls{PCBM} absorption. The polythio\-phene absorption peak is located at \SI{455}{\nm}, thus it was somewhat blue shifted respect to our thin films from spin coating from pure \cmpd+{ig2-8}, this is a reasonable consequence of the disorder introduced with \gls{PCBM}. 

\begin{SCfigure}[][tbp]%cella
\includegraphics[width=0.4\textwidth]{cella-bn.png}
\caption{Testing cell layers scheme.}
\label{fig:cella}
\end{SCfigure}

The structure of the realized cells is reported in Figure~\ref{fig:cella}. 
The utilized architecture was the direct one, that is: \gls{FTO}\-/\acrshort{PEDOT:PSS} (\SI{50}{\nm})\-/active layer\-/\ch{Ca} (\SI{20}{\nm})\-/\ch{Al} (\SI{100}{\nm}) where \gls{FTO} is the anode and aluminium is the cathode. 
Refer to experimental section for cell preparation procedure. 
Eight cells were produced and a mean of the results is reported here. 
\gls{P3HT}:\gls{PCBM} was used in the active layer for reference cells. These reference cells were annealed at \SI{150}{\celsius} for \SI{10}{\minute} while the \cmpd+{ig2-8} based cells weren't annealed. 

\begin{table}[tbp]%pv
\centering
\caption[Characteristic parameters of solar cells from polymer-PCBM blends.]{Characteristic parameters of solar cells from polymer-\gls{PCBM} blends.}\label{tab:pv}
\begin{tabular}{c|c|c|c|c}
\toprule
polymer&$V_{oc}$ (\SI{}{\mV})&$J_{sc}$ (\SI{}{\mA\per\square\cm})&FF (\%)&$\eta$ (\%) \\ \cmidrule{1-5}
\gls{P3HT}&$0.567 \pm 0.007$&$-10.8 \pm 0.7$&$50 \pm 20$	& $3 \pm 1$ \\ 
\cmpd+{ig2-8}&$0.68 \pm 0.05$&$-3.2 \pm 0.5$&$31 \pm 3$&$0.7 \pm 0.1$\\ 
\bottomrule
\end{tabular}
\end{table}

From results in Table~\ref{tab:pv} we notice that \cmpd+{ig2-8}:\gls{PCBM} cells under illumination showed an open circuit voltage $V_{oc}$ of \SI{0.68}{\mV}, higher than that of \gls{P3HT} (\SI{0.57}{\mV}). This value is in accordance with \citeauthor*{Zoombelt2008}\superfootcite{Zoombelt2008} where a similar, but non branched, polythio\-phene side group was used. 
This is caused by the electron withdrawing effect of the oxy\-gen atom in side chain respect to the methylene in \gls{P3HT}. The measured short circuit current density $J_{sc}$ from \cmpd+{ig2-8} was lower than the current from \gls{P3HT}. We think that this is caused by worse crystallization due the presence of branching in side group and by the presence of the oxy\-gen atom, as reported by \citeauthor*{Zoombelt2008}.\superfootcite{Zoombelt2008} 
The obtained \gls{FF} was rather low (31~\%) and the resulting power conversion efficiency $\eta$ resulted 0.7~\% versus and efficiency of 
3~\% from the reference cells. 
Considering the high $V_{oc}$ observed this is a promising candidate for block copolymer synthesis and its use in bulk heterojunction solar cells active layer. 

\end{subsection}
\end{section}
\clearpage
\begin{section}{Second Block}

Two main approaches are common in the synthesis of diblock copolymers: grafting onto, which means joining two blocks after their polymerization, and grafting from, which means starting the polymerization of the second block from a termination of the first block. The grafting onto way allows one to polymerize both blocks in a clean reaction environment and permits a better characterization of the single blocks prior to coupling. 
But in the present case the grafting onto approach is hampered by the chelating behavior of the \gls{P4VP} block that would interfere with the joining reaction. 
Thus a grafting from approach was followed, polymerizing the \gls{P4VP} block starting from an adequately functionalized end of \cmpd+{ig2-8}. 
Two poly\-thio\-phene-\textit{block}-poly\-vinyl\-pyridine syntheses were already present in the literature. \Citeauthor*{Lohwasser2012}\superfootcite{Lohwasser2012} synthesized an alkyne-terminated poly\-thio\-phene, then, using a click reaction, they functionalized the alkyne termination with an initiator for poly\-vinyl\-pyridine polymerization. 
We decided to don't follow this route because the formation of the alkyne end capper is reported to produce both mono and di terminated poly\-thio\-phene chains.\superfootcite{Jeffries-EL2004}\superfootcite{Jeffries-El2005} Moreover the triple bond is sensible to the standard quenching procedure and there are questions on the efficiency of the click reaction on a triple bond conjugated to a poly\-thio\-phene chain.\superfootcite{Urien2008} 
\Citeauthor*{Sary2010a}\superfootcite{Sary2010a} synthesized an \ch{H/Br} terminated poly\-thio\-phene, then they employed the halogenated termination to functionalize the polymer with an aldehyde group. Then the polythio\-phene-aldehyde was used as a quenching agent for vinyl\-pyridine anionic polymerization. This route was discarded because of the strict reaction conditions needed by an anionic polymerization.

For the synthesis of the first polythio\-phenic block we followed polymerization methods which were reported producing high yields in \ch{H/Br} terminated chains. 
Then we tried to functionalize the halogenated endings of our polymers with a boronic ester. We reacted \cmpd+{ig2-8} with \textit{n}-butyl\-lithium in \gls{THF} at low temperature, then we added 2-iso\-prop\-oxy-4,4,5,5-tetra\-methyl-1,3,2-di\-oxa\-borolane in order to produce a pinacol ester of the poly\-thio\-phene boronic acid. Although trying thrice this functionalization with different \textit{n}-butyl\-lithium quantities, we didn't get the boron atom bonded to the polymer (this was confirmed by \gls{FTIR}, {\HNMR} and {\CNMR} spectroscopies and later by \gls{MALDI}). 

So we followed a procedure similar to the strategy from \citeauthor*{Sary2010a}\superfootcite{Sary2010a} but avoiding the anionic polymerization. In the first step we functionalized the polymer \cmpd+{ig2-8} on the halogenated sites using a Suzuki coupling. The functionalizing group was chosen as a di\-thio\-phene bearing a boronic acid group and a formyl group. We preferred using a thio\-phenic end capper because the similarity of its electronic levels to the levels of the rest of the chain should avoid the possible charge trapping on the end capper. Then this aldehyde-terminated chiral poly\-thio\-phene was reacted with the lithium activated polymerization initiator. Finally a poly\-(4-vinyl\-pyridine) block was polymerized starting from this group. 

\begin{subsection}{Radical Mediator Synthesis and Characterization}
The polymerization of \gls{P4VP} was a living radical polymerization mediated by the \gls{TIPNO} radical.\superfootcite{Lohmeijer2005} This radical is a commercial product but we need to build an initiator molecule (containing the radical) suitable to be attached on a poly\-thio\-phene end. We synthesized \gls{BrPhEtTIPNO} \cmpd+{ig2-21} which will be lithiated on the bromine atom for the linking to an aldehyde-terminated poly\-thio\-phene. 

\begin{figure}[tbp]%syn-tipno
\centering
\includegraphics[scale=0.7]
{syn6-7-riduzione-bromurazione.pdf}
\medskip
\includegraphics[scale=0.7]
{syn8-tipno-br.pdf} 
\caption[PhEt-TIPNO synthesis.]{\Acrfull{BrPhEtTIPNO} \cmpd+{ig2-21} synthesis.}
\label{fig:syn-tipno}
\end{figure}

As shown in Figure~\ref{fig:syn-tipno} we started reducing 4-bromo\-aceto\-phenone to 1-(4-bromo\-phenyl)\-ethanol \cmpd+{ig2-16} via a simple reaction with sodium boro\-hydride.\superfootcite{Wiles2006} Then we brominated\superfootcite{Stalmach2001} 1-(4-bromo\-phenyl)\-ethanol \cmpd+{ig2-16} with phosphorus tri\-bromide (\ch{PBr3}) obtaining a mixture of 1-bromo-4-(1-bromo\-ethyl)\-benzene \cmpd+{ig2-17} and 4-bromo\-styrene. This mixture was reacted with \gls{TIPNO} radical in a copper(I) catalyzed coupling.\superfootcite{Benoit1999} The purification was carried on basic allumina (in order to avoid the evenience of a breaking of the nitr\-oxyde group on the slightly acid silica). 

The resulting product was characterized with poor results using \gls{GCMS} with a modified on purpose method (injector and transfer line temperatures were lowered), then the structure was confirmed by {\HNMR}, the assignation is reported in experimental section and the spectra is reported in appendix.

\end{subsection}
\begin{subsection}{Functionalization of Polythio\-phene with Radical Mediator}
\begin{figure}[tbp]%syn9-suzuki
\centering
\includegraphics[width=1\textwidth]
{syn9-suzuki.pdf}
\caption{End-functionalization of poly\-thio\-phene with aldehyde using a Suzuki reaction.}
\label{fig:syn9-suzuki}
\end{figure}

The poly\-thio\-phene \cmpd+{ig2-8} was functionalized with an aldehyde group via palladium catalyzed Suzuki coupling with B-(5'-formyl\-[2,2'-bi\-thio\-phen]\-5-yl)\-boronic acid, as schematized in Figure~\ref{fig:syn9-suzuki}. We followed a procedure reported in literature\superfootcite{Zhang2011}\superfootcite{Imae2012} but with a significant variation: we substituted the tetra\-kis\-(tri\-phenyl\-phosphine)\-palladium(0) (\ch{Pd(PPh3)4}) catalyst with the more stable bis\-(tri\-phenyl\-phosphine)\-palladium(II) di\-chloride (\ch{Pd(PPh3)2Cl2}). We found (in a test reaction involving a low molecular weight halogenated thio\-phene) that this convenient catalyst works for our coupling. 
The functionalized polymer \cmpd+{ig2-19} was purified via precipitation, washing and redissolved in various solvents. The presence of an aldehyde group was confirmed by \gls{FTIR} spectroscopy and {\HNMR} (in \gls{TCE}-d$_2$) by the appearance of a peak at \SI{9.80}{\ppm} attributable to carbonylic proton. 

Then the \gls{BrPhEtTIPNO} \cmpd+{ig2-21} was lithiated with one equivalent of \textit{n}-buthyl\-lithium in \gls{THF} and poured in well dissolved (in \gls{THF} at \SI{50}{\celsius}) poly\-thio\-phene-aldehyde \cmpd+{ig2-19} obtaining the macroinitiator poly\-thio\-phene-\gls{TIPNO} \cmpd+{ig2-22}. The reaction yield was estimated from the decrease of the aldehyde peak in {\HNMR} spectrum and resulted $\approx90$~\% (see spectra in appendix, also new peaks in aromatic region can be noticed). The lithiation completeness was verified by quenching with water an aliquot of lithium activated \gls{BrPhEtTIPNO}, the \gls{GCMS} analysis showed the absence of brominated compounds.

\begin{figure}[tbp]%syn10-pt-tipno
\centering
\includegraphics[scale=0.5]
{syn10-pt-tipno.pdf}
\caption[Functionalization of poly\-thio\-phene-aldehyde \cmpd+{ig2-19} with PhEt-TIPNO \cmpd+{ig2-21} obtaining the macroinitiator \cmpd+{ig2-22}.]{Functionalization of poly\-thio\-phene-aldehyde \cmpd+{ig2-19} with \gls{BrPhEtTIPNO} \cmpd+{ig2-21} obtaining the macroinitiator \cmpd+{ig2-22}.}
\label{fig:syn10-pt-tipno}
\end{figure}

\end{subsection}
\begin{subsection}{Second Block Synthesis and Characterization}
\label{sec:nmrp}

\begin{figure}[tbp]%syn11-p4vp
\centering
\includegraphics[width=1\textwidth]
{syn11-p4vp.pdf}
\caption[NMRP polymerization of P4VP starting from macroinitiator \cmpd+{ig2-22} and obtaining block copolymer \cmpd+{sz17}.]{\gls{NMRP} polymerization of \gls{P4VP} starting from macroinitiator \cmpd+{ig2-22} and obtaining block copolymer \cmpd+{sz17}.}
\label{fig:syn11-p4vp}
\end{figure}

\Acrfull{NMRP} polymerization of vinyl\-pyridine was performed following procedure from \citeauthor*{Lohwasser2012}.\superfootcite{Lohwasser2012} The radical polymerization was initiated by heating at \SI{125}{\celsius} the solution of macroinitiator \cmpd+{ig2-22} in ortho-di\-chloro\-benzene containing 4-vinyl\-pyridine in a molar ratio of 1:3000 with addition of free radical \gls{TIPNO} (enhances the control on the polymerization). The polymerization was run for \SI{4}{\hour}. The molar ratio and the duration were modified with respect to the reference article in order to reduce the formation of \gls{P4VP} homo\-polymer. Then the reaction was cooled down, diluted in chloroform and poured in hexane. 
Hexane is a poor solvent for both poly\-thio\-phene and poly\-vinyl\-pyridine, thus the solid was filtered and recovered in chloroform, that is a good solvent for both blocks. The chloroform solution was poured again in methanol, that is a good solvent only for poly\-vinyl\-pyridine block. {\HNMR} characterization of this methanol fraction showed the presence of vinyl\-pyridine homo\-polymers and small molecules, no peaks for thio\-phenic units. 
The solid was filtered and washed with hot toluene (\SI{70}{\celsius}) that is a good solvent (at room temperature is poor) for the poly\-thio\-phene block only. With {\HNMR} we verified the absence of poly\-vinyl\-pyridine aromatic peaks looking both for aromatic and aliphatic protons (spectra with expansions in appendix). The remaining solid was recovered in chloroform obtaining the block copolymer \cmpd+{sz17}. {\HNMR} analysis on this fraction, reported in Figure~\ref{fig:sz17-nmr-h-zoom-py}, showed the presence of broad aromatic (\SI{8.26}{\ppm} and \SI{6.35}{\ppm} in \gls{TCE}-d$_2$) and aliphatic (\SI{0.83}{\ppm} and \SI{0.79}{\ppm} in \gls{TCE}-d$_2$) peaks assigned to a short block of poly\-vinyl\-pyridine together with the usual peaks from poly\-thio\-phene block. 
Comparing with \gls{NMR} spectrum from \cmpd+{ig2-22} we can notice also variations in peaks around \SI{7.4}{\ppm} that are related to the \gls{TIPNO} group. A confirmation of the removal of unreacted thio\-phenic homo\-polymer come from an aldehydic {\HNMR} peak: The impurities of unreacted polythio\-phene-aldehyde \cmpd+{ig2-19} (in the functionalization reaction with lithiated \gls{BrPhEtTIPNO} \cmpd+{ig2-21}) can't work as a macroinitiator for a \gls{P4VP} block, the fact that in the toluene fraction we can see an aldehyde \gls{NMR} peak at \SI{9.80}{\ppm} that is 5 times more intense than the corresponding peak in the purified fraction means that the majority of homo\-polymer was washed in hot toluene. 

The purification procedure reported by the reference article\superfootcite{Lohwasser2012} was more complex and not suitable for small quantities (we worked with less than \SI{20}{\mg} of product). 
Indeed our simpler purification method 
was effective. The relative quantity of pyridine rings can be estimated by integration of aromatic protons, from this calculation results that for every pyridine ring there are 5 thio\-phene rings, thus around 10 pyridine rings per polythio\-phene chain. The length of this block is inferior to that reported from similar syntheses.\superfootcite{Lohwasser2012}\superfootcite{Mougnier2012}

\begin{figure}[tbp]%sz17-nmr-h-zoom-py
\includegraphics[width=0.99\textwidth]{sz17-nmr-h-zoom-py.pdf}
\caption[{\HNMR} of poly\-thio\-phene-poly\-vinyl\-pyridine \cmpd+{sz17}.]{{\HNMR} (\SI{600}{\MHz}) of poly\-thio\-phene-poly\-vinyl\-pyridine \cmpd+{sz17} in \gls{TCE}-d2.}
\label{fig:sz17-nmr-h-zoom-py}
\end{figure}

Preliminary optical characterizations were performed on the new di\-block copolymer \cmpd+{sz17}. \Gls{UVvis} spectroscopy in chloroform showed a spectra similar to \cmpd+{ig2-8} but the intensity was halved, maybe due the presence of \gls{P4VP} (or due to non proper drying). 
Adding 40~\% methanol to a chloroform solution resulted in a \SI{12}{\nm} red shift in main \gls{UVvis} peak, related to a polythio\-phene $\pi-\pi^*$ transition. With additional 10~\% of methanol resulted in a further red shift of \SI{10}{\nm} (see data in Table~\ref{tab:sz17-cd}). We don't observe the usual quenching of the peak at \SI{260}{\nm} because in that region there's also the main absorption peak of pyridine. 

\begin{table}%sz17-cd
\centering
\caption[UV-vis and CD data of \cmpd+{sz17} solutions.]{\gls{UVvis} and \gls{CD} data of \cmpd+{sz17} solutions. Concentration \SI{137}{\mg\per\liter}. Pathlength \SI{1}{\mm}.}\label{tab:sz17-cd}
\begin{tabular}{c|c|c|c}
\toprule
processing&\tallcell{UV-vis\\ peak\\(\SI{}{nm})}&\tallcell{CD\\ inflection\\(\SI{}{nm})}&\tallcell{CD\\ amplitude\\(\SI{}{mdeg})}\\ \cmidrule{1-4}
\ch{CHCl3} solution 				&440&&\\
\ch{CHCl3} 80 \% - \ch{MeOH} 20 \%	&438&&\\
\ch{CHCl3} 70 \% - \ch{MeOH} 30 \%	&440&&\\
\ch{CHCl3} 70 \% - \ch{MeOH} 30 \% \SI{1}{\day}&440&$\approx448$&0.36\\
\ch{CHCl3} 60 \% - \ch{MeOH} 40 \%	&452&464&1.4\\
\ch{CHCl3} 50 \% - \ch{MeOH} 50 \%	&462&461&3.2\\
\ch{CHCl3} 30 \% - \ch{MeOH} 70 \%	&464&460&4.7\\

\bottomrule
\end{tabular}
\end{table}

\begin{figure}[tbp]%sz17-uvvis
\begin{tikzpicture}
\begin{axis}[axis x line=bottom,axis y line=left,enlarge x limits=false,enlarge y limits=false,width=1\textwidth,height=8cm,xlabel=Wavelength (nm),ylabel=Absorbance,xmin=240,xmax=650,legend style={at={(0.05,1.1)},anchor=north west,font=\footnotesize},cycle list name=linestyles*]
\addplot table {img/results/sz17-uvvis-137-chcl3-70-meoh-30.txt};
\addlegendentry{in \ch{CHCl3} 70 \% - \ch{MeOH} 30 \%};
\addplot table {img/results/sz17-uvvis-137-chcl3-60-meoh-40.txt};
\addlegendentry{in \ch{CHCl3} 60 \% - \ch{MeOH} 40 \%};
\addplot table {img/results/sz17-uvvis-137-chcl3-50-meoh-50.txt};
\addlegendentry{in \ch{CHCl3} 50 \% - \ch{MeOH} 50 \%};
\addplot table {img/results/sz17-uvvis-137-chcl3-30-meoh-70.txt};
\addlegendentry{in \ch{CHCl3} 30 \% - \ch{MeOH} 70 \%};
\end{axis}
\end{tikzpicture} \caption[UV-vis spectra of block copolymer \cmpd+{sz17} in mixtures of good-poor solvent.]{UV-vis spectra of polymer \cmpd+{sz17} in chloroform-methanol mixtures. Concentration \SI{137}{\mg\per\liter}. Pathlength \SI{1}{\mm}.}
\label{fig:sz17-uvvis}
\end{figure}

Comparing \gls{CD} spectra in chloroform-methanol solutions of \cmpd+{sz17} with spectra from \cmpd+{ig2-8} we can observe some differences. With \cmpd+{ig2-8}, up to 30~\% of methanol no \gls{CD} signal can be observed but waiting one day, \cmpd+{ig2-8} in 30~\% of methanol showed a very intense bisignate couplet (see Figure~\ref{fig:ig2-8-cd-meoh-1}) while \cmpd+{sz17} showed only an extremely weak signal in the same conditions. 
Adding more methanol resulted in an increase of the \gls{CD} intensity. Comparing intensities of \gls{CD} spectra of \cmpd+{sz17} and \cmpd+{ig2-8} in similar concentrations of methanol, \cmpd+{sz17} showed an optical activity around 6 times weaker. 
These differences should be related to the \gls{P4VP} block that is soluble both in chloroform and in methanol and could help the polythio\-phene block to remain in solution or could interfere with its supramolecular organization. We hypothesize the formation of direct micelles. 

\begin{figure}[tbp]%sz17-cd
\begin{tikzpicture}
\begin{axis}[axis x line=bottom,axis y line=left,enlarge x limits=false,enlarge y limits=true,width=1\textwidth,height=8cm,xlabel=Wavelength (nm),ylabel=Circular Dichroism (mdeg),xmin=240,xmax=600,legend pos=north west,cycle list name=linestyles*]
\addplot table {img/results/sz17-cd-137-chcl3-70-meoh-30-1day.txt};
\addlegendentry{in \ch{CHCl3} 70 \% - \ch{MeOH} 30 \% \SI{1}{\day}};
\addplot table {img/results/sz17-cd-137-chcl3-60-meoh-40.txt};
\addlegendentry{in \ch{CHCl3} 60 \% - \ch{MeOH} 40 \%};
\addplot table {img/results/sz17-cd-137-chcl3-50-meoh-50.txt};
\addlegendentry{in \ch{CHCl3} 50 \% - \ch{MeOH} 50 \%};
\addplot table {img/results/sz17-cd-137-chcl3-30-meoh-70.txt};
\addlegendentry{in \ch{CHCl3} 30 \% - \ch{MeOH} 70 \%};
\draw[help lines] (axis cs:240,0) -- (axis cs:600,0);
\end{axis}
\end{tikzpicture} \caption{Circular dichroism spectra of block copolymer \cmpd+{sz17} in chloroform-methanol solvent mixtures. Concentration \SI{137}{\mg\per\liter}. Pathlength \SI{1}{\mm}.}
\label{fig:sz17-cd}
\end{figure}

\end{subsection}
\end{section} 


\clearpage
\begin{section}{Further Studies}

Further studies will focus on the origin of the optical activity investigating the conformation of small aggregates and bulk solid via \gls{NMR} and probably solid state \gls{NMR}.  
Then we plan to study \cmpd+{sz17} to further confirm that it is a block copolymer rather than a blend of polymers. Then we could use \gls{CD} spectroscopy to study the formation of inverse micelles, this study can give useful information on the intra\-chain or inter\-chain origin of optical activity. 
Obviously we plan to test our diblock copolymer in solar cells active layer and to study improvements of \gls{PCBM} and nano\-crystals dispersion before and after an annealing process. We'll try to find a suitable matrix for the \gls{MALDI} analysis of both polythiophene and polyvinylpyridine, needed for the \gls{MALDI} characterization of the block copolymer. 

\end{section}

\chapter{Conclusions}

We synthesized three chiral polythiophenes using two different enantiomeric excesses and two different polymerization procedures, exploiting the living chain-growth polymerization of polythiophenes. 
Our polythiophenes are unreported in the literature.

From {\HNMR}, {\CNMR} and \gls{NMR2D} we verified a high regioregularity (see Section~\ref{sec:regioregolarita}). We observed (via \gls{SEC} analysis) that the addition of \ch{LiCl} to the polymerization resulted in a higher polymerization degree but increased also the polydispersity giving rise to a small amount of dimerization of the growing chain (see Section~\ref{sec:weight}). The higher weight was confirmed by \gls{UVvis} and \gls{FTIR} spectroscopies and by \gls{SEC}\-/\gls{MALDI}. Furthermore from \gls{MALDI} we saw that our polymers have various terminating groups, conversely to the results reported in the literature for well-defined polythiophenes (see Section~\ref{sec:sintesi-terminazione}). 

The chiral center, incorporated in the side groups, allowed us to perform a chiroptical study on the polymers in solution, in thin film and in solid powder state. The aggregation behavior in good-poor solvent mixtures and the supramolecular structure of these polymers were studied by means of \gls{UVvis} absorption, emission (\gls{PL}) and \acrfull{CD} spectroscopy. From \gls{PL} in good-poor solvents we can clearly notice differences between enantiopure and enantioenriched polymers (see Section~\ref{sec:pl}). 
The \gls{CD} results indicate the presence of more than one aggregation morphology, giving rise to similar \gls{CD} signals. 
Measuring the \gls{CD} intensity at various polymer concentrations we noticed a small nonlinearity which could be related to a multi\-chain event (see page~\pageref{nonlinearita}). 
Comparing the circular dichroism from enantiopure and from enantioenriched polymers we have verified that majority-rules effect is not present or it's too weak to be observed in our data (see page~\pageref{majority}). The solid state structure was investigated by means of \gls{XRD} and \acrshort{XRD2D} showing the presence of both amorphous domains and crystallized domains with the characteristic distances of an interdigitated regioregular poly\-(3-alkyl\-thio\-phene) (see Section~\ref{xrd}). No long range repeating structure was observed using \acrshort{SAXS}. Hence we tend to exclude the presence of large crystalline domains with a small tilt between lamellae giving rise to circular dichroism. More likely there is some extent of chiral structuring in amorphous regions. 

The longest enantiopure polymer was employed in the active layer of a solar cell as donor material in blend with \gls{PCBM}. Preliminary results show an open-circuit voltage $V_{oc}$ higher than for the standard \gls{P3HT}:\-\gls{PCBM} blend (see Section~\ref{cella}). This is caused by the electron withdrawing effect of the oxygen in the side chain of our polymer. 

Starting from the halogenated chain terminations we followed a novel convenient strategy to functionalize our higher weight enantiopure polythiophene \cmpd+{ig2-8} with \gls{BrPhEtTIPNO}, a \acrfull{NMRP} initiator and mediator. Then a short \acrfull{P4VP} block was polymerized obtaining an amphiphilic diblock copolymer (see Section~\ref{sec:nmrp}). {\HNMR} characterization of properly purified material confirms the formation of a block copolymer and allows us to determine the amount of inserted 4-vinyl\-pyridine. 
In spite of the smallness of the \gls{P4VP}, an effect can be observed in the preliminary \gls{CD} studies on aggregation from good-poor solvent mixtures. The coordinating property of \gls{P4VP} block with respect to \gls{PCBM} or semi\-conducting nano\-crystals can allow one to obtain a bi\-continuous pathway in a stable nano\-segregated active layer of a polymeric solar cell. Thus the present approach can unlock the possibility to characterize by circular dichroism spectroscopy the active layer of a stable bulk heterojunction solar cell. 
To the best of our knowledge we present here the first ``chiral thiophene-non thiophene'' copolymer. 

\chapter{Experimental Section}

\begin{section}{Materials and Methods}

All chemicals were purchased from Sigma-Aldrich and used without further purification, except: tetrahydrofuran and diethyl ether used in the syntheses were distilled on potassium and benzophenone under nitrogen atmosphere; B-(5'-formyl[2,2'-bi\-thio\-phen]\-5-yl)\-boronic acid was purchased from Tokyo Chemical Industry (\textsmaller{TCI}) and purified on a \SI{3}{\cm} silica gel column washing with dichloromethane and recovering with ethyl acetate; scalemic 2-methyl-1-butanol was purchased from Merck-Schuchardt and distilled under reduced pressure prior use.

All manipulations involving air-sensitive reagents or dry solvents were performed under an atmosphere of dry, filtered nitrogen. 

Thin layer chromatographies were performed using glass supported silica gel.

For flash chromatographies silica gel Merck 60 (0.04-\SI{0.06}{\mm}, flash) was used.

De\-oxy\-gen\-ated sodium carbonate aqueous solution was prepared by bubbling nitrogen in distilled water and after that sodium carbonate was added to obtain a \SI{1}{\Molar} solution.

Degassed ortho-di\-chloro\-benzene was bubbled with nitrogen and passed through three freeze-pump-thaw cycles.

Glassware used for reactions in nitrogen atmosphere were passed through three cycles of vacuum by mechanical pump, heating with heating gun and filling with nitrogen.

Common impurities found in tetra\-chloro\-ethane-d$_2$ used in \gls{NMR} analysis: {\HNMR} $\delta$ (ppm) 0.03 (s), 1.20 (s), 2.10 (s), 3.64 (dd), 5.14 (q); {\CNMR} $\delta$ (ppm) 29.62, 73.80 (t, solvent peak), 74.06 (s, solvent peak), 120.20. 

\ch{KCl} disks for \gls{UVvis} and \gls{CD} analysis were prepared grinding the solid in a mortar and consolidated with a press. 

Samples for photoluminescence measurements were prepared diluting with adequate amounts of methanol and chloroform a concentrated solution of the desired polymer in chloroform. After \SI{20}{\minute} samples were analyzed using in a glass cuvette. Some samples were measured again after further \SI{40}{\minute}.

\end{section}
\begin{section}{Instrumentation}

\Acrfull{FTIR} spectra were recorded using a \emph{Bruker \gls{FTIR} Tensor 27} spectrometer, samples were prepared as thin films cast on \ch{KBr} disk. Abbreviations used in reported spectra are: symm.\ (symmetrical), antisymm.\ (antisymmetrical), deform.\ (deformation), degen.\ (degenerate), arom.\ (aromatic).

Ultraviolet--visible (\gls{UVvis} ) spectra were recorded using using \emph{Jasco V-650} spectrophotometer and \SI{1}{\mm} path length quartz cuvette.

\label{fitting}
\Acrfull{CD} spectra were recorded using \emph{Jasco J-710} and a spectral window from \SI{230}{\nm} to \SI{600}{\nm}. Each measurement is repeated at least four times and a mean is reported. For measurements on films or solid a mean between various orientations of the substrate was reported, this helped to reduce artifacts from linear dichroism. Measurements on solutions were recorded using a \SI{1}{\mm} path length quartz cuvette. For estimating bisignate \gls{CD} peaks' center and amplitude the peaks were fitted with a sine waveform ($\mathrm{CD}= A+Bsin(C(\mathrm{wavelength}-D))$) using the Gnuplot software and considering only data points between 400 and \SI{500}{\nm}.

Spectra from \acrfull{GCMS} were recorded using GC \emph{Agilent 7890A} combined with detector 
\emph{Agilent 5975C MSD} and electron ionization source.

\Acrfull{DSC} analysis were accomplished using a \emph{Perkin Elmer Differential Scanning Calorimeter Pyris 1} instrument in absence of oxy\-gen. Heating rate: \SI{20}{\celsius\per\minute}, cooling rate: \SI{10}{\celsius\per\minute}.

\Acrfull{PL} spectra were collected with a setup consisting of \emph{Osram XBO} \SI{450}{\W} xenon lamp followed by a \emph{Jobin Yvon Horiba Gemini 180} monochromator, used for excitation wavelength selection, and a \emph{Spex 270M} monochromator combined with CCD, employed as a detector. The spectra were corrected for the instrument response functions. 

\Acrfull{NMR} spectra were collected at room temperature in 1,2-di\-deutero-1,1,2,2-tetra\-chloro\-ethane (\gls{TCE}-d2) and deuterated chloroform (\ch{CDCl3}) at the \gls{NMR} facility in ISMAC-CNR, Milan. All chemical shifts are reported in the standard $\delta$ notation of parts per million (ppm) using the peak of residual proton and carbon signals of \ch{CDCl3} ($\delta _{^1\mathrm{H}} = 7.26$, $\delta_{^{13}\mathrm{C}}=77.16$) and \gls{TCE}-d2 ($\delta _{^{1}\mathrm{H}}=5.94$, $\delta _{^{13}\mathrm{C}}=73.80$) as internal reference.

\Acrfull{SEC} measurements in Milan were performed using a \emph{GPCV 2000} pump and polystyrene standards. For measurements in chloroform at \SI{35}{\celsius} two columns \emph{PLgel Mixed C} (\SI{5}{\um} particle size) \emph{Polymer Laboratories} were used with a flux of \SI{0.8}{\mL\per\minute} and an injection volume of \SI{218.5}{\uL}. Two detectors were used: a \gls{UVvis} photo\-diode array \emph{DRI996} and a differential refractive index detector \emph{DRI2414}. 
For measurements in ortho-di\-chloro\-benzene at \SI{145}{\celsius} three columns \emph{PLgel Olexis Polymer Laboratories} were used with a flux of \SI{0.8}{\mL\per\minute} and an injection volume of \SI{218.5}{\uL}. Two detectors were used: a differential viscometer and a differential refractive index detector \emph{DRI2414}.

In \gls{SEC} measurements in Catania as a first step, relative averages ($M_n$, $M_w$, $M_z$) and polydispersity index (PDI = $M_w/M_n$) of all polymers were determined by a conventional size exclusion chromatography system calibrated toward polystyrene standards, using tetrahydrofuran as mobile phase, \SI{0.8}{\mL\per\minute} of flow rate and a temperature of \SI{35}{\celsius}. The conventional chromatographic system consisted of an integrated \emph{Alliance 2695} chromatographic system (de\-gasser, pump and injector) from \emph{Waters} (Milford, MA, USA) and a 2414 differential refractometer (DRI) used as concentration detector. The column set was composed of two PLgel Mixed C columns (\SI{5}{\um} of particle size) from \emph{Polymer Laboratories} (Shropshire, UK).

\Acrfull{MALDI} were recorded in the Institute of Chemistry and Technology of Polymers (\textsmaller{ICTP-CNR}, Catania) in reflector or linear delayed extraction mode, using a \emph{Voyager-DE STR} instrument (\emph{Perseptive Biosystem}) mass spectrometer, equipped with a nitrogen laser ($\lambda =$ \SI{337}{\nm}, pulse width = \SI{3}{\ns}), working in a positive ion mode. The accelerating voltage was \SI{20}{\kV}, grid voltage and delay time (delayed extraction, time lag), were optimized for each sample to achieve the higher mass resolution ($M / \Delta M$). Laser irradiance was maintained slightly above threshold. $\alpha$-Cyano-4-hydr\-oxy\-cinnamic acid and 1,8-di\-hydr\-oxy-9,10-di\-hydro\-anthracen-9-one (di\-thranol) \SI{0.1}{\Molar} in \gls{THF} were used as matrices. Then \SI{1}{\mL} of each analyte/matrix mixture was spotted on the \textsmaller{MALDI} sample holder and slowly dried to allow analyte/matrix co-crystallization. 

In \gls{SEC}\-/\gls{MALDI} combined analysis \SI{1}{\mL} fractions of each selected concentrated \gls{SEC} solutions was mixed with \SI{1}{\mL} or \SI{3}{\mL} of the matrix solution (trans-2-[3-(4-tert-butyl\-phenyl)-2-methyl-2-propenylidene]\-malononitrile, \textsmaller{DCTB}). Then, the same procedure used for the \gls{MALDI} analyses of the poly\-disperse polythiophene sample was followed. 

The \acrlong{WAXD} (\gls{WAXD} or \gls{XRD}) data were obtained at \SI{25}{\celsius} using a \emph{Siemens D-500} diffracto\-meter equipped with a \emph{Siemens FK 60-10} \SI{2000}{\W} tube (Cu K$_a$ radiation, $\lambda =$ \SI{0.154}{\nm}). 
The operating voltage and current were \SI{40}{\kV} and \SI{40}{\mA}, respectively. The data were collected from $2\theta =$ \SI{3}{\degree} to \SI{30}{\degree} at \SI{0.05}{\degree} intervals. 
\Acrfull{SAXS} measurements were conducted at \SI{25}{\celsius} with a \emph{Kratky Compact Camera}. Mono\-chromatized \ch{Cu} K$_a$ radiation 
($\lambda =$ \SI{0.154}{\nm}) was supplied by a stabilized \emph{Siemens Krystalloflex 710} generator and a \emph{Siemens FK 60-04}, \SI{1500}{\W} \ch{Cu} target tube operated at \SI{40}{\kV} and \SI{30}{\mA}.

\Acrfull{XRD2D} spectrum was recorded at \emph{Electra - Sincrotrone Trieste} synchrotron radiation facility in Trieste.

Spin coating were performed using \emph{Electron mec PRS 5V} spinner.

Annealing was performed using \emph{Mettler FP5} melting point apparatus.

Tests in solar cells were performed in Centre for Hybrid and Organic Solar Energy (\textsmaller{CHOSE}) of the University of Rome Tor Vergata. The photo\-voltaic devices were realized on \textsmaller{FTO}-coated soda-lime glass substrate (\emph{Pilkington}, $\approx8~\Omega/\square$) patterned by wet etching. Glass/\textsmaller{FTO} substrates were cleaned by ultrasonic bath with acetone and isopropyl alcohol (\SI{15}{\minute} each step). \textsmaller{PEDOT}:\textsmaller{PSS} (\emph{VPAI4083, Clevios}) was spin coated in a glove box at \SI{5000}{\rpm} for \SI{60}{\s} and thermally treated at \SI{150}{\celsius} for \SI{10}{\minute}. 
The active layer was spin coated at \SI{400}{\rpm} for \SI{120}{\s} and then dried by slow evaporation in nitrogen atmosphere. Calcium (\ch{Ca}) was evaporated on active layer in high vacuum (\SI{1E-6}{\milli\bar}) at \SI{0.06}{\nm\per\s}. Finally, a aluminum layer of \SI{100}{\nm} was thermally evaporated as electrode in high vacuum at \SI{0.1}{\nm\per\s}. A shadow mask was used in order to define the active area of the solar cells. 
The thickness of the as-deposited films was measured with a profilometer (\emph{Dektak 150}); transmittances of films were acquired by a \textsmaller{UV}-vis-\textsmaller{NIR} spectrophotometer (\emph{Shimadzu UV2550}). 
Electrical performance of the devices was evaluated outside the glove box under a solar simulator \emph{KHS Solar Constant 1200 AM1.5 Class B} (\SI{100}{\mW\per\square\cm}) with a parameter analyzer (\emph{Agilent E5291A}); for the measurements, we used a shadow mask to avoid performance contributions outside the active area. 

The utilized chem\-informatics resources were: OpenEnventory for laboratory report and chemicals inventory, OpenEnventory 2013-05-21 (\url{www.open-enventory.de}); Gaussian for computations, Gaussian 09, Revision C.01 (\url{www.gaussian.com}); OpenChrom for analyzing \gls{GCMS} data, OpenChrom 0.8.0 (\url{www.openchrom.net}); MarvinSketch for drawing chemical structures and reactions, MarvinSketch 6.0.2\_b97, 2013, ChemAxon (\url{www.chemaxon.com}); Mendeley for literature management, Mendeley Desktop 1.10.1 (\url{www.mendeley.com}); {\LaTeX} for writing this thesis, TeX Live 2013 (\url{www.tug.org/texlive}); Debian Linux the universal operating system, Debian Linux Jessie/Sid (\url{www.debian.org}); laptop for writing and computations, Acer Extensa 5220, Intel\textsuperscript{\textregistered} Core\texttrademark2 Duo CPU T7700 \SI{2.40}{\GHz}, RAM memory \SI{3}{\gibi\byte}. All the images (except the logo of Pisa university) were made on purpose for this thesis.

\end{section}
\clearpage
\begin{section}{Synthesis of Monomers}

\subsection[2-bromo-3-[(2-methyl\-butyl\-oxy)\-methyl{]}\-thio\-phene (\cmpd+{ig2-9})]{2-bromo-3-[(2-methyl\-butyl\-oxy)\-methyl{]}\-thio\-phene (\cmpd+{ig2-9}) \\ Etherification of 2-bromo-3-(bromo\-methyl)\-thio\-phene with 2-methyl-1-butanol} 
\label{sec:ig2-9}

\begin{figure}[H]%syn1-eterificazione
\centering
\includegraphics[scale=0.5]
{syn1-eterificazione.pdf}
\end{figure}

2-methyl-1-butanol (\SI{0.26}{\mL}, \SI{210}{\mg}, \SI{2.4}{\mmol}) distilled on calcium oxide was dissolved in dry tetrahydrofuran (\SI{6}{\mL}) in a Schlenk tube under nitrogen atmosphere. The reaction was cooled to \SI{0}{\celsius} and sodium hydride (\SI{55}{\mg}, \SI{2.3}{\mmol}) was added in portions. After stirring at room temperature for \SI{30}{\minute} the reaction was cooled to \SI{0}{\celsius} and 2-bromo-3-(bromo\-methyl)\-thio\-phene (\SI{0.26}{\mL}, \SI{510}{\mg}, \SI{2}{\mmol}) was added. The reaction mixture was stirred for \SI{40}{\minute} at \SI{0}{\celsius}, then at room temperature overnight. 
The reaction was quenched by adding water and extracted with dichloromethane (3 $\times$ \SI{8}{\mL}). Organic fractions were washed with water and purified through flash column chromatography (silica gel, heptane:\-di\-chloro\-methane 4:1, dry deposition) obtaining 2-bromo-3-[(2-methyl\-butyl\-oxy)\-methyl{]}\-thio\-phene \cmpd+{ig2-9} as a colorless oil (\SI{348.7}{\mg}, yield 66~\%, \gls{GCMS} pure). 

\gls{GCMS} (EI): 262, 264 (\ch{M+}), 174.9, 176.9 (\ch{M+ $-$ OC5H11}), 113.0 (\ch{M+ $-$ C5H11 $-$ Br + H}), 97.0 (\ch{M+ $-$ OC5H11 $-$ Br + H}).

\subsection[2-bromo-3-[((S)-2-methyl\-butyl\-oxy)\-methyl{]}\-thio\-phene (\cmpd+{ig2-1})]{2-bromo-3-[((S)-2-methyl\-butyl\-oxy)\-methyl{]}\-thio\-phene (\cmpd+{ig2-1}) \\ Etherification of 2-bromo-3-(bromo\-methyl)\-thio\-phene with (S)-2-methyl-1-butanol}%SDIG2-1
\label{sec:ig2-1}
The synthetic procedure for 2-bromo-3-[(2-methyl\-butyl\-oxy)\-methyl{]}\-thio\-phene \cmpd+{ig2-9} was followed using 2-bromo-3-(bromo\-methyl)\-thio\-phene (\SI{0.67}{\mL}, \SI{1.3}{\g}, \SI{5.2}{\mmol}), (S)-2-methyl-1-butanol (without further purification, \SI{0.704}{\mL}, \SI{0.577}{\g}, \SI{6.54}{\mmol}), sodium hydride (\SI{0.140}{\g}, \SI{5.82}{\mmol}) and tetrahydrofuran (\SI{15}{\mL}). The crude product was extracted with dichloromethane and water, then was purified by silica gel column chromatography (hexane:\-chloro\-form 4:1) affording 2-bromo-3-[((S)-2-methyl\-butyl\-oxy)\-methyl{]}\-thio\-phene \cmpd+{ig2-1} as a colorless oil (\SI{0.959}{\g}, yield 71~\%, \gls{GCMS} pure).

\gls{GCMS} (EI): 262, 264 (\ch{M+}), 174.9, 176.9 (\ch{M+ $-$ OC5H11}), 113.0 (\ch{M+ $-$ C5H11 $-$ Br + H}), 97.0 (\ch{M+ $-$ OC5H11 $-$ Br + H}).

\subsection[2-bromo-5-iodo-3-[(2-methyl\-butyl\-oxy)\-methyl{]}\-thio\-phene (\cmpd+{ig2-10})]{2-bromo-5-iodo-3-[(2-methyl\-butyl\-oxy)\-methyl{]}\-thio\-phene (\cmpd+{ig2-10}) \\ Iodination of 2-bromo-3-[(2-methyl\-butyl\-oxy)\-methyl{]}\-thio\-phene}
\label{sec:ig2-10}

\begin{figure}[H]%syn2-iodurazione
\centering
\includegraphics[scale=0.5]
{syn2-iodurazione.pdf}
\end{figure}

In a Schlenk tube 2-bromo-3-[(2-methyl\-butyl\-oxy)\-methyl{]}\-thio\-phene \cmpd+{ig2-1} (\SI{349}{\mg}, \SI{1.32}{\mmol}) was dried with mechanical pump and kept under nitrogen atmosphere. Dichloromethane (\SI{3}{\mL}, degassed by bubbling nitrogen) was added and the flask was cooled to \SI{0}{\celsius}. Iodine (\SI{219}{\mg}, \SI{0.864}{\mmol}) and iodo\-benzene di\-acetate (\ch{PhI(OAc)2}, \SI{255}{\mg}, \SI{0.79}{\mmol}) were added and the reaction was stirred at \SI{0}{\celsius} for \SI{4}{\hour}. 
The reaction was quenched with saturated aqueous sodium thio\-sulfate (\SI{5}{\mL}) keeping the temperature at \SI{0}{\celsius}. Then was extracted with dichloromethane (4 $\times$ \SI{10}{\mL}) and the organic fractions were washed with water. The product was purified by a flash chromatography on silica gel (hexane:\-chloro\-form 9:1, dry deposition). 
The selected fractions were concentrated and purified through a second flash column chromatography (hexane:\-ethyl\-acetate 69:1, dry deposition) to obtain a pale yellow oil (\SI{484.6}{\mg}). \gls{GCMS} analysis on the obtained oil showed the presence of the desired product 2-bromo-5-iodo-3-[(2-methyl\-butyl\-oxy)\-methyl{]}\-thio\-phene \cmpd+{ig2-10} (94~\% in oil, 93~\% overall yield) together with an impurity of 2-bromo-4,5-diiodo-3-[(2-methyl\-butyl\-oxy)\-methyl{]}\-thio\-phene (6~\% in oil).

{\HNMR} (\SI{600}{\MHz}, \gls{TCE}-d$_2$) $\delta$ (ppm) 7.13 (s, 1H, arom.\ H), 4.31 (s, 2H, \ch{CCH2O}), 3.26 (dd, $J =$ \SI{9.1}{\Hz} (geminal), \SI{6.3}{\Hz}, 1H, \ch{OC\underline{H}2CH}), 3.17 (dd, $J =$ \SI{8.9}{\Hz} (geminal), \SI{7.0}{\Hz}, 1H, \ch{OC\underline{H}2CH}), 1.61 (m, 1H, \ch{OCH2C\underline{H}}), 1.40 (m, 1H, \ch{OCH2CHC\underline{H}2}), 1.09 (m, 1H, \ch{OCH2CHC\underline{H}2}), 0.85 (d, $J =$ \SI{6.6}{\Hz}, 3H, \ch{CHC\underline{H}3}), 0.84 (t, $J =$ \SI{7.4}{\Hz}, 3H, \ch{CH2C\underline{H}3}).
{\CNMR} (\SI{150}{\MHz}, \gls{TCE}-d$_2$) $\delta$ (ppm) 141.01 (arom.\ \ch{\underline{C}CH2}), 137.81 (arom.\ \ch{CH}), 113.05 (\ch{CBr}), 75.80 (\ch{O\underline{C}H2CH}), 71.86 (\ch{CI}), 66.36 (\ch{C\underline{C}H2O}), 34.69 (\ch{OCH2\underline{C}H}), 26.08 (\ch{OCH2CH\underline{C}H2}), 16.52 (\ch{CH\underline{C}H3}), 11.30 (\ch{CH2\underline{C}H3}).
\gls{GCMS} (EI): 387.9, 389.9 (\ch{M+}), 300.8, 302.8 (\ch{M+ $-$ OC5H11}), 222.9 (\ch{M+ $-$ C5H11 $-$ Br + H}), 174.9, 176.9 (\ch{M+ $-$ C5H11 $-$ I + H}).
\gls{FTIR} (film on \ch{KBr}) $\nu$ (\SI{}{\per\cm}) 3093 (weak, arom.\ CH stretching), 2959 (strong, \ch{CH3} degen.\ stretching), 2928 (strong, \ch{CH2} antisymm.\ stretching), 2873 (strong, \ch{CH3} symm.\ stretching), 2858 (strong, \ch{CH2} symm.\ stretching), 1461 (medium, \ch{C=C} thiophene symm.\ stretching, \ch{CH3} degen.\ deform.), 1413 (medium), 1178 (medium), 1098 (strong, \ch{CO} antisymm.\ stretching), 994 (medium, ring breathing), 833 (medium, arom.\ H out-of-plane bending), 471.

\subsection[2-bromo-5-iodo-3-[((S)-2-methyl\-butyl\-oxy)\-methyl{]}\-thio\-phene (\cmpd+{ig2-2})]{2-bromo-5-iodo-3-[((S)-2-methyl\-butyl\-oxy)\-methyl{]}\-thio\-phene (\cmpd+{ig2-2}) \\ Iodination of 2-bromo-3-[((S)-2-methyl\-butyl\-oxy)\-methyl{]}\-thio\-phene}
\label{sec:ig2-2}
The synthetic procedure for 2-bromo-5-iodo-3-[(2-methyl\-butyl\-oxy)\-methyl{]}\-thio\-phene \cmpd+{ig2-10} was followed using 2-bromo-3-[((S)-2-methyl\-butyl\-oxy)\-methyl{]}\-thio\-phene \cmpd+{ig2-1} (\SI{0.959}{\g}, \SI{3.64}{\mmol}), iodine (\ch{I2}, \SI{0.601}{\g}, \SI{2.37}{\mmol}), iodo\-benzene di\-acetate (\ch{PhI(OAc)2}, \SI{0.700}{\g}, \SI{2.17}{\mmol}) and dichloromethane (\SI{8}{\mL}). 
The crude product was extracted with dichloromethane and water, then was purified by two following silica gel column chromatographies (hexane:\-chloro\-form 9:1, then hexane:\-ethyl acetate 69:1) affording 2-bromo-5-iodo-3-[((S)-2-methyl\-butyl\-oxy)\-methyl{]}\-thio\-phene \cmpd+{ig2-2} as a pale yellow oil (\SI{1.343}{\g}, yield 95~\%, \gls{GCMS} pure, {\HNMR} pure).

{\HNMR} (\SI{600}{\MHz}, \gls{TCE}-d$_2$) $\delta$ (ppm) 7.12 (s, 1H, arom.\ H), 4.31 (s, 2H, \ch{CCH2O}), 3.25 (dd, $J =$ \SI{9.1}{\Hz} (geminal), \SI{6.2}{\Hz}, 1H, \ch{OC\underline{H}2CH}), 3.17 (dd, $J =$ \SI{9.1}{\Hz} (geminal), \SI{6.9}{\Hz}, 1H, \ch{OC\underline{H}2CH}), 1.61 (m, 1H, \ch{OCH2C\underline{H}}), 1.39 (m, 1H, \ch{OCH2CHC\underline{H}2}), 1.08 (m, 1H, \ch{OCH2CHC\underline{H}2}), 0.85 (d, $J =$ \SI{6.7}{\Hz}, 3H, \ch{CHC\underline{H}3}), 0.84 (t, $J =$ \SI{7.4}{\Hz}, 3H, \ch{CH2C\underline{H}3}).
{\HNMR} (\SI{600}{\MHz}, \ch{CDCl3}) $\delta$ (ppm) 7.15 (s, 1H, arom.\ H), 4.37 (s, 2H, \ch{CCH2O}), 3.31 (dd, $J =$ \SI{9.0}{\Hz} (geminal), \SI{6.2}{\Hz}, 1H, \ch{OC\underline{H}2CH}), 3.22 (dd, $J =$ \SI{9.0}{\Hz} (geminal), \SI{6.7}{\Hz}, 1H, \ch{OC\underline{H}2CH}), 1.66 (m, 1H, \ch{OCH2C\underline{H}}), 1.45 (m, 1H, \ch{OCH2CHC\underline{H}2}), 1.14 (m, 1H, \ch{OCH2CHC\underline{H}2}), 0.91 (d, $J =$ \SI{6.7}{\Hz}, 3H, \ch{CHC\underline{H}3}), 0.89 (t, $J =$ \SI{7.5}{\Hz}, 3H, \ch{CH2C\underline{H}3}). 
{\CNMR} (\SI{150}{\MHz}, \ch{CDCl3}) $\delta$ (ppm) 141.10 (arom.\ \ch{\underline{C}CH2}), 137.97 (arom.\ \ch{CH}), 113.41 (\ch{CBr}), 76.05 (\ch{O\underline{C}H2CH}), 71.97 (\ch{CI}), 66.68 (\ch{C\underline{C}H2O}), 35.09 (\ch{OCH2\underline{C}H}), 26.34 (\ch{OCH2CH\underline{C}H2}), 16.73 (\ch{CH\underline{C}H3}), 11.43 (\ch{CH2\underline{C}H3}).
\gls{GCMS} (EI): 387.9, 389.9 (\ch{M+}), 300.8, 302.8 (\ch{M+ $-$ OC5H11}), 222.9 (\ch{M+ $-$ C5H11 $-$ Br + H}), 174.9, 176.9 (\ch{M+ $-$ C5H11 $-$ I + H}).

\end{section}
\begin{section}{Polymerization}

\subsection[Poly\{3-[((S)-2-methyl\-butyl\-oxy)\-methyl{]}\-thio\-phene\} (\cmpd+{ig2-4})]{Poly\{3-[((S)-2-methyl\-butyl\-oxy)\-methyl{]}\-thio\-phene\} (\cmpd+{ig2-4}) \\ Polymerization of 2-bromo-5-iodo-3-[((S)-2-methyl\-butyl\-oxy)\-methyl{]}\-thio\-phene}
\label{sec:ig2-4}

\begin{figure}[H]%syn3-attivazione-nolicl
\centering
\includegraphics[scale=0.5]
{syn3-attivazione-nolicl.pdf}
\smallskip

\includegraphics[scale=0.5]
{syn4-polimerizzazione-nolicl.pdf}

\includegraphics[scale=0.5]
{syn5-quenching.pdf}
\end{figure}
2-bromo-5-iodo-3-[((S)-2-methyl\-butyl\-oxy)\-methyl{]}\-thio\-phene \cmpd+{ig2-2} (\SI{266}{\mg}, \SI{0.684}{\mmol}) was placed in a Schlenk tube under nitrogen atmosphere. Then dry tetrahydrofuran (\SI{4}{\mL}) was added, the reaction mixture was cooled to \SI{0}{\celsius} and isopropyl-magnesium chloride (\SI{2.0}{\Molar} solution in tetrahydrofuran, \SI{0.34}{\mL}, \SI{0.68}{\mmol}) was added dropwise. The mixture was stirred at \SI{0}{\celsius} for \SI{1}{\hour}, then a suspension of \acrlong{nidppp} (\acrshort{nidppp}, \SI{6.0}{\mg}, \SI{0.011}{\mmol}%
) in dry tetrahydrofuran (\SI{3}{\mL}) was added in one portion, then was stirred at room temperature for \SI{24}{\hour}.
The reaction was quenched by adding aqueous hydrochloric acid (\SI{5}{\Molar}, \SI{11}{\mL}), then was extracted with chloroform (\SI{200}{\mL}). The organic fraction was reduced to few milliliters, precipitated from methanol (\SI{0.5}{\L}) and filtered (\gls{PTFE} membrane, \SI{0.45}{\um}).
The solid was washed with hot methanol (\SI{50}{\celsius}), dried under reduced pressure and recovered mechanically from the filter, redissolved in chloroform, re\-precipitated in methanol and filtered again to obtain poly\{3-[((S)-2-methyl\-butyl\-oxy)\-methyl{]}\-thio\-phene\} \cmpd+{ig2-4} as a dark red solid (\SI{35}{\mg}, yield 22~\%). 

{\HNMR} (\SI{600}{\MHz}, \gls{TCE}-d$_2$) $\delta$ (ppm) 7.22 (s, 1H, arom.\ H), 4.54 (s, 2H, \ch{CCH2O}), 3.40 (m, 1H, \ch{OC\underline{H}2CH}), 3.31 (m, 1H, \ch{OC\underline{H}2CH}), 1.69 (m, 1H, \ch{OCH2C\underline{H}}), 1.47 (m, 1H, \ch{OCH2CHC\underline{H}2}), 1.15 (m, 1H, \ch{OCH2CHC\underline{H}2}), 0.94 -- 0.86 (m, 6H, \ch{CHC\underline{H}3 + CH2C\underline{H}3}).
{\HNMR} (\SI{600}{\MHz}, \ch{CDCl3}) $\delta$ (ppm) 7.25 (s, 1H, arom.\ H), 4.58 (s, 2H, \ch{CCH2O}), 3.45 -- 3.40 (m, 1H, \ch{OC\underline{H}2CH}), 3.36 -- 3.31 (m, 1H, \ch{OC\underline{H}2CH}), 1.85 -- 1.69 (m, 1H, \ch{OCH2C\underline{H}}), 1.57 -- 1.49 (m, 1H, \ch{OCH2CHC\underline{H}2}), 1.24 -- 1.15 (m, 1H, \ch{OCH2CHC\underline{H}2}), 0.99 -- 0.91 (m, 6H, \ch{CHC\underline{H}3 + CH2C\underline{H}3}).
{\CNMR} (\SI{150}{\MHz}, \gls{TCE}-d$_2$) $\delta$ (ppm) 136.54 (arom.\ \ch{\underline{C}CH2}), 133.68 (arom.\ \ch{S\underline{C}CCH2}), 132.91 (arom.\ \ch{S\underline{C}CH}), 129.28 (arom.\ \ch{CH}), 75.88 (\ch{O\underline{C}H2CH}), 66.71 (\ch{C\underline{C}H2O}), 34.86 (\ch{OCH2\underline{C}H}), 26.25 (\ch{OCH2CH\underline{C}H2}), 16.73 (\ch{CH\underline{C}H3}), 11.38 (\ch{CH2\underline{C}H3}).
{\CNMR} (\SI{100}{\MHz}, \ch{CDCl3}) $\delta$ (ppm) 136.61 (arom.\ \ch{\underline{C}CH2}), 134.25 (arom.\ \ch{S\underline{C}CCH2}), 133.36 (arom.\ \ch{S\underline{C}CH}), 129.52 (arom.\ \ch{CH}), 76.13 (\ch{O\underline{C}H2CH}), 67.02 (\ch{C\underline{C}H2O}), 35.23 (\ch{OCH2\underline{C}H}), 26.50 (\ch{OCH2CH\underline{C}H2}), 16.94 (\ch{CH\underline{C}H3}), 11.52 (\ch{CH2\underline{C}H3}).
$^1$H-COSY, $^1$H-TOCSY, $^1$H,$^{13}$C-HSQC and $^1$H,$^{13}$C-HMBC characterizations results are reported in table~\ref{tab:ig2-4-nmr2d}.
\gls{FTIR} (film on \ch{KBr}) $\nu$ (\SI{}{\per\cm}) 3064 (weak, arom.\ CH stretching), 2961 (strong, \ch{CH3} degen.\ stretching), 2930 (strong, \ch{CH2} antisymm.\ stretching), 2874 (strong, \ch{CH3} symm.\ stretching), 2858 (strong, \ch{CH2} symm.\ stretching), 2805 (medium), 2729 (weak), 1517 (medium), 1461 (medium, \ch{C=C} thiophene symm.\ stretching, \ch{CH3} degen.\ deform.), 1378 (medium, \ch{CH3} symm.\ deform.), 1180 (medium), 1091 (strong, \ch{CO} antisymm.\ stretching), 838 (medium, arom.\ H out-of-plane bending).

\vfill

\subsection[Poly\{3-[(2-methyl\-butyl\-oxy)\-methyl{]}\-thio\-phene\} (\cmpd+{ig2-15})]{Poly\{3-[(2-methyl\-butyl\-oxy)\-methyl{]}\-thio\-phene\} (\cmpd+{ig2-15}) \\ Polymerization of 2-bromo-5-iodo-3-[(2-methyl\-butyl\-oxy)\-methyl{]}\-thio\-phene in presence of \ch{LiCl}}
\label{sec:ig2-15}

\begin{figure}[H]%syn3-attivazione-licl
\centering
\includegraphics[scale=0.5]
{syn3-attivazione-licl.pdf}

\smallskip
\includegraphics[scale=0.5]
{syn4-polimerizzazione-licl.pdf}

\includegraphics[scale=0.5]
{syn5-quenching.pdf}
\end{figure}
A Schlenk tube equipped with a septum and stirring bar was charged with lithium chloride (\SI{85}{\mg}, \SI{2}{\mmol}, hygroscopic) and was heated by a heat-gun under reduced pressure. After the tube was cooled to room temperature under a nitrogen atmosphere, 2-bromo-5-iodo-3-[(2-methyl\-butyl\-oxy)\-methyl{]}\-thio\-phene \cmpd+{ig2-10} (\SI{152}{\mg}, \SI{0.391}{\mmol}, optical purity 59~\%) was added and dried under reduced pressure (mechanical pump). 
Then dry tetrahydrofuran (\SI{6}{\mL}) was added, the reaction mixture was cooled to \SI{0}{\celsius} and iso\-propyl\-magnesium chloride (\SI{2.0}{\Molar} solution in tetrahydrofuran, \SI{0.20}{\mL}, \SI{0.40}{\mmol}) was added dropwise (\SI{2}{\minute}) while stirring. 
The mixture was stirred at \SI{0}{\celsius} for \SI{50}{\minute}, then a suspension of \acrlong{nidppp} (\acrshort{nidppp}, \SI{4.84}{\mg}, \SI{8.93}{\umol}) in dry tetrahydrofuran (\SI{1.29}{\mL}) was added in one portion. 
The reaction was stirred at \SI{0}{\celsius} for \SI{2}{\minute} and at room temperature overnight. The polymerization was quenched by the rapid addition of aqueous hydrochloric acid (\SI{5}{\Molar}, \SI{4}{\mL}) and extracted with chloroform (4 $\times$ \SI{2}{\mL}). The collected organic fractions were washed with water (2 $\times$ \SI{7}{\mL}) and the solvent removed under reduced pressure. 
The resulting red solid was fractionated by Soxhlet extraction using methanol (extracted \SI{26}{\mg}), heptane (extracted \SI{3}{\mg}) and acetone (extracted \SI{3}{\mg}), and finally recovered using chloroform. The solvent was removed under reduced pressure to give poly\{3-[(2-methyl\-butyl\-oxy)\-methyl{]}\-thio\-phene\} \cmpd+{ig2-15} as a dark red solid (\SI{43.3}{\mg}, yield 61~\%).

{\HNMR} (\SI{600}{\MHz}, \gls{TCE}-d$_2$) $\delta$ (ppm) 7.21 (s, 1H, arom.\ H), 4.55 (s, 2H, \ch{CCH2O}), 3.40 (dd, $J =$ \SI{8.8}{\Hz} (geminal), \SI{6.2}{\Hz}, 1H, \ch{OC\underline{H}2CH}), 3.31 (dd, $J =$ \SI{9.1}{\Hz} (geminal), \SI{6.8}{\Hz}, 1H, \ch{OC\underline{H}2CH}), 1.74 -- 1.63 (m, 1H, \ch{OCH2C\underline{H}}), 1.53 -- 1.41 (m, 1H, \ch{OCH2CHC\underline{H}2}), 1.18 -- 1.08 (m, 1H, \ch{OCH2CHC\underline{H}2}), 0.92 (d, $J =$ \SI{6.7}{\Hz}, 3H, \ch{CHC\underline{H}3}), 0.88 (t, $J =$ \SI{7.4}{\Hz}, 3H, \ch{CH2C\underline{H}3}).
{\CNMR} (\SI{150}{\MHz}, \gls{TCE}-d$_2$) $\delta$ (ppm) 136.55 (arom.\ \ch{\underline{C}CH2}), 133.68 (arom.\ \ch{S\underline{C}CCH2}), 132.90 (arom.\ \ch{S\underline{C}CH}), 129.28 (arom.\ \ch{CH}), 75.88 (\ch{O\underline{C}H2CH}), 66.71 (\ch{C\underline{C}H2O}), end of spectral window.
\gls{FTIR} (film on \ch{KBr}) $\nu$ (\SI{}{\per\cm}) 3061, 2961, 2927, 2874, 2858, 2805, 1462, 1378, 1261, 1091, 1019, 841, 801.

\subsection[Poly\{3-[((S)-2-methyl\-butyl\-oxy)\-methyl{]}\-thio\-phene\} (\cmpd+{ig2-8})]{Poly\{3-[((S)-2-methyl\-butyl\-oxy)\-methyl{]}\-thio\-phene\} (\cmpd+{ig2-8}) \\ Polymerization of 2-bromo-5-iodo-3-[((S)-2-methyl\-butyl\-oxy)\-methyl{]}\-thio\-phene in presence of \ch{LiCl}}
\label{sec:ig2-8}
The synthetic procedure for poly\{3-[(2-methyl\-butyl\-oxy)\-methyl{]}\-thio\-phene\} \cmpd+{ig2-15} was followed using 2-bromo-5-iodo-3-[((S)-2-methyl\-butyl\-oxy)\-methyl{]}\-thio\-phene \cmpd+{ig2-2} (\SI{312}{\mg}, \SI{0.802}{\mmol}%
), lithium chloride (\SI{176}{\mg}, \SI{4.15}{\mmol}), iso\-propyl\-magnesium chloride (\SI{0.400}{\mL}, \SI{0.800}{\mmol}%
), \acrlong{nidppp} (\acrshort{nidppp}, \SI{9.8}{\mg}, \SI{0.018}{\mmol}%
), and tetrahydrofuran (\SI{20}{\mL}). 
The crude product was extracted with chloroform and water, the collected organic fractions were dried under reduced pressure and fractionated by Soxhlet extraction using methanol (extracted \SI{52.5}{\mg}), hexane (extracted \SI{9}{\mg}) and acetone (extracted \SI{4}{\mg}), and finally recovered using chloroform. After solvent evaporation poly\{3-[((S)-2-methyl\-butyl\-oxy)\-methyl{]}\-thio\-phene\} was obtained as a dark red solid \cmpd+{ig2-8} (\SI{113.6}{\mg}, yield 78~\%%
).

{\HNMR} (\SI{600}{\MHz}, \gls{TCE}-d$_2$) $\delta$ (ppm) 7.21 (s, 1H, arom.\ H), 4.54 (s, 2H, \ch{CCH2O}), 3.40 (dd, $J =$ \SI{9.0}{\Hz} (geminal), \SI{6.1}{\Hz}, 1H, \ch{OC\underline{H}2CH}), 3.31 (dd, $J =$ \SI{8.8}{\Hz} (geminal), \SI{7.0}{\Hz}, 1H, \ch{OC\underline{H}2CH}), 1.77 -- 1.59 (m, 1H, \ch{OCH2C\underline{H}}), 1.52 -- 1.36 (m, 1H, \ch{OCH2CHC\underline{H}2}), 1.19 -- 1.05 (m, 1H, \ch{OCH2CHC\underline{H}2}), 0.92 (d, $J =$ \SI{6.5}{\Hz}, 3H, \ch{CHC\underline{H}3}), 0.88 (t, $J =$ \SI{7.3}{\Hz}, 3H, \ch{CH2C\underline{H}3}).
{\CNMR} (\SI{150}{\MHz}, \gls{TCE}-d$_2$) $\delta$ (ppm) 136.55 (arom.\ \ch{\underline{C}CH2}), 133.68 (arom.\ \ch{S\underline{C}CCH2}), 132.90 (arom.\ \ch{S\underline{C}CH}), 129.28 (arom.\ \ch{CH}), 75.88 (\ch{O\underline{C}H2CH}), 66.71 (\ch{C\underline{C}H2O}), 34.86 (\ch{OCH2\underline{C}H}), 26.25 (\ch{OCH2CH\underline{C}H2}), 16.73 (\ch{CH\underline{C}H3}), 11.38 (\ch{CH2\underline{C}H3}).
\gls{FTIR} (film on \ch{KBr}) $\nu$ (\SI{}{\per\cm}) 3062, 2961, 2928, 2874, 2858, 2805, 1462, 1378, 1091, 841. 

\end{section}
\begin{section}{Synthesis of Mediator for Radical Polimerization}

\subsection[1-(4-Bromo\-phenyl)\-ethanol (\cmpd+{ig2-16})]{1-(4-Bromo\-phenyl)\-ethanol (\cmpd+{ig2-16}) \\ Reduction of 4-bromo\-aceto\-phenone}%SDIG2-16

\begin{figure}[H]%syn6-riduzione
\centering
\includegraphics[scale=0.5]
{syn6-riduzione.pdf}
\end{figure}

Sodium boro\-hydride (\ch{NaBH4}, \SI{57}{\mg}, \SI{1.5}{\mmol}) was added to a stirred solution of 4-bromo\-aceto\-phenone (\SI{203}{\mg}, \SI{1}{\mmol}) in ethanol (reagent grade, \SI{10}{\mL}). After stirring for \SI{4}{\hour} the reaction mixture was concentrated in mild vacuum and the residue dissolved in dichloromethane (\SI{30}{\mL}). The organic layer was subsequently washed with distilled water (3 $\times$ \SI{20}{\mL}), the organic extract was concentrated in mild vacuum to afford 1-(4-bromo\-phenyl)\-ethanol \cmpd+{ig2-16} (\SI{155}{\mg}, yield 76~\%, \gls{GCMS} pure). 

\gls{GCMS} (EI): 202.0, 204.0 (\ch{M+}), 184.9, 186.9 (\ch{M+ $-$ CH3}), 156.9, 158.9, 121.1 (\ch{M+ $-$ Br}), 103 (\ch{M+ $-$ Br $-$ H2O}).

\subsection[1-Bromo-4-(1-bromo\-ethyl)\-benzene (\cmpd+{ig2-17})]{1-Bromo-4-(1-bromo\-ethyl)\-benzene (\cmpd+{ig2-17}) \\ Bromination of 1-(4-bromo\-phenyl)\-ethanol}

\begin{figure}[H]%syn7-bromurazione
\centering
\includegraphics[scale=0.5]
{syn7-bromurazione.pdf}
\end{figure}

In a \SI{25}{\mL} Schlenk tube phosphorus tri\-bromide (\ch{PBr3}, \SI{60}{\uL}, \SI{0.64}{\mmol}) was dissolved in toluene dry (\SI{4}{\mL}, this solution showed some white precipitate at first, then turned to pale yellow). 
\SI{2}{\mL} of the previously prepared phosphorus tri\-bromide solution were added drop wise (\SI{10}{\minute}) to a solution of 1-(4-bromo\-phenyl)\-ethanol \cmpd+{ig2-16} (\SI{155}{\mg}, \SI{0.771}{\mmol} in a \SI{25}{\mL} Schlenk tube) in toluene dry (about \SI{2}{\mL}) at \SI{-5}{\celsius} (cooled using an immersion cooler in a external bath). The reaction was stirred for \SI{4}{\hour} keeping at \SI{-5}{\celsius}, then it was quenched with water.
The organic layer was separated and the aqueous layer was extracted with diethyl ether (2 $\times$ \SI{15}{\mL}). The organic layers were collected and washed with water (\SI{10}{\mL}), dried over anhydrous magnesium sulfate, filtered and concentrated under reduced pressure (\SI{40}{\celsius}, \SI{30}{\milli\bar}, then nitrogen flux) obtaining a pale pink oil (\SI{138}{\mg}). From \gls{GCMS} analysis the product results to be a mixture of 1-bromo-4-(1-bromo\-ethyl)\-benzene \cmpd+{ig2-17} (peak area 55~\%, \SI{76}{\mg}, \SI{0.29}{\mmol},
yield 37~\%) and 4-bromo\-styrene (peak area 45~\%, \SI{62}{\mg}, \SI{0.34}{\mmol}, 
yield 44~\%).

\gls{GCMS} (EI): 263.9 (\ch{M+}), 248.8 (\ch{M+ $-$ CH3}), 183.0, 185.0 (\ch{M+ $-$ Br}), 169, 171 (\ch{M+ $-$ Br $-$ CH2}), 155.9, 157.9 (\ch{M+ $-$ Br $-$ CCH3}), 104.1 (\ch{M+ $-$ Br $-$ Br}).

\subsection[\textit{N-Tert}-butyl-\textit{N}-(2-methyl-1-phenyl\-propyl)-\textit{O}-(1-phenyl\-ethyl)\-hydroxyl\-amine (\cmpd+{ig2-21})]{\textit{N-Tert}-butyl-\textit{N}-(2-methyl-1-phenyl\-propyl)-\textit{O}-(1-phenyl\-ethyl)\-hydroxyl\-amine (\cmpd+{ig2-21}) \\ Condensation of 1-bromo-4-(1-bromo\-ethyl)\-benzene with 2,2,5-tri\-methyl-4-phenyl-3-aza\-hexane-3-nitroxide}

\begin{figure}[H]%syn8-tipno-br
\centering
\includegraphics[scale=0.5]
{syn8-tipno-br.pdf}
\end{figure}

In a \SI{10}{\mL} Schlenk tube copper(I) bromide (\SI{17}{\mg}, \SI{0.12}{\mmol}), copper powder (\SI{547}{\mg}, \SI{8.61}{\mmol}) and 2,2'-bi\-pyridine (\SI{84}{\mg}, \SI{0.54}{\mmol}) were degassed under reduced pressure and dry toluene was added (\SI{8}{\mL}) obtaining a suspension.
To a solution of the crude 1-bromo-4-(1-bromo\-ethyl)\-benzene \cmpd+{ig2-17} (\SI{530}{\mg}, according \gls{GCMS} contains 55~\% 1-bromo-4-(1-bromo\-ethyl)\-benzene, \SI{1.10}{\mmol} and 45~\% of 4-bromo\-styrene, \SI{1.30}{\mmol}, in \SI{3}{\mL} toluene dry) toluene dry (\SI{10}{\mL}) and 2,2,5-tri\-methyl-4-phenyl-3-aza\-hexane-3-nitr\-oxide (\gls{TIPNO} radical, \SI{0.600}{\mL}, \SI{571}{\mg}, \SI{2.60}{\mmol}) were added. 
Then this mixture was degassed by bubbling nitrogen and \SI{1.88}{\mL} of the previously prepared suspension was added dropwise. After stirring the reaction mixture for \SI{24}{\hour} at \SI{75}{\celsius} (the suspension remained dark yellow rather turning to green as expected), the reaction was cooled down, adsorbed on basic allumina and filtered over basic alumina (column diameter \SI{2}{\cm}, height \SI{10}{\cm}), using hexane as the eluent to obtain \acrfull{BrPhEtTIPNO} \cmpd+{ig2-21} as a viscous colorless oil (\SI{333}{\mg}, from \gls{GCMS} and {\HNMR} analysis the product resulted to be a mixture with 10~\% of precursor \cmpd+{ig2-17} and 20~\% of di\-merized precursor isomers \ch{(BrC6H4C2H4)2}, so the yield in \cmpd+{ig2-21} was $\approx 50$~\%).

{\HNMR} (\SI{600}{\MHz}, \gls{TCE}-d$_2$) $\delta$ (ppm) 7.46 – 7.09 (m, 18H, arom.\ H), 4.83 (q, $J =$ \SI{6.6}{\Hz}, 1H, \ch{CHO}), 4.81 (q, $J =$ \SI{6.6}{\Hz}, 1H, \ch{CHO}), 3.35 (d, $J =$ \SI{10.6}{\Hz} (geminal), 1H, \ch{CHN}), 3.26 (d, $J =$ \SI{10.8}{\Hz} (geminal), 1H, \ch{CHN}), 2.24 (m, 1H, \ch{C\underline{H}(CH3)2}), 1.55 (d, $J =$ \SI{6.7}{\Hz}, 3H, \ch{CH(C\underline{H}3)2}), 1.47 (d, $J =$ \SI{6.6}{\Hz}, 3H, \ch{CH(C\underline{H}3)2}), 1.37 (m, 1H, \ch{C\underline{H}(CH3)2}), 1.24 (d, $J =$ \SI{6.3}{\Hz}, 3H, \ch{OCHC\underline{H}3}), 0.98 (s, 9H, \ch{C(C\underline{H}3)3}), 0.89 (d, $J =$ \SI{6.2}{\Hz}, 3H, \ch{OCHC\underline{H}3}), 0.72 (s, 9H, \ch{C(C\underline{H}3)3}), 0.48 (d, $J =$ \SI{6.6}{\Hz}, 3H, \ch{CH(C\underline{H}3)2}), 0.21 (d, $J =$ \SI{6.6}{\Hz}, 3H, \ch{CH(C\underline{H}3)2}).

\end{section}
\begin{section}[Second Block Polymerization]{Functionalization of Polythiophene as Macroinitiator and Second Block Polymerization}

\subsection{End-functionalization of Polythiophene with Aldehyde (\cmpd+{ig2-19})}

\begin{figure}[H]%syn9-suzuki
\centering
\includegraphics[scale=0.5]
{syn9-suzuki.pdf}
\end{figure}
Poly\{3-[((S)-2-methyl\-butyl\-oxy)\-methyl{]}\-thio\-phene\} \cmpd+{ig2-8} (\SI{14.4}{\mg}, quantity of chains estimated using $M_w$ \SI{0.976}{\umol}) was added in a \SI{100}{\mL} Schlenk tube with B-(5'-formyl[2,2'-bi\-thio\-phen]\-5-yl)\-boronic acid (\SI{1.37}{\mg}, \SI{5.75}{\umol}, solution in distilled tetrahydrofuran \SI{170}{\uL}) and bis(tri\-phenyl\-phosphine)\-palladium(II) dichloride (\ch{Pd(PPh3)2Cl2}, \SI{0.56}{\mg}, solution in distilled tetrahydrofuran \SI{175}{\uL}), then the mixture was dried on mechanical pump. 
Under nitrogen atmosphere dry tetrahydrofuran (\SI{6}{\mL}) was added. De\-oxy\-gen\-ated sodium carbonate aqueous solution (\SI{1}{\Molar}, \SI{1}{\mL}) was added and the reaction was refluxed for \SI{24}{\hour} under stirring in nitrogen atmosphere.
At room temperature the reaction was quenched by adding distilled water (\SI{10}{\mL}) and extracted with dichloromethane (3 $\times$ \SI{20}{\mL}). The collected organic fractions were washed with water, concentrated then poured in methanol (\SI{100}{\mL}). The precipitated was filtered (\gls{PTFE} membrane, \SI{0.45}{\um} porosity), washed with diethyl ether and solubilized again in chloroform. The chloroform solution was precipitated in methanol and filtered a second time obtaining a red solid \cmpd+{ig2-19} (\SI{11}{\mg}, 76~\% yield).

{\HNMR} (\SI{600}{\MHz}, \gls{TCE}-d$_2$) $\delta$ (ppm) 9.80 (s, 0.015H, \ch{H-C=O}), 7.21 (s, 1H, arom.\ H), 4.54 (s, 2H, \ch{CCH2O}), 3.40 (dd, $J =$ \SI{9.0}{\Hz} (geminal), \SI{6.1}{\Hz}, 1H, \ch{OC\underline{H}2CH}), 3.33 -- 3.29 (m, 1H, \ch{OC\underline{H}2CH}), 1.76 -- 1.65 (m, 1H, \ch{OCH2C\underline{H}}), 1.52 -- 1.40 (m, 1H, \ch{OCH2CHC\underline{H}2}), 1.20 -- 1.11 (m, 1H, \ch{OCH2CHC\underline{H}2}), 0.92 (d, $J =$ \SI{6.1}{\Hz}, 3H, \ch{CHC\underline{H}3}), 0.88 (t, $J =$ \SI{6.7}{\Hz}, 3H, \ch{CH2C\underline{H}3}).
\gls{FTIR} (film on \ch{KBr}) $\nu$ (\SI{}{\per\cm}) 2960, 2930, 2875, 1668, 1461, 1378, 1091, 840.

\subsection{Functionalization of Aldehyde on Polythiophene with Radical Polymerization Initiator (\cmpd+{ig2-22})}

\begin{figure}[H]%syn10-pt-tipno
\centering
\includegraphics[scale=0.5]
{syn10-pt-tipno.pdf}
\end{figure}

The previously synthesized bromo-alkoxy\-amine \cmpd+{ig2-21} (\SI{33}{\mg}, mixture containing \SI{23}{\mg} of N-tert-butyl-O-[1-[4-bromo\-phenyl]\-ethyl]-N-(2-methyl-1-phenyl\-propyl)\-hydroxyl\-amine, \gls{BrPhEtTIPNO}, \SI{57}{\umol}; \SI{3}{\mg} of precursor \cmpd+{ig2-17}, \SI{11}{\umol}; \SI{7}{\mg} of dimerized precursor isomers \ch{(BrC6H4C2H4)2}, \SI{19}{\umol}) was dissolved in dry tetrahydrofuran (\SI{3}{\mL}) and cooled to \SI{-40}{\celsius}. 
A solution of \textit{n}-butyl\-lithium (\SI{1.6}{\Molar} in hexane, \SI{75}{\uL}, \SI{120}{\umol}) was added dropwise while stirring, then was allowed to warm to room temperature. In another flask polythiophene-aldehyde (\SI{10}{\mg}) was dried under vacuum and dissolved in dry tetrahydrofuran (\SI{3}{\mL}) at \SI{50}{\celsius}. 
The solution of lithiated alkoxy\-amine was transferred via syringe into a solution of polythiophene-aldehyde solution, the reaction was kept at \SI{50}{\celsius} for one hour, then at room temperature it was diluted with chloroform and poured in acetone (\SI{100}{\mL}). 
The red solid was collected by filtration (\gls{PTFE} membrane, \SI{0.45}{\um}), dissolved in chloroform and washed with methanol obtaining polythiophene-\gls{TIPNO} \cmpd+{ig2-22} as a red powder (\SI{9}{\mg}, yield $\approx90$~\%).

{\HNMR} (\SI{500}{\MHz}, \gls{TCE}-d$_2$) $\delta$ (ppm) 7.21 (s, 1H, arom.\ H), 4.54 (s, 2H, \ch{CCH2O}), 3.43 -- 3.36 (m, 1H, \ch{OC\underline{H}2CH}), 3.34 -- 3.27 (m, 1H, \ch{OC\underline{H}2CH}), 1.76 -- 1.64 (m, 2H, \ch{OCH2C\underline{H}}), 1.51 -- 1.42 (m, 1H, \ch{OCH2CHC\underline{H}2}), 1.19 -- 1.10 (m, 1H, \ch{OCH2CHC\underline{H}2}), 0.92 (d, $J =$ \SI{6.4}{\Hz}, 3H, \ch{CHC\underline{H}3}), 0.88 (t, $J =$ \SI{7.2}{\Hz}, 3H, \ch{CH2C\underline{H}3}).

\subsection[Poly\{3-[((S)-2-methyl\-butyl\-oxy)\-methyl{]}\-thio\-phene\}-\textit{block}-poly(4-vinyl\-pyridine) (\cmpd+{sz17})]{Poly\{3-[((S)-2-methyl\-butyl\-oxy)\-methyl{]}\-thio\-phene\}-\textit{block}-poly(4-vinyl\-pyridine) (\cmpd+{sz17}) \\ NMRP polymerization of 4-vinyl\-pyridine}

\begin{figure}[H]%syn11-p4vp
\centering
\includegraphics[scale=0.5]
{syn11-p4vp.pdf}
\end{figure}

Polythiophene-\gls{TIPNO} \cmpd+{ig2-22} (\SI{9}{\mg}, \SI{0.9}{\umol} of polymeric chains) in a Schlenk tube was dissolved in degassed ortho-di\-chloro\-benzene (\SI{2}{\mL}) and sonicated. Then 4-vinyl\-pyridine (\SI{291}{\uL}, \SI{2.7}{\mmol}) and \acrlong{TIPNO} (\gls{TIPNO} radical, \SI{10}{\uL} of a 1~\% solution in ortho-di\-chloro\-benzene, \SI{0.4}{\umol}) were added. The whole reaction mixture was passed through three freeze/pump/thaw cycles and heated to \SI{125}{\celsius}. 
After \SI{4}{\hour} the mixture showed no changes in color or viscosity and the reaction was stopped cooling in liquid nitrogen. Then at room temperature was diluted with chloroform (\SI{2}{\mL}), poured in hexane (\SI{400}{\mL}) and stirred overnight. 
The so-obtained solid product was filtered (\gls{PTFE} membrane, \SI{0.45}{\um}), recovered in chloroform and precipitated again in methanol. The orange solid was washed first with hot toluene (\SI{70}{\celsius}) and the remaining product was recovered with chloroform to achieve two different fractions (\SI{6}{\mg} in toluene and \SI{6}{\mg} in chloroform).

{\HNMR} (\SI{600}{\MHz}, \ch{CDCl3}) $\delta$ (ppm) 8.36 (broad s, \gls{P4VP} arom.\ near N), 7.24 (s, thiop.\ arom.\ H), 6.48 (broad s, \gls{P4VP} arom.\ far N), 4.58 (s, thioph.\ \ch{CCH2O}), 3.49 -- 3.40 (m, thioph.\ \ch{OC\underline{H}2CH}), 3.39 -- 3.29 (m, thioph.\ \ch{OC\underline{H}2CH}), 1.78 -- 1.69 (m, thioph.\ \ch{OCH2C\underline{H}}), 1.55 -- 1.49 (m, thioph.\ \ch{OCH2CHC\underline{H}2}), 0.97 (d, $J =$ \SI{6.3}{\Hz}, thioph.\ \ch{CHC\underline{H}3}), 0.92 (t, $J =$ \SI{7.5}{\Hz}, thioph.\ \ch{CH2C\underline{H}3}), 0.90 -- 0.83 (m, \gls{P4VP} aliph.).

{\HNMR} (\SI{600}{\MHz}, \gls{TCE}) $\delta$ (ppm) 8.26 (broad s, \gls{P4VP} arom.\ near N), 7.21 (s, thiop.\ arom.\ H), 6.35 (broad s, \gls{P4VP} arom.\ far N), 4.54 (s, thioph.\ \ch{CCH2O}), 3.44 -- 3.35 (m, thioph.\ \ch{OC\underline{H}2CH}), 3.34 -- 3.25 (m, thioph.\ \ch{OC\underline{H}2CH}), 1.73 -- 1.65 (m, thioph.\ \ch{OCH2C\underline{H}}), 1.57 -- 1.38 (m, thioph.\ \ch{OCH2CHC\underline{H}2}), 0.92 (d, $J =$ \SI{6.5}{\Hz}, thioph.\ \ch{CHC\underline{H}3}), 0.88 (t, $J =$ \SI{7.4}{\Hz}, thioph.\ \ch{CH2C\underline{H}3}), 0.85 -- 0.77 (m, H, \gls{P4VP} aliph.).
\end{section}

\chapter*{Ringraziamenti}

Giunto alla fine di questa tesi voglio ringraziare chi mi ha aiutato, supportato o tenuto compagnia in questo lungo periodo di vita milanese.

Ringrazio Stefania Zappia e Marco Toscanini che m'hanno insegnato moltissimo in laboratorio e m'hanno aiutato moltissimo, senza di loro non sarei riuscito a far nulla.

Ringrazio i miei relatori: Silvia Destri e Lorenzo Di Bari per l'infinita pazienza che hanno avuto con me, per tutto il periodo di tesi ma soprattutto nelle ultime settimane quando la consegna si avvicinava.

Ringrazio tutti coloro che mi hanno analizzato i miei materiali e mi hanno spiegato un mucchio di cose interessantissime: William Porzio, Filippo Samperi \& Salvo Battiato, Andrea Zampetti, Alberto Giacometti Scheroni \& Daniele Piovani, Antonella Boccia, Maurizio Canetti, Fabio Bertini \& Adriana Cacciamani, Wojtek Mroz, Francesco Zinna.

Ringrazio i colleghi milanesi che sono stati la mia ancora di salvezza: Francesca, Elisa, Laura, Giorgio, Stefania, Marco, Wojtek, Sajjad.

Ringrazio i milanesi non colleghi, che m'hanno accolto facendomi sentire a casa: tutto il gruppo di Manitese di Milano, la bella ggente del Macao, Accio e Alessandro in particolare, i volontari di Fà la cosa giusta, Valentina e Jean Paul in particolare, Daniela, Salvo, Ottilia e Alberto.

Ringrazio moltisssssimo i pisani che m'hanno ospitato a Pisa nelle mie frequenti visite lampo e non: Svalo, Lucrezia, Ambra, Francesca, Ax, Costa, Bolzo, Soba, Yak, Pietro, Reca, Ciera.

Ringrazio gli altri bellissimi personi pisani e veronesi che m'hanno sopportato nelle serate, nelle mailing list, d'estate e nei week end e a cui voglio un sacco di bbene: tutti i Lannari Mannari con l'aggiunta della Cò e della Giulietta, tutta la gioventù di eigenLab che m'ha tenuto compagnia perfino a distanza, tutto il mio anno di scienze della Normale, i compagni del corso di chimica organica, il compagno di lab Leo, tutte le bimbe e i bimbi chimici del dipartimento, il Picconi che non può mai mancare, Lorenzo, Filippo, il Cieratto e tutti gli ex del glorioso Timpano.

Infine ringrazio mamma (ciao mamma!) e ringrazio la Normale soprattutto per avermi fatto conoscere un sacco di gente ganza.

(sono sicuro di dimenticarmi qualcuno\ldots vabbè\ldots)

\appendix \chapter{Supplementary Spectra}
\graphicspath{ {./img/spectra/} }

\renewcommand{\bottomfraction}{0.99} %aggiunto per far stare gli spettri nella prima pagina delle sezioni
\setcounter{bottomnumber}{3} %aggiunto per far stare gli spettri nella prima pagina delle sezioni

\clearpage
\section*{Gas Chromatography--Mass Spectrometry}

\begin{figure}%ig2-12-gcms
\includegraphics[width=\textwidth]{ig2-12-gcms.png}
\caption[GCMS of Grignard metathesis and quenching reaction.]{GCMS TIC of Grignard metathesis and quenching reaction described in Section~\ref{sec:monomero-regioselettivita}. Product \cmpd+{ig2-9} peak \SI{15.23}{\minute}, area 94~\%; precursor \cmpd+{ig2-10} peak \SI{20.92}{\minute}, area 6~\%.}
\label{fig:ig2-12-gcms}
\end{figure}%

\begin{figure}%ig2-13-gcms
\includegraphics[width=\textwidth]{ig2-13-gcms.png}
\caption[GCMS of Grignard metathesis in presence of lithium chloride and quenching reaction.]{GCMS TIC of Grignard metathesis in presence of lithium chloride and quenching reaction described in Section~\ref{sec:monomero-regioselettivita}. Product \cmpd+{ig2-9} peak \SI{15.22}{\minute}, area 78~\%; precursor \cmpd+{ig2-10} peak \SI{20.91}{\minute}, area 4~\%.}
\label{fig:ig2-13-gcms}
\end{figure}

\FloatBarrier
\clearpage

\FloatBarrier
\clearpage
\section*{Nuclear Magnetic Resonance}

\label{sec:hnmr}

\begin{figure}%ig2-4-nmr-h
\centering
\includegraphics[width=1\textwidth]{ig2-4-nmr-h.pdf}
\caption[{\HNMR} of polymer \cmpd+{ig2-4}.]{{\HNMR} (\SI{600}{\MHz}) of polymer \cmpd+{ig2-4} (\SI{35}{\mg}, \SI{5}{\mm} tube) in \gls{TCE}-d2.}
\label{fig:ig2-4-nmr-h}
\end{figure}

\begin{figure}%ig2-4-nmr-c
\centering
\includegraphics[trim=0 0 0 2.5cm,clip,width=1\textwidth]{ig2-4-nmr-c.pdf}
\caption[{\CNMR} of polymer \cmpd+{ig2-4}.]{{\CNMR} (\SI{150}{\MHz}) of polymer \cmpd+{ig2-4} (\SI{35}{\mg}, \SI{5}{\mm} tube) in \gls{TCE}-d2.}
\label{fig:ig2-4-nmr-c}
\end{figure}

\begin{figure}%ig2-15-nmr-c
\centering
\includegraphics[width=\textwidth]{ig2-15-nmr-c.pdf}
\caption[{\CNMR} of polymer \cmpd+{ig2-15}.]{{\CNMR} (\SI{150}{\MHz}) of polymer \cmpd+{ig2-15} in \gls{TCE}-d2.}
\label{fig:ig2-15-nmr-c}
\end{figure}

\begin{figure}%ig2-8-nmr-h
\centering
\includegraphics[width=1\textwidth]{ig2-8-nmr-h.pdf}
\caption[{\HNMR} of polymer \cmpd+{ig2-8}.]{{\HNMR} (\SI{600}{\MHz}) of polymer \cmpd+{ig2-8} in \gls{TCE}-d2.}
\label{fig:ig2-8-nmr-h}
\end{figure}

\begin{figure}%ig2-8-nmr-c
\centering
\includegraphics[trim=0 0 0 0.5cm,clip,width=1\textwidth]{ig2-8-nmr-c.pdf}
\caption[{\CNMR} of polymer \cmpd+{ig2-8}.]{{\CNMR} (\SI{150}{\MHz}) of polymer \cmpd+{ig2-8} in \gls{TCE}-d2.}
\label{fig:ig2-8-nmr-c}
\end{figure}

\begin{figure}%ig2-21-nmr-h
\centering
\includegraphics[width=0.95\textwidth]{ig2-21-nmr-h.pdf}
\caption[{\HNMR} of NMRP initiator PhEt-TIPNO \cmpd+{ig2-21}.]{{\HNMR} (\SI{600}{\MHz}) of \gls{NMRP} initiator \gls{BrPhEtTIPNO} \cmpd+{ig2-21} in \gls{TCE}-d2.}
\label{fig:ig2-21-nmr-h}
\end{figure}

\begin{figure}%ig2-19-nmr-h
\centering
\includegraphics[width=0.95\textwidth]{ig2-19-nmr-h.pdf}
\caption[{\HNMR} of polythiophene-aldehyde \cmpd+{ig2-19}.]{{\HNMR} (\SI{600}{\MHz}) of polythiophene-aldehyde \cmpd+{ig2-19} in \gls{TCE}-d2.}
\label{fig:ig2-19-nmr-h}
\end{figure}

\begin{figure}%ig2-22-nmr-h
\centering
\includegraphics[width=0.95\textwidth]{ig2-22-nmr-h.pdf}
\caption[{\HNMR} of polythiophene-TIPNO \cmpd+{ig2-22}.]{{\HNMR} (\SI{600}{\MHz}) of polythiophene-TIPNO \cmpd+{ig2-22} in \gls{TCE}-d2.}
\label{fig:ig2-22-nmr-h}
\end{figure}

\begin{figure}%sz17-nmr-h-fraz-toluene
\centering
\includegraphics[width=0.95\textwidth]{sz17-nmr-h-fraz-toluene.pdf}
\caption[{\HNMR} of \cmpd+{sz17} washing in toulene.]{{\HNMR} (\SI{600}{\MHz}) of \cmpd+{sz17} washing in toulene in \gls{TCE}-d2.}
\label{fig:sz17-nmr-h-fraz-toluene}
\end{figure}

\begin{figure}%sz17-nmr-h-fraz-cloroformio
\centering
\includegraphics[width=1\textwidth]{sz17-nmr-h-fraz-cloroformio.pdf}
\caption[{\HNMR} of \cmpd+{sz17} purified fraction.]{{\HNMR} (\SI{600}{\MHz}) of \cmpd+{sz17} purified fraction in \ch{CDCl3}.}
\label{fig:sz17-nmr-h-fraz-cloroformio}
\end{figure}

\FloatBarrier
\clearpage

\section*{Mass Spectrometry}

\begin{figure}%ig2-8-maldi-lineare
\begin{tikzpicture}
\begin{axis}[axis x line=bottom,axis y line=left,enlarge x limits=false,enlarge y limits=false,width=1\textwidth,height=6cm,yticklabels={,,},xlabel=Mass-to-charge ratio $m/Q$,ylabel=Intensity (a.\ u.),xmax=14000,cycle list name=linestyles*]
\addplot table {img/results/ig2-8-maldi-lineare-compresso.txt};
\end{axis}
\end{tikzpicture} 

\bigskip

\bigskip

\begin{tikzpicture}
\begin{axis}[axis x line=bottom,axis y line=left,enlarge x limits=false,enlarge y limits=false,width=1\textwidth,height=8cm,yticklabels={,,},xlabel=Mass-to-charge ratio $m/Q$,ylabel=Intensity (a.\ u.),xmin=2720,xmax=2960,cycle list name=linestyles*]
\addplot table {img/results/ig2-8-maldi-lineare-cut.txt};
\end{axis}
\end{tikzpicture} 
\caption[MALDI-TOF MS in linear mode of polymer \cmpd+{ig2-8}.]{\Acrfull{MALDI} in linear mode of polymer \cmpd+{ig2-8}.}
\label{fig:ig2-8-maldi-lineare}
\end{figure}

\FloatBarrier
\clearpage
\section*{Infrared Spectroscopy}

\begin{figure}%ig2-10-ir
\begin{tikzpicture}
\begin{axis}[enlarge x limits=false,enlarge y limits=false,width=1\textwidth,height=7cm,x dir=reverse,xminorgrids=true,xmajorgrids=true,ymajorgrids=true,minor x tick num=4,/pgf/number format/1000 sep={},max space between ticks=80pt,xlabel=Wavenumber (cm$^{-1}$),ylabel=Normalized transmittance (\%)]
\addplot[line width=0.5pt] file {img/spectra/ig2-10-ir.txt};
\end{axis}
\end{tikzpicture}
\caption[FT-IR of pure liquid \cmpd+{ig2-10}.]{FT-IR of pure viscous liquid \cmpd+{ig2-10} on \ch{KBr} disk, discussion in Section~\ref{sec:ir}.}
\label{fig:ig2-10-ir}
\end{figure}

\begin{figure}%ig2-4-ir
\begin{tikzpicture}
\begin{axis}[enlarge x limits=false,enlarge y limits=false,width=1\textwidth,height=7cm,x dir=reverse,xminorgrids=true,xmajorgrids=true,ymajorgrids=true,minor x tick num=4,/pgf/number format/1000 sep={},max space between ticks=80pt,xlabel=Wavenumber (cm$^{-1}$),ylabel=Normalized transmittance (\%)]
\addplot[line width=0.5pt] file {img/spectra/ig2-4-ir.txt};
\end{axis}
\end{tikzpicture}
\caption[FT-IR of polymer \cmpd+{ig2-4}.]{FT-IR of polymer \cmpd+{ig2-4} as a film drop cast from dichloromethane on \ch{KBr} disk, discussion in Section~\ref{sec:ir}.}
\label{fig:ig2-4-ir}
\end{figure}

\begin{figure}%ig2-15-ir
\begin{tikzpicture}
\begin{axis}[enlarge x limits=false,enlarge y limits=false,width=1\textwidth,height=7cm,x dir=reverse,xminorgrids=true,xmajorgrids=true,ymajorgrids=true,minor x tick num=4,/pgf/number format/1000 sep={},max space between ticks=80pt,xlabel=Wavenumber (cm$^{-1}$),ylabel=Normalized transmittance (\%)]
\addplot[line width=0.5pt] file {img/spectra/ig2-15-ir.txt};
\end{axis}
\end{tikzpicture}
\caption[FT-IR of polymer \cmpd+{ig2-15}.]{FT-IR of polymer \cmpd+{ig2-15} as a film drop cast from chloroform on \ch{KBr} disk, discussion in Section~\ref{sec:ir}.}
\label{fig:ig2-15-ir}
\end{figure}

\begin{figure}%ig2-8-ir
\begin{tikzpicture}
\begin{axis}[enlarge x limits=false,enlarge y limits=false,width=1\textwidth,height=7cm,x dir=reverse,xminorgrids=true,xmajorgrids=true,ymajorgrids=true,minor x tick num=4,/pgf/number format/1000 sep={},max space between ticks=80pt,xlabel=Wavenumber (cm$^{-1}$),ylabel=Normalized transmittance (\%)]
\addplot[line width=0.5pt] file {img/spectra/ig2-8-ir.txt};
\end{axis}
\end{tikzpicture}
\caption[FT-IR of polymer \cmpd+{ig2-8}.]{FT-IR of polymer \cmpd+{ig2-8} as a film drop cast from dichloromethane on \ch{KBr} disk, discussion in Section~\ref{sec:ir}.}
\label{fig:ig2-8-ir}
\end{figure}

\FloatBarrier
\clearpage
\section*{Ultraviolet--Visible Spectroscopy}

\begin{figure}%ig2-4-uvvis
\begin{tikzpicture}
\begin{axis}[axis x line=bottom,axis y line=left,yticklabels={,,},enlarge x limits=false,enlarge y limits=false,width=1\textwidth,height=8cm,xlabel=Wavelength (nm),ylabel=Normalized absorbance, xmin=240,xmax=700]
\addplot[line width=0.5pt]
table[y expr=\thisrowno{1}/0.701684] {img/results/ig2-4-uvvis-chcl3.txt};
\addlegendentry{\cmpd+{ig2-4} in \ch{CHCl3}};
\addplot[dashed]
table[y expr=\thisrowno{1}/0.398487-0.3] {img/results/ig2-4-uvvis-chcl3-50-meoh-50.txt};
\addlegendentry{\ch{CHCl3} 50 \% - \ch{MeOH} 50 \%};
\addplot[dashdotted] 
table[y expr=\thisrowno{1}/0.437508-0.7] {img/results/ig2-4-uvvis-chcl3-40-meoh-60.txt};
\addlegendentry{\ch{CHCl3} 40 \% - \ch{MeOH} 60 \%};
\addplot[dotted]
table[y expr=\thisrowno{1}/0.28404-1] {img/results/ig2-4-uvvis-chcl3-30-meoh-70.txt};
\addlegendentry{\ch{CHCl3} 30 \% - \ch{MeOH} 70 \%};
\end{axis}
\end{tikzpicture}
\caption[UV-vis spectra of polymer \cmpd+{ig2-4} in chloroform-methanol mixtures.]{UV-vis spectra of polymer \cmpd+{ig2-4} in chloroform-methanol mixtures.}
\label{fig:ig2-4-uvvis}
\end{figure}

\begin{figure}%ig2-15-uvvis
\begin{tikzpicture}
\begin{axis}[axis x line=bottom,axis y line=left,yticklabels={,,},enlarge x limits=false,enlarge y limits=false,width=1\textwidth,height=8cm,xlabel=Wavelength (nm),ylabel=Normalized absorbance, xmin=240,xmax=600,legend style={at={(0.05,1.25)},anchor=north west,font=\footnotesize}]
\addplot[line width=0.5pt]
table[y expr=\thisrowno{1}/1.04169] {img/results/ig2-15-uvvis-chcl3.txt};
\addlegendentry{\cmpd+{ig2-15} in \ch{CHCl3}};
\addplot[densely dashed]
table[y expr=\thisrowno{1}/1.0759-0.3] {img/results/ig2-15-uvvis-chcl3-50-ch3cn-50.txt};
\addlegendentry{\ch{CHCl3} 50 \% - \ch{CH3CN} 50 \%};
\addplot[dotted]
table[y expr=\thisrowno{1}/0.581413-0.6] {img/results/ig2-15-uvvis-thf.txt};
\addlegendentry{\cmpd+{ig2-15} in \gls{THF}};
\addplot[densely dotted] 
table[y expr=\thisrowno{1}/1.8789-0.9] {img/results/ig2-15-uvvis-thf-55-ch3cn-45.txt};
\addlegendentry{\gls{THF} 55 \% - \ch{CH3CN} 45 \%};
\addplot[dashdotted]
table[y expr=\thisrowno{1}/1.43846-1.25] {img/results/ig2-15-uvvis-thf-70-meoh-30.txt};
\addlegendentry{\gls{THF} 70 \% - \ch{MeOH} 30 \%};
\addplot[dashdotdotted]
table[y expr=\thisrowno{1}/0.983174-1.5] {img/results/ig2-15-uvvis-thf-60-meoh-40.txt};
\addlegendentry{\gls{THF} 60 \% - \ch{MeOH} 40 \%};
\end{axis}
\end{tikzpicture} \caption{UV-vis spectra of polymer \cmpd+{ig2-15} in mixtures of good-poor solvents.}
\label{fig:ig2-15-uvvis}
\end{figure}

\begin{figure}%ig2-8-uvvis-ch3cn
\begin{tikzpicture}
\begin{axis}[axis x line=bottom,axis y line=left,enlarge x limits=false,enlarge y limits=false,width=1\textwidth,height=8cm,xlabel=Wavelength (nm),ylabel=Molar absorptivity (\SI{}{\per\Molar\per\cm}),xmin=240,legend pos=north west,cycle list name=linestyles*]
\addplot table[y expr=\thisrowno{1}/0.0000753] {img/results/ig2-8-uvvis-137-chcl3-for-ch3cn.txt};
\addlegendentry{\cmpd+{ig2-8} in \ch{CHCl3}};
\addplot table[y expr=\thisrowno{1}/0.0000753] {img/results/ig2-8-uvvis-137-chcl3-40-ch3cn-60.txt};
\addlegendentry{\ch{CHCl3} 40 \% - \ch{CH3CN} 60 \%};
\addplot table[y expr=\thisrowno{1}/0.0000753] {img/results/ig2-8-uvvis-137-chcl3-30-ch3cn-70.txt};
\addlegendentry{\ch{CHCl3} 30 \% - \ch{CH3CN} 70 \%};
\addplot table[y expr=\thisrowno{1}/0.0000753] {img/results/ig2-8-uvvis-137-chcl3-20-ch3cn-80.txt};
\addlegendentry{\ch{CHCl3} 20 \% - \ch{CH3CN} 80 \%};
\addplot table[y expr=\thisrowno{1}/0.0000753] {img/results/ig2-8-uvvis-137-chcl3-10-ch3cn-90.txt};
\addlegendentry{\ch{CHCl3} 10 \% - \ch{CH3CN} 90 \%};
\end{axis}
\end{tikzpicture} \caption[UV-vis spectra of polymer \cmpd+{ig2-8} in chloroform-acetonitrile mixtures.]{UV-vis spectra of polymer \cmpd+{ig2-8} in chloroform-acetonitrile mixtures. Concentration \SI{137}{\mg\per\liter}. Pathlength \SI{1}{\mm}. Molar absorptivity is referred to the quantity of thiophene rings; $\epsilon = \mathrm{abs} / ( l (\mathrm{cm}) \cdot c (\mathrm{mol/L}))$.}
\label{fig:ig2-8-uvvis-ch3cn}
\end{figure}

\begin{figure}%ig2-4-uvvis-film
\begin{tikzpicture}
\begin{axis}[axis x line=bottom,axis y line=left,enlarge x limits=false,enlarge y limits=false,width=1\textwidth,height=8cm,xlabel=Wavelength (nm),ylabel=Absorbance,xmin=240,xmax=650,cycle list name=linestyles*]
\addplot table {img/results/ig2-4-uvvis-chcl3.txt};
\addlegendentry{\cmpd+{ig2-4} in \ch{CHCl3}};
\addplot table {img/results/ig2-4-uvvis-chcl3-30-meoh-70.txt};
\addlegendentry{\ch{CHCl3} 30 \% - \ch{MeOH} 70 \%};
\addplot table {img/results/ig2-4-uvvis-spin700-chcl3.txt};
\addlegendentry{spin \SI{700}{\rpm} from \ch{CHCl3}};
\end{axis}
\end{tikzpicture} \caption[UV-vis spectra of polymer \cmpd+{ig2-4} in solution and thin film.]{UV-vis spectra of polymer \cmpd+{ig2-4} in chloroform, in chloroform-methanol solution and thin film spin coated from a chloroform solution at \SI{700}{\rpm}.}
\label{fig:ig2-4-uvvis-film}
\end{figure}

\begin{figure}%ig2-8-uvvis-film
\begin{tikzpicture}
\begin{axis}[axis x line=bottom,axis y line=left,enlarge x limits=false,enlarge y limits=false,width=1\textwidth,height=8cm,xlabel=Wavelength (nm),ylabel=Absorbance,xmin=240,legend style={at={(1,1.2)},anchor=north east,
font=\footnotesize},cycle list name=linestyles*]
\addplot table {img/results/ig2-8-uvvis-cast-chcl3.txt};
\addlegendentry{cast from \ch{CHCl3}};
\addplot table {img/results/ig2-8-uvvis-cast-thf.txt};
\addlegendentry{cast from \gls{THF}};
\addplot table {img/results/ig2-8-uvvis-cast-ch2cl2.txt};
\addlegendentry{cast from \ch{CH2Cl2}};
\addplot table {img/results/ig2-8-uvvis-spin700-chcl3.txt};
\addlegendentry{spin \SI{700}{\rpm} from \ch{CHCl3}};
\addplot table {img/results/ig2-8-uvvis-spin700-thf.txt};
\addlegendentry{spin \SI{700}{\rpm} from \gls{THF}};
\addplot[densely dotted] table {img/results/ig2-8-uvvis-spin1000-chcl3.txt};
\addlegendentry{spin \SI{1000}{\rpm} from \ch{CHCl3}};
\end{axis}
\end{tikzpicture} \caption[UV-vis spectra of polymer \cmpd+{ig2-8} in thin films.]{UV-vis spectra of polymer \cmpd+{ig2-8} in thin films cast and spin coated from various solutions.}
\label{fig:ig2-8-uvvis-film}
\end{figure}

\begin{figure}%ig2-8-uvvis-film-dip
\begin{tikzpicture}
\begin{axis}[axis x line=bottom,axis y line=left,enlarge x limits=false,enlarge y limits=false,width=1\textwidth,height=8cm,xlabel=Wavelength (nm),ylabel=Absorbance,xmin=245,legend style={at={(1,1.2)},anchor=north east,
font=\footnotesize},cycle list name=linestyles*]
\addplot table {img/results/ig2-8-uvvis-spin700-chcl3.txt};
\addlegendentry{spin \SI{700}{\rpm} from \ch{CHCl3}};
\addplot table {img/results/ig2-8-uvvis-spin700-chcl3-dip-pentane.txt};
\addlegendentry{spin \SI{700}{\rpm} from \ch{CHCl3} dip pentane};
\addplot table[y expr=\thisrowno{1}-0.1] {img/results/ig2-8-uvvis-spin700-thf.txt};
\addlegendentry{spin \SI{700}{\rpm} from \gls{THF}};
\addplot table[y expr=\thisrowno{1}-0.1] {img/results/ig2-8-uvvis-spin700-thf-dip-pentane.txt};
\addlegendentry{spin \SI{700}{\rpm} from \gls{THF} dip pentane};
\addplot table[y expr=\thisrowno{1}-0.2] {img/results/ig2-8-uvvis-spin1000-chcl3.txt};
\addlegendentry{spin \SI{1000}{\rpm} from \ch{CHCl3}};
\addplot[densely dotted] table[y expr=\thisrowno{1}-0.2] {img/results/ig2-8-uvvis-spin1000-chcl3-dip-pentane.txt};
\addlegendentry{spin \SI{1000}{\rpm} from \ch{CHCl3} dip pentane};
\end{axis}
\end{tikzpicture} \caption[UV-vis spectra of polymer \cmpd+{ig2-8} in thin films spin coated and dipped in pentane.]{Stacked UV-vis spectra of polymer \cmpd+{ig2-8} in thin film spin coated from various solutions and dipped in pentane.}
\label{fig:ig2-8-uvvis-film-dip}
\end{figure}

\begin{figure}%uvvis-kcl
\begin{tikzpicture}
\begin{axis}[axis x line=bottom,axis y line=left,enlarge x limits=false,enlarge y limits=false,width=1\textwidth,height=8cm,xlabel=Wavelength (nm),ylabel=Absorbance,xmin=240,legend style={at={(1,0.75)},anchor=north east},
cycle list name=linestyles*]
\addplot table {img/results/ig2-4-uvvis-kcl.txt};
\addlegendentry{\cmpd+{ig2-4}};
\addplot table {img/results/ig2-15-uvvis-kcl-3.txt};
\addlegendentry{\cmpd+{ig2-15}};
\addplot table[y expr=(\thisrowno{1}*0.00961)+0.340] {img/results/ig2-8-ht-kcl-2.txt};
\addlegendentry{\cmpd+{ig2-8}};
\end{axis}
\end{tikzpicture} \caption[UV-vis spectra of polymers in \ch{KCl} disks.]{UV-vis spectra of polymers \cmpd+{ig2-4}, \cmpd+{ig2-15} and \cmpd+{ig2-8} in \ch{KCl} disks.}
\label{fig:uvvis-kcl}
\end{figure}
\FloatBarrier
\clearpage
\section*{Photoluminescence}

\begin{figure}%pl-cloroformio
\begin{tikzpicture}
\begin{axis}[axis x line=bottom,axis y line=left,enlarge x limits=false,enlarge y limits=false,width=1\textwidth,height=6cm,xlabel=Wavelength (nm),ylabel=Photoluminescence (a.\ u.),xmin=480,legend style={font=\footnotesize},cycle list name=linestyles*]
\addplot table {img/results/ig2-4-pl-137-chcl3.txt};
\addlegendentry{\cmpd{ig2-4} \SI{137}{\mg\per\liter} in \ch{CHCl3}};
\addplot table {img/results/ig2-4-pl-5480-chcl3.txt};
\addlegendentry{\cmpd{ig2-4} \SI{5.48}{\g\per\liter} in \ch{CHCl3}};
\addplot table {img/results/ig2-15-pl-137-chcl3.txt};
\addlegendentry{\cmpd{ig2-15} \SI{137}{\mg\per\liter} in \ch{CHCl3}};
\addplot table {img/results/ig2-15-pl-5480-chcl3.txt};
\addlegendentry{\cmpd{ig2-15} \SI{5.48}{\g\per\liter} in \ch{CHCl3}};
\addplot table {img/results/ig2-8-pl-137-chcl3.txt};
\addlegendentry{\cmpd{ig2-8} \SI{137}{\mg\per\liter} in \ch{CHCl3}};
\addplot[densely dotted] table {img/results/ig2-8-pl-5480-chcl3.txt};
\addlegendentry{\cmpd{ig2-8} \SI{5.48}{\g\per\liter} in \ch{CHCl3}};
\end{axis}
\end{tikzpicture} \caption[Photoluminescence of polymers in chloroform solution.]{Photoluminescence of polymers \cmpd+{ig2-4}, \cmpd+{ig2-15} and \cmpd+{ig2-8} in chloroform solution.}
\label{fig:pl-cloroformio}
\end{figure}

\begin{figure}%ig2-4-pl
\begin{tikzpicture}
\begin{axis}[axis x line=bottom,axis y line=left,enlarge x limits=false,enlarge y limits=false,width=1\textwidth,height=6cm,xlabel=Wavelength (nm),ylabel=Photoluminescence (a.\ u.),xmin=480,cycle list name=linestyles*]
\addplot table {img/results/ig2-4-pl-137-chcl3.txt};
\addlegendentry{\cmpd{ig2-4} \SI{137}{\mg\per\liter} in \ch{CHCl3}};
\addplot table {img/results/ig2-4-pl-137-chcl3-50-meoh-50.txt};
\addlegendentry{in \ch{CHCl3} 50 \% - \ch{MeOH} 50 \%};
\end{axis}
\end{tikzpicture} \caption{Photoluminescence of polymer \cmpd+{ig2-4} in chloroform-methanol solution.}
\label{fig:ig2-4-pl}
\end{figure}

\FloatBarrier
\clearpage
\section*{Circular Dichroism Spectroscopy}

\begin{figure}%cd-goodsolvent
\begin{tikzpicture}
\begin{axis}[axis x line=bottom,axis y line=left,enlarge x limits=false,enlarge y limits=true,width=1\textwidth,height=6cm,xlabel=Wavelength (nm),ylabel=Circular Dichroism (mdeg),xmin=240,legend pos=north west]
\addplot[line width=0.5pt] table {img/results/ig2-15-cd-thf.txt};
\addlegendentry{\cmpd+{ig2-15} in \gls{THF}};
\addplot[line width=1.7pt] table {img/results/ig2-8-cd-137-chcl3.txt};
\addlegendentry{\cmpd+{ig2-8} in \ch{CHCl3}};
\draw[help lines] (axis cs:240,0) -- (axis cs:600,0);
\end{axis}
\begin{axis}[width=1\textwidth,
enlarge x limits=false,enlarge y limits=true,
axis y line=right,
ylabel=Absorbance,
yticklabels={,,},
axis x line=none,
xmin=240,
height=6cm]
\addplot[line width=0.5pt] table {img/results/ig2-15-ht-thf.txt};
\addplot[line width=1.7pt] table {img/results/ig2-8-ht-137-chcl3.txt};
\end{axis}
\end{tikzpicture} \caption[Circular dichroism of polymer \cmpd+{ig2-15} and \cmpd+{ig2-8} in good solvents solutions.]{Circular dichroism and UV-vis absorbance (as recorder by circular dichroism spectrometer) of polymer \cmpd+{ig2-15} and \cmpd+{ig2-8} in good solvents solutions. Concentration for polymer \cmpd+{ig2-8} was \SI{137}{\mg\per\liter}, for \cmpd+{ig2-15} concentration is unknown. Pathlength \SI{1}{\mm}.}
\label{fig:cd-goodsolvent}
\end{figure}

\begin{figure}%ig2-4-cd-meoh
\begin{tikzpicture}
\begin{axis}[axis x line=bottom,axis y line=left,enlarge x limits=false,enlarge y limits=true,width=1\textwidth,height=6cm,xlabel=Wavelength (nm),ylabel=Molar circular dichroism $\Delta\epsilon$,xmin=240,legend style={at={(0.28,1.05)},anchor=north west,
font=\footnotesize},
cycle list name=linestyles*]
\addplot table[y expr=\thisrowno{1}*0.0806] {img/results/ig2-4-cd-137-chcl3-60-meoh-40.txt};
\addlegendentry{\ch{CHCl3} 60 \% - \ch{MeOH} 40 \%};
\addplot table[y expr=\thisrowno{1}*0.0806] {img/results/ig2-4-cd-137-chcl3-50-meoh-50.txt};
\addlegendentry{\ch{CHCl3} 50 \% - \ch{MeOH} 50 \%};
\addplot table[y expr=\thisrowno{1}*0.0806] {img/results/ig2-4-cd-137-chcl3-40-meoh-60.txt};
\addlegendentry{\ch{CHCl3} 40 \% - \ch{MeOH} 60 \%};
\addplot[line width=1.6pt] table[y expr=\thisrowno{1}*0.0806] {img/results/ig2-4-cd-137-chcl3-20-meoh-80.txt};
\addlegendentry{\ch{CHCl3} 20 \% - \ch{MeOH} 80 \%};
\end{axis}
\end{tikzpicture} \caption[Circular dichroism of polymer \cmpd+{ig2-4} in chloroform-methanol mixtures.]{Circular dichroism of polymer \cmpd+{ig2-4} in chloroform-methanol mixtures. Concentration \SI{137}{\mg\per\liter}. Pathlength \SI{5}{\mm}. Molar circular dichroism is referred to the quantity of thiophene rings; $\Delta\epsilon = \theta (\mathrm{mdeg}) / (32982 \cdot l (\mathrm{cm}) \cdot c (\mathrm{mol/L}))$.}
\label{fig:ig2-4-cd-meoh}
\end{figure}

\begin{figure}%ig2-15-g-thf-ch3cn
\begin{tikzpicture}
\begin{axis}[axis x line=bottom,axis y line=left,enlarge x limits=false,enlarge y limits=true,width=1\textwidth,height=8cm,xlabel=Wavelength (nm),ylabel=Dissymmetry factor $g$,xmin=250,xmax=550,legend pos=north west,cycle list name=linestyles*]
\addplot table {img/results/ig2-15-g-137-thf-50-ch3cn-50.txt};
\addlegendentry{in \gls{THF} 50 \% - \ch{CH3CN} 50 \%};
\addplot table {img/results/ig2-15-g-274-thf-50-ch3cn-50.txt};
\addlegendentry{in \gls{THF} 50 \% - \ch{CH3CN} 50 \% \SI{274}{\mg\per\liter}};
\addplot table {img/results/ig2-15-g-548-thf-50-ch3cn-50.txt};
\addlegendentry{in \gls{THF} 50 \% - \ch{CH3CN} 50 \% \SI{548}{\mg\per\liter}};
\addplot table {img/results/ig2-15-g-137-chcl3-50-meoh-50.txt};
\addlegendentry{in \ch{CHCl3} 50 \% - \ch{MeOH} 50 \%};
\end{axis}
\end{tikzpicture} \caption[Dissymmetry factor of polymer \cmpd+{ig2-15} at various concentrations in THF - \ch{CHCl3}.]{Dissymmetry factor of polymer \cmpd+{ig2-15} in \gls{THF} 50~\% - \ch{CH3CN} 50~\% at various concentrations and comparison with polymer in \ch{CHCl3} 50~\% - \ch{MeOH} 50~\% mixture. Where not specified a concentration of \SI{137}{\mg\per\liter} was used. Pathlength \SI{1}{\mm}.}
\label{fig:ig2-15-g-thf-ch3cn}
\end{figure}

\begin{figure}%ig2-15-cd-anomali
\begin{tikzpicture}
\begin{axis}[axis x line=bottom,axis y line=left,enlarge x limits=false,enlarge y limits=true,width=1\textwidth,height=8cm,xlabel=Wavelength (nm),ylabel=Circular Dichroism (mdeg),xmin=240,legend pos=south west,cycle list name=linestyles*]
\addplot table {img/results/ig2-15-cd-chcl3-30-ch3cn-70-anomalo.txt};
\addlegendentry{in \ch{CHCl3} 30 \% - \ch{CH3CN} 70 \%};
\addplot table {img/results/ig2-15-cd-thf-50-ch3cn-50-anomalo.txt};
\addlegendentry{in \gls{THF} 50 \% - \ch{CH3CN} 50 \%};
\addplot table {img/results/ig2-15-cd-thf-50-ch3cn-50-anomalo-1day.txt};
\addlegendentry{in \gls{THF} 50 \% - \ch{CH3CN} 50 \% \SI{1}{\day}};
\end{axis}
\end{tikzpicture} \caption[Unusual circular dichroism found for polymer \cmpd+{ig2-15}.]{Unusual circular dichroism found for polymer \cmpd+{ig2-15} in chloroform-acetonitrile and tetrahydrofuran-acetonitrile mixture. Concentration not known.}
\label{fig:ig2-15-cd-anomali}
\end{figure}

\begin{figure}%ig2-8-cd-ch3cn
\begin{tikzpicture}
\begin{axis}[axis x line=bottom,axis y line=left,enlarge x limits=false,enlarge y limits=true,width=1\textwidth,height=6cm,xlabel=Wavelength (nm),ylabel=Molar circular dichroism $\Delta\epsilon$,xmin=240,legend pos=north west,cycle list name=linestyles*]
\addplot table[y expr=\thisrowno{1}*0.403] {img/results/ig2-8-cd-137-chcl3-40-ch3cn-60.txt};
\addlegendentry{in \ch{CHCl3} 40 \% - \ch{CH3CN} 60 \%};
\addplot table[y expr=\thisrowno{1}*0.403] {img/results/ig2-8-cd-137-chcl3-30-ch3cn-70.txt};
\addlegendentry{in \ch{CHCl3} 30 \% - \ch{CH3CN} 70 \%};
\addplot table[y expr=\thisrowno{1}*0.403] {img/results/ig2-8-cd-137-chcl3-20-ch3cn-80.txt};
\addlegendentry{in \ch{CHCl3} 20 \% - \ch{CH3CN} 80 \%};
\addplot table[y expr=\thisrowno{1}*0.403] {img/results/ig2-8-cd-137-chcl3-10-ch3cn-90.txt};
\addlegendentry{in \ch{CHCl3} 10 \% - \ch{CH3CN} 90 \%};
\end{axis}
\end{tikzpicture} \caption[Circular dichroism of polymer \cmpd+{ig2-8} in chloroform-acetonitrile.]{Circular dichroism of polymer \cmpd+{ig2-8} in chloroform-acetonitrile mixtures. Concentration \SI{137}{\mg\per\liter}. Pathlength \SI{1}{\mm}.}
\label{fig:ig2-8-cd-ch3cn}
\end{figure}

\begin{figure}%ig2-8-cd-thf-ch3cn
\begin{tikzpicture}
\begin{axis}[axis x line=bottom,axis y line=left,enlarge x limits=false,enlarge y limits=true,width=1\textwidth,height=6cm,xlabel=Wavelength (nm),ylabel=Molar circular dichroism $\Delta\epsilon$,xmin=240,legend pos=north west,cycle list name=linestyles*]
\addplot table[y expr=\thisrowno{1}*0.403] {img/results/ig2-8-cd-137-thf-60-ch3cn-40.txt};
\addlegendentry{in \gls{THF} 60 \% - \ch{CH3CN} 40 \%};
\addplot table[y expr=\thisrowno{1}*0.403] {img/results/ig2-8-cd-137-thf-50-ch3cn-50.txt};
\addlegendentry{in \gls{THF} 50 \% - \ch{CH3CN} 50 \%};
\addplot table[y expr=\thisrowno{1}*0.403] {img/results/ig2-8-cd-137-thf-40-ch3cn-60.txt};
\addlegendentry{in \gls{THF} 40 \% - \ch{CH3CN} 60 \%};
\addplot table[y expr=\thisrowno{1}*0.403] {img/results/ig2-8-cd-137-thf-20-ch3cn-80.txt};
\addlegendentry{in \gls{THF} 20 \% - \ch{CH3CN} 80 \%};
\end{axis}
\end{tikzpicture} \caption[Circular dichroism of polymer \cmpd+{ig2-8} in tetrahydrofuran-acetonitrile.]{Circular dichroism of polymer \cmpd+{ig2-8} in tetrahydrofuran-acetonitrile mixtures. Concentration \SI{137}{\mg\per\liter}. Pathlength \SI{1}{\mm}.}
\label{fig:ig2-8-cd-thf-ch3cn}
\end{figure}

\begin{figure}%ig2-4-cd-film
\begin{tikzpicture}
\begin{axis}[axis x line=bottom,axis y line=left,enlarge x limits=false,enlarge y limits=true,width=1\textwidth,height=6cm,xlabel=Wavelength (nm),ylabel=Dissymmetry factor $g$,xmin=250,xmax=550,legend style={at={(0.2,1)},anchor=north west},
cycle list name=linestyles*]
\addplot table {img/results/ig2-4-g-spin700-chcl3.txt};
\addlegendentry{\cmpd+{ig2-4} spin \SI{700}{\rpm} from \ch{CHCl3}};
\addplot table {img/results/ig2-4-g-137-chcl3-40-meoh-60-5mm.txt};
\addlegendentry{in \ch{CHCl3} 40 \% - \ch{MeOH} 60 \%};
\end{axis}
\end{tikzpicture} \caption[Dissymmetry factor \cmpd+{ig2-4} in thin film spin coated.]{Dissymmetry factor of polymer \cmpd+{ig2-4} in thin film spin coated from chloroform and comparison with chloroform-methanol mixture (rescaled). $g(\lambda) = \Delta\epsilon(\lambda) / \epsilon(\lambda)$.}
\label{fig:ig2-4-cd-film}
\end{figure}

\begin{figure}%ig2-15-cd-film
\begin{tikzpicture}
\begin{axis}[axis x line=bottom,axis y line=left,enlarge x limits=false,enlarge y limits=true,width=1\textwidth,height=6cm,xlabel=Wavelength (nm),ylabel=Dissymmetry factor $g$,xmin=250,xmax=550,legend style={at={(0.25,1)},anchor=north west},
cycle list name=linestyles*]
\addplot table {img/results/ig2-15-g-cast-ch2cl2.txt};
\addlegendentry{\cmpd+{ig2-15} cast from \ch{CH2Cl2}};
\addplot table {img/results/ig2-15-g-137-chcl3-50-meoh-50.txt};
\addlegendentry{in \ch{CHCl3} 50 \% - \ch{MeOH} 50 \%};
\end{axis}
\end{tikzpicture} \caption[Dissymmetry factor of polymer \cmpd+{ig2-15} in thin film cast from dichloromethane.]{Dissymmetry factor of polymer \cmpd+{ig2-15} in thin film cast from dichloromethane and comparison with polymer in chloroform-methanol mixture (rescaled).}
\label{fig:ig2-15-cd-film}
\end{figure}

\begin{figure}%ig2-8-cd-film-dip
\begin{tikzpicture}
\begin{axis}[axis x line=bottom,axis y line=left,enlarge x limits=false,enlarge y limits=true,width=1\textwidth,height=8cm,xlabel=Wavelength (nm),ylabel=Circular Dichroism (mdeg),xmin=240,legend style={at={(0.05,1.4)},anchor=north west,font=\footnotesize},
cycle list name=linestyles*]
\addplot table {img/results/ig2-8-cd-spin700-thf.txt};
\addlegendentry{spin \SI{700}{\rpm} frpm \gls{THF}};
\addplot table {img/results/ig2-8-cd-spin700-thf-dip-pentane.txt};
\addlegendentry{spin \SI{700}{\rpm} frpm \gls{THF} dip pentane};
\addplot table[y expr=\thisrowno{1}-1] {img/results/ig2-8-cd-spin700-chcl3.txt};
\addlegendentry{spin \SI{700}{\rpm} frpm \ch{CHCl3}};
\addplot table[y expr=\thisrowno{1}-1] {img/results/ig2-8-cd-spin700-chcl3-dip-pentane.txt};
\addlegendentry{spin \SI{700}{\rpm} frpm \ch{CHCl3} dip pentane};
\addplot table[y expr=\thisrowno{1}-2] {img/results/ig2-8-cd-spin1000-chcl3.txt};
\addlegendentry{spin \SI{1000}{\rpm} frpm \ch{CHCl3}};
\addplot[densely dotted] table[y expr=\thisrowno{1}-2] {img/results/ig2-8-cd-spin1000-chcl3-dip-pentane.txt};
\addlegendentry{spin \SI{1000}{\rpm} frpm \ch{CHCl3} dip pentane};
\end{axis}
\end{tikzpicture} \caption[Circular dichroism of polymer \cmpd+{ig2-8} spin coated and dipped in pentane.]{Circular dichroism of polymer \cmpd+{ig2-8} in thin film from spin coating and dipping in pentane. Spectra are stacked.}
\label{fig:ig2-8-cd-film-dip}
\end{figure}

\FloatBarrier
\clearpage
\section*{Differential Scanning Calorimetry}

\begin{figure}%ig2-15-dsc
\begin{tikzpicture}
\begin{axis}[axis x line=bottom,axis y line=left,enlarge x limits=false,enlarge y limits=true,width=1\textwidth,height=8cm,yticklabels={,,},xlabel=Sample Temperature (\SI{}{\celsius}),ylabel=Heat Flow,legend style={at={(0.28,1)},anchor=north west},cycle list name=linestyles*]
\addplot table {img/results/ig2-15-dsc-risc1.txt};
\addlegendentry{first heating \SI{20}{\celsius\per\minute}};
\addplot table {img/results/ig2-15-dsc-raff1.txt};
\addlegendentry{cooling \SI{10}{\celsius\per\minute}};
\addplot table {img/results/ig2-15-dsc-risc2.txt};
\addlegendentry{second heating \SI{20}{\celsius\per\minute}};
\end{axis}
\end{tikzpicture} \caption[Differential scanning calorimetry of polymer \cmpd+{ig2-15}.]{Differential scanning calorimetry of polymer \cmpd+{ig2-15} in absence of oxygen.}
\label{fig:ig2-15-dsc}
\end{figure}

\begin{figure}%ig2-8-dsc
\begin{tikzpicture}
\begin{axis}[axis x line=bottom,axis y line=left,enlarge x limits=false,enlarge y limits=true,width=1\textwidth,height=6cm,yticklabels={,,},xlabel=Sample Temperature (\SI{}{\celsius}),ylabel=Heat Flow,ymin=34.7,legend pos=north west,cycle list name=linestyles*]
\addplot table {img/results/ig2-8-dsc-risc1.txt};
\addlegendentry{heating \SI{20}{\celsius\per\minute}};
\addplot table {img/results/ig2-8-dsc-raff1.txt};
\addlegendentry{cooling \SI{10}{\celsius\per\minute}};
\end{axis}
\end{tikzpicture} \caption[Differential scanning calorimetry of polymer \cmpd+{ig2-8}.]{Differential scanning calorimetry of polymer \cmpd+{ig2-8} in absence of oxygen.}
\label{fig:ig2-8-dsc}
\end{figure}

\FloatBarrier
\clearpage
\section*{X--Ray Diffraction}

\begin{figure}%xrd-films
\begin{tikzpicture}
\begin{axis}[axis x line=bottom,axis y line=left,enlarge x limits=false,enlarge y limits=true,width=1\textwidth,height=8cm,xlabel=Degrees $2\theta$,ylabel=Intensity (a.\ u.),legend style={font=\footnotesize},legend pos=north west]
\addplot[line width=0.3pt] table {img/results/ig2-4-xrd-spin700.txt};
\addlegendentry{\cmpd{ig2-4} spin \SI{700}{\rpm}};
\addplot[line width=0.8pt] table[y expr=\thisrowno{1}-30] {img/results/ig2-4-xrd-spin1500.txt};
\addlegendentry{\cmpd{ig2-4} spin \SI{1500}{\rpm}};
\addplot[line width=1.3pt] table[y expr=\thisrowno{1}-60] {img/results/ig2-8-xrd-cast-ch2cl2.txt};
\addlegendentry{\cmpd{ig2-8} cast from \ch{CH2Cl2}};
\addplot[line width=1.8pt] table[y expr=\thisrowno{1}-90] {img/results/ig2-8-xrd-cast-ch2cl2-anneal.txt};
\addlegendentry{\cmpd{ig2-8} cast from \ch{CH2Cl2} and annealed};
\addplot[dotted] table[y expr=\thisrowno{1}-120] {img/results/quartz-xrd.txt};
\addlegendentry{quartz baseline};
\end{axis}
\end{tikzpicture} \caption[X-ray diffraction of polymer \cmpd+{ig2-4} and \cmpd+{ig2-8} thin films]{X-ray diffraction of polymer \cmpd+{ig2-4} and \cmpd+{ig2-8} thin films from spin coating or coating on quartz.}
\label{fig:xrd-films}
\end{figure}

\begin{figure}%xrd-smallangle
\begin{tikzpicture}
\begin{axis}[axis x line=bottom,axis y line=left,enlarge x limits=false,enlarge y limits=true,width=1\textwidth,height=6cm,xlabel=Degrees $2\theta$,ylabel=Intensity (a.\ u.),xmin=0.2,legend pos=north west,cycle list name=linestyles*]
\addplot table {img/results/ig2-8-xrdsmallangle.txt};
\end{axis}
\end{tikzpicture} \caption[Small angle X--ray scattering of polymer \cmpd+{ig2-8} in bulk solid.]{\Acrfull{SAXS} of polymer \cmpd+{ig2-8} in bulk solid.}
\label{fig:xrd-smallangle}
\end{figure}

\vfill
\begin{figure}[bh]
\centering
Ilario Gelmetti - \textit{Master Thesis} - 15 October 2013\\
\medskip 
Synthesis and Characterization of Chiral Polythiophenes.
\bigskip

\includegraphics[width=0.3\textwidth]{img/cc-by-sa.pdf}
\end{figure}


\end{document}
